\chapter{Dermatology}

\medterm{Bowen Disease}-- Squamous cell carcinoma in situ confined to the epidermis, presenting as a persistent scaly, sharply demarcated red or brown patch or plaque. It is associated with ultraviolet exposure, human papillomavirus, arsenic, and immunosuppression; treatment includes excision, curettage and cautery, cryotherapy, topical therapy, or photodynamic therapy according to site and lesion features.

\medterm{Brazilian Pemphigus Foliaceus}-- An endemic form of pemphigus foliaceus, also called fogo selvagem, occurring mainly in parts of Brazil and other tropical regions. Autoantibodies against desmoglein 1 cause superficial fragile blisters, erosions, crusting, and widespread scaling without the prominent mucosal disease typical of pemphigus vulgaris.

\medterm{Alopecia Celsi}-- See \hyperlink{Alopecia Areata}{alopecia areata}.
\medterm{Alopecia Circumscripta}-- See \hyperlink{Alopecia Areata}{alopecia areata}.
\medterm{Aplasia Cutis Congenita}-- A congenital absence of skin, most often on the scalp, presenting as a sharply demarcated ulcer, erosion, or membranous scar. Small lesions may heal with conservative wound care, while large or deep defects can involve the skull or dura and require specialist surgical management. Associated anomalies and syndromes should be assessed.

\medterm{ABCDE Rule for Melanoma}-- A clinical screening mnemonic for suspicious pigmented lesions: Asymmetry, Border irregularity, Color variation, Diameter greater than approximately 6 mm, and Evolution or change. It is a screening aid rather than a diagnostic test; suspicious lesions require dermoscopic assessment and biopsy.

\medterm{Acanthosis Nigricans}-- Hyperpigmented, velvety, thickened plaques typically in skin flexures (neck, axillae, groin) and often associated with insulin resistance, obesity, type 2 diabetes mellitus, and metabolic syndrome. May be drug-induced (nicotinic acid, glucocorticoids) or paraneoplastic (gastric and other gastrointestinal adenocarcinomas--malignant acanthosis nigricans). Treatment addresses the underlying cause; weight loss and metformin may improve insulin-resistant forms.

\medterm{Acne Rosacea Subtypes}-- Rosacea subtypes include erythematotelangiectatic (persistent facial erythema and flushing), papulopustular (red bumps and pus-filled pimples), phymatous (skin thickening, rhinophyma), and ocular (eye irritation, blepharitis). Demodex mite overgrowth contributes to pathogenesis. Triggers include sun, heat, spicy foods, alcohol, and stress. Treatment includes metronidazole cream, ivermectin cream, oral doxycycline, and laser therapy for telangiectasias.

\medterm{Acne Vulgaris}-- A chronic inflammatory disease of the pilosebaceous unit (hair follicle and associated sebaceous gland), affecting approximately 85 percent of adolescents and young adults between ages 12 and 24, with a prevalence of 12 to 22 percent in young adult women and 3 percent in men, making it the most common skin condition worldwide. Pathogenesis is multifactorial: excess sebum production (driven by androgens, particularly dihydrotestosterone [DHT] converted by 5-alpha-reductase from ttestosterone by 5-alpha-reductase type 1 in the sebaceous gland), follicular hyperkeratinisation (abnormal desquamation of follicular epithelial cells leading to microcomedone formation), Cutibacterium acnes colonisation and proliferation within the follicle, and inflammation mediated by both innate and adaptive immune responses targeting C. acnes antigens. Comedones (open blackheads and closed whiteheads) are the primary non-inflammatory lesions, while papules, pustules, and nodules represent inflammatory lesions.

\medterm{Acral lentiginous melanoma}-- A subtype of melanoma arising on acral skin (palms, soles, or under nails), appearing as a dark, irregularly pigmented patch or plaque. It is more common in people with darker skin tones and often presents at a later stage due to less frequent surveillance of these sites. Prognosis depends on thickness and stage at diagnosis.

\medterm{Acrodermatitis Enteropathica}-- A zinc-deficiency disorder causing periorificial and acral dermatitis, alopecia, diarrhoea, poor wound healing, and increased susceptibility to infection. It may be inherited from impaired intestinal zinc absorption or acquired through malnutrition, malabsorption, or prolonged parenteral nutrition. Diagnosis is clinical and supported by low zinc levels; treatment is zinc replacement and correction of the underlying cause.

\medterm{Actinic granuloma}-- A benign granulomatous skin condition arising on sun-damaged skin, presenting as annular or ring-shaped plaques with central atrophy and an elevated, erythematous border. It is thought to represent a foreign-body reaction to degenerated elastic tissue from chronic ultraviolet exposure. It is most common in older adults on sun-exposed areas.

\medterm{Actinic Keratosis}-- A premalignant lesion caused by chronic UV exposure, appearing as rough, scaly, erythematous papules on sun-exposed skin. AK is the precursor to squamous cell carcinoma, with approximately 10 percent progressing to invasive cancer. Risk factors include fair skin, outdoor occupation, immunosuppression, and human papillomavirus co-infection. Field cancerization describes widespread actinic damage. Treatment includes cryotherapy, topical 5-fluorouracil, imiquimod, diclofendiclofenac gel, and photodynamic therapy. Regular surveillance for progression to squamous cell carcinoma is necessary.

\medterm{Actinic Keratosis Management}-- Actinic keratosis management requires field-directed and lesion-directed therapy. Lesion-directed: cryotherapy (liquid nitrogen, 5 to 10 second freeze-thaw cycle), curettage, electrodessication. Field-directed: 5-fluorouracil cream (1 or 5 percent daily for 2 to 4 weeks), imiquimod cream (3 or 5 percent), diclofenac gel (twice daily for 60 to 90 days), photodynamic therapy (ALA or MAL followed by blue or red light), and topical tirbanibulin. Regular skin cancer screening is essential for patientpatients with regular follow-up and lesion mapping.

\medterm{Acute Generalized Exanthematous Pustulosis}-- An acute drug reaction characterised by dozens to thousands of small, sterile, nonfollicular pustules arising on an erythematous base, often with fever and neutrophilia. It usually begins rapidly after exposure to an antibiotic or other medication and resolves after withdrawal of the trigger. Severe mucosal disease, organ involvement, or diagnostic uncertainty requires urgent specialist assessment to distinguish it from pustular psoriasis and severe cutaneous adverse reactions.

\medterm{Albinism (Oculocutaneous Albinism)}-- A hereditary, congenital pigmentary disorder characterized by a complete or partial absence of melanin in the skin, hair, and eyes, most commonly caused by mutations in the tyrosinase gene (e.g., OCA1) or other genes involved in melanin synthesis. It presents with white or very light hair and skin, nystagmus, photophobia, and reduced visual acuity. Individuals have a significantly increased risk of developing ultraviolet (UV) radiation-induced skin malignancies, particularly squamous cell carcinoma. Management focuses on strict sun protection, regular skin cancer screening, and ophthalmologic care.

\medterm{Allergic Dermatitis}-- An inflammatory skin reaction caused by an immune response to an external allergen, most commonly allergic contact dermatitis. It is a delayed, T-cell-mediated hypersensitivity reaction that typically develops after prior sensitization. Clinical features include pruritus, erythema, oedema, papules, vesicles, oozing, crusting, or chronic scaling and lichenification. Diagnosis is based on the distribution and exposure history and may be confirmed by patch testing. See \hyperlink{Contact Dermatitis}{contact dermatitis}.

\medterm{Alopecia}-- Hair loss from any cause. Androgenetic alopecia (male/female pattern hair loss): most common, related to DHT sensitivity. Alopecia areata: autoimmune, patchy hair loss with potential progression to totalis (scalp) or universalis (entire body). Telogen effluvium: diffuse shedding 2-3 months after a stressor (illness, childbirth, weight loss, medication). Traction alopecia: from tight hairstyles. Scarring alopecias: lichen planopilaris, discoid lupus, frontal fibrosing alopecia. Treatment: minoxidil, finasteride, corticosteroids, JAK inhibitors (for alopecia areata), and hair transplantation.

\medterm{Alopecia Areata}-- A common, autoimmune, non-scarring hair loss disorder affecting approximately 1 to 2 percent of the population worldwide, with a lifetime risk of 2 percent, occurring equally across sexes and ethnicities, with peak onset between 15 and 29 years. Pathogenesis involves loss of immune privilege of the hair follicle: in normal physiology, the hair follicle bulb maintains an immune-privileged microenvironment (low MHC class I expression, production of TGF-beta, alpha-MSH, and IL-15), leading to collapse of the immune-privileged state, with CD8+ cytotoxic T cells attacking the hair follicle bulb (the anagen-stage hair follicle), causing premature catagen entry and cessation of hair growth. The condition presents as well-circumscribed, round to oval, smooth patches of non-scarring hair loss on the scalp or any body hair, and severity ranges from a single patch (patchy AA) to complete scalp hair loss (alopecia totalis) or total body hair loss (alopecia universalis).

\medterm{Amelanotic melanoma}-- A variant of melanoma that lacks significant melanin production, presenting as a pink, red, or flesh-coloured papule or nodule rather than a typical dark lesion. Its paucity of pigment can delay diagnosis, making it a diagnostic challenge. Histopathology with immunohistochemical markers confirms the diagnosis.

\medterm{Amiodarone-induced skin discoloration}-- A slate-grey to bluish hyperpigmentation of sun-exposed skin caused by long-term amiodarone therapy, due to deposition of the drug and its metabolites in dermal macrophages. It is dose-related, more pronounced in photodistributed areas, and may be partially reversible upon drug cessation. Ocular and pulmonary toxicity should also be monitored.

\medterm{Amyloid purpura}-- Purpuric discoloration of the eyelids and periorbital skin (raccoon eyes) caused by cutaneous amyloid deposition, which weakens vessel walls and leads to extravasation of erythrocytes. It may be a sign of systemic amyloidosis, particularly AL amyloidosis, and warrants further workup including serum free light chains and tissue biopsy. The lesions are non-blanching and may extend to other sun-exposed areas.

\medterm{Androgenetic Alopecia}-- The most common cause of hair loss, caused by progressive miniaturization of genetically susceptible hair follicles under androgen influence. Male pattern commonly affects the frontal hairline and vertex; female pattern commonly causes diffuse crown thinning with relative preservation of the frontal hairline. Treatment is individualized and may include minoxidil, selected antiandrogens, or hair transplantation; finasteride and dutasteride require particular pregnancy-related precautions.

\medterm{Angioedema}-- Rapid swelling of the deep dermis/subcutaneous and submucosal tissues (lips, eyelids, tongue, throat, hands, GI tract) due to mast cell mediator release (histamine—allergic/idiopathic, urticaria associated) or bradykinin-mediated (hereditary angioedema—C1 esterase inhibitor deficiency; ACE-inhibitor-induced). Features: non-pitting, non-pruritic swelling; airway involvement (laryngeal edema) is life-threatening. Diagnosis: clinical; C4, C1-INH level/function for hereditary; drug history (ACE inhibitors). Management: acute—airway management, antihistamines/corticosteroids/epinephrine for histaminergic; for bradykinin: C1-INH concentrate, icatibant, ecallantide; stop ACE inhibitors; prophylaxis for hereditary (danazol, C1-INH, lanadelumab).

\medterm{Anti-MDA5 dermatomyositis}-- A clinically distinct subtype of dermatomyositis associated with anti-melanoma differentiation-associated gene 5 antibodies, characterised by skin-predominant disease, mechanic's hands, and a high risk of rapidly progressive interstitial lung disease. Cutaneous manifestations include diffuse papules, mucocutaneous ulceration, and palmar papules. Early recognition and aggressive immunosuppression are critical due to the risk of fatal pulmonary complications.

\medterm{Aquagenic keratoderma}-- Transient, painful whitish thickening and wrinkling of the palms and soles occurring within minutes of water immersion, accompanied by a prickling or burning sensation. It is often associated with cystic fibrosis or other sweat gland disorders but can also be idiopathic or drug-induced. Symptoms resolve shortly after drying the skin.

\medterm{Athlete's Foot}-- A common superficial fungal infection of the plantar and interdigital foot skin (tinea pedis), most frequently caused by anthropophilic dermatophytes including \textit{Trichophyton rubrum}, \textit{Trichophyton mentagrophytes} var. \textit{interdigitale}, and \textit{Epidermophyton floccosum}. Transmission occurs through direct contact with desquamated keratinocytes shed on warm, moist surfaces in public showers, locker rooms, and swimming pools. Clinical variants comprise: (1) interdigital tinea pedis, presenting with maceration, fissuring, scaling, and intense pruritus between the toes, most commonly in the fourth interdigital space; (2) moccasin-type hyperkeratotic tinea pedis, characterized by chronic, fine, silvery-white scaling and erythema distributed over the soles and lateral margins of the feet, frequently co-occurring with onychomycosis and palmar involvement (the 'two feet, one hand' syndrome); and (3) acute vesiculobullous tinea pedis, manifesting with painful, inflammatory vesicles on the instep. Diagnosis is confirmed by potassium hydroxide (KOH) preparation demonstrating branching septate fungal hyphae. Topical allylamines (terbinafine) or azoles are curative for localized interdigital disease, whereas extensive moccasin-type infections require oral terbinafine or itraconazole.

\medterm{Atopic Dermatitis}-- A chronic, relapsing inflammatory skin disease characterized by pruritus, xerosis, and eczematous lesions. It reflects epidermal-barrier dysfunction and immune dysregulation, with flares triggered by irritants, allergens, infection, stress, or climate. It commonly begins in childhood and may be associated with food allergy, asthma, and allergic rhinitis. Treatment includes skin-barrier care, topical anti-inflammatory therapy, trigger reduction, and specialist-directed systemic therapy for severe disease.

\medterm{Bacterial dermatitis}-- Inflammation of the skin caused by bacterial infection, typically presenting with erythema, edema, warmth, pain, and purulent discharge. Common causative organisms include Staphylococcus aureus and Streptococcus pyogenes, and the condition may arise de novo or as a superinfection of pre-existing dermatitis. Treatment involves topical or systemic antibiotics as appropriate.

\medterm{Balanitis}-- Inflammation of the glans penis, often involving the prepuce (balanoposthitis). Causes include poor hygiene, irritants, infection, and dermatologic conditions such as lichen planus, psoriasis, and lichen sclerosus. It may cause erythema, swelling, discharge, pruritus, and discomfort; treatment depends on the cause.

\medterm{Baldness}-- The common term for hair loss or alopecia, most frequently referring to androgenetic alopecia (male or female pattern hair loss). Affects up to 50\% of men by age 50 and 25\% of women by age 60. Driven by dihydrotestosterone (DHT)-mediated miniaturization of hair follicles in susceptible individuals. Other causes include alopecia areata, telogen effluvium, traction alopecia, and scarring alopecias. Treatments: topical minoxidil, oral finasteride (men), low-level laser therapy, platelet-rich plasma injections, and hair transplantation.

\medterm{Barrier Creams}-- Topical preparations applied to the skin to create a physical protective layer against irritants, moisture, and pathogens. Common formulations include zinc oxide, petrolatum, dimethicone, and titanium dioxide. Indicated for preventing irritant contact dermatitis (occupational exposure to water, detergents, chemicals), managing incontinence-associated dermatitis, protecting periwound skin from maceration, and preventing moisture-associated skin damage. Barrier creams complement but do not replace regular cleansing and moisturizing.

\medterm{Basal Cell Carcinoma}-- The most common malignant skin tumor, arising from basal keratinocyte-lineage cells and associated with cumulative ultraviolet exposure and Hedgehog-pathway abnormalities. It commonly presents as a pearly papule or nodule with telangiectasia, ulceration, or a rolled border. It grows locally and rarely metastasizes; diagnosis requires biopsy and treatment is selected according to subtype, size, site, and recurrence risk.

\medterm{Beau Lines}-- Transverse, horizontal grooves or band-like depressions spanning the entire width of the nail plate, resulting from a temporary, acute arrest or profound deceleration of mitotic activity and protein synthesis in the proximal nail matrix. The width and depth of the furrow reflect the duration and severity of the matrix insult, while its anatomical distance from the proximal nail fold allows chronological estimation of the inciting physiological event (based on normal fingernail growth rates of approximately $2\text{--}3\text{ mm/month}$). Beau lines appear symmetrically across multiple digits following severe systemic physiological stress, including prolonged high-grade fevers, septic shock, major surgical trauma, myocardial infarction, acute severe viral illnesses (such as COVID-19 or hand-foot-and-mouth disease), severe malnutrition, or systemic cytotoxic chemotherapy. Isolated involvement of a single digit indicates localized trauma or acute paronychia. When the matrix arrest is complete and prolonged across all digits, the furrow may penetrate the entire thickness of the plate, leading to full proximal detachment and shedding of the nail (onychomadesis). Management is entirely directed at addressing the underlying systemic illness, as the lines gradually grow out distally and resolve spontaneously without specific intervention.

\medterm{Becker’s nevus}-- A benign hamartomatous lesion appearing as a unilateral, hyperpigmented, often hairy patch, typically developing in adolescence on the shoulder or upper trunk of males. It results from a localized androgen-dependent increase in dermal smooth muscle and pigmentation. It is cosmetically concerning but carries no malignant potential.

\medterm{Bedbug Bites}-- Cutaneous reactions to bites of Cimex lectularius (bedbugs), nocturnal blood-feeding insects. Bites are painless but cause erythematous, pruritic papules, often in a linear arrangement (breakfast, lunch, dinner sign) or clusters on exposed skin (face, neck, arms, hands). Delayed hypersensitivity develops over days. Some have bullous or urticarial reactions. Diagnosis: clinical; finding bedbugs or blood spots on sheets. Treatment: symptomatic (topical corticosteroids, oral antihistamines). Secondary infection may require antibiotics. Management: professional pest control, wash bedding in hot water, encase mattresses.

\medterm{Bedsore}-- Pressure ulcer (decubitus): localized skin/soft tissue injury from prolonged pressure over bony prominences (sacrum, heels, occiput, trochanters), usually in immobile patients. Stages: 1 (non-blanching erythema), 2 (partial-thickness skin loss/abrasion/blister), 3 (full-thickness to subcutaneous fat), 4 (deep to muscle/bone); unstageable with eschar. Risk: immobility, incontinence, malnutrition, sensory loss, diabetes, poor perfusion. Complications: infection (cellulitis, osteomyelitis, sepsis). Prevention: repositioning (2-hourly), pressure-relieving surfaces, nutrition, skin care. Management: offloading, wound care (debridement, dressings), infection control, and surgical repair for advanced ulcers.

\medterm{Biologic Therapy in Dermatology}-- Treatment with targeted biological agents, usually monoclonal antibodies or receptor proteins, that inhibit selected immune pathways. Biologics are used for diseases such as psoriasis, atopic dermatitis, hidradenitis suppurativa, and urticaria; infection risk, vaccination status, screening, and monitoring must be considered.

\medterm{Birthmarks}-- Congenital or neonatal cutaneous lesions presenting as circumscribed discolorations, vascular malformations, or hamartomatous structural abnormalities of the skin, broadly classified into vascular anomalies and pigmented melanocytic or epidermal lesions. Vascular birthmarks include infantile hemangiomas (benign vascular tumors that express GLUT1, proliferate rapidly during the first months of life, and subsequently undergo slow spontaneous involution, treatable with oral propranolol when ulcerated or visually threatening) and vascular malformations (e.g., port-wine stains or capillary malformations, which persist and grow proportionally with the child and may signify Sturge-Weber syndrome when involving the ophthalmic $V_1$ trigeminal distribution). Pigmented birthmarks encompass congenital melanocytic nevi (proliferations of melanocytes requiring serial dermatoscopic surveillance due to elevated melanoma risks in large lesions $>20\text{ cm}$), dermal melanocytosis (Mongolian spots), café-au-lait macules (suggestive of neurofibromatosis type 1 when six or more are present), and sebaceous or epidermal nevi. Diagnostic evaluation relies on clinical morphology, dermoscopy, and selective ultrasound or MRI to rule out associated syndromic neurocutaneous disorders.

\medterm{Birt-Hogg-Dube Syndrome}-- Autosomal dominant disorder caused by mutations in the FLCN gene (17p11.2) encoding folliculin. Classical triad of cutaneous fibrofolliculomas (multiple small flesh-colored papules on face, neck, upper chest), pulmonary cysts predisposing to spontaneous pneumothorax, and increased risk of bilateral, multifocal renal tumors (chromophobe renal cell carcinoma, hybrid oncocytic tumors, oncocytoma). Most patients reach adulthood before manifestations appear. Diagnosis: clinical fibrofolliculomas with histologic confirmation, family history, FLCN genetic testing. Surveillance: annual/periodic renal MRI from age 18-20; baseline chest CT. Treatment: surgical excision of skin lesions; chest tube for pneumothoraces; surgical management of renal tumors; genetic counseling of at-risk relatives.

\medterm{Blackheads}-- Open comedones; non-inflammatory primary lesions of acne vulgaris resulting from the impaction and distension of the pilosebaceous follicular infundibulum by an adherent plug of desquamated keratinocytes and oxidized sebum. The follicular orifice remains dilated and patent to the skin surface, exposing the compacted lipid-keratin material to atmospheric oxygen. The characteristic dark black or brown coloration is caused by the oxidation of melanin granules and lipid components within the plug, rather than trapped environmental dirt. Open comedones represent a key morphologic stage in acne pathophysiology, creating an anaerobic, lipid-rich microenvironment within the deeper follicle that permits the proliferation of \textit{Cutibacterium acnes}, potentially triggering follicular wall rupture and progressing into inflammatory papules, pustules, or nodules. Medical management centers on topical retinoids (tretinoin, adapalene, tazarotene), which normalize follicular epithelial desquamation and loosen existing microcomedones, often combined with topical keratolytic agents such as salicylic acid, azelaic acid, or alpha-hydroxy acids.

\medterm{Blastomyces dermatitidis}-- A dimorphic fungus that causes blastomycosis, a systemic pyogranulomatous infection acquired by inhalation of conidia from soil in endemic regions such as the Ohio and Mississippi River valleys. Cutaneous manifestations include verrucous or ulcerated plaques, and pulmonary involvement is common. Diagnosis requires identification of broad-based budding yeast in tissue.

\medterm{Body Lice}-- Cutaneous infestation caused by \textit{Pediculus humanus humanus} (the human body louse), an obligate hematophagous ectoparasite that resides, feeds, and lays eggs (nits) primarily in the seams and folds of clothing and bedding rather than directly on human host hair. Body lice classically afflict populations living under conditions of severe poverty, homelessness, crowding, social displacement, or refugee crises where laundering facilities and regular hygiene are unavailable. Clinical manifestations result from hypersensitivity reactions to louse saliva, causing intense nocturnal pruritus, linear excoriations, erythematous macules, wheals, and, in chronic cases, marked lichenification and post-inflammatory hyperpigmentation on the trunk, axillae, and waist (termed vagabond's disease). Medically, the body louse is of critical epidemiological significance as the biological vector for three life-threatening systemic bacterial infections: epidemic typhus (\textit{Rickettsia prowazekii}), trench fever (\textit{Bartonella quintana}), and louse-borne epidemic relapsing fever (\textit{Borrelia recurrentis}). Definitive treatment requires laundering and high-heat machine drying of clothing and bedding ($>54^\circ\text{C}$ or $>130^\circ\text{F}$ for at least 30 minutes); topical permethrin $5\%$ cream or oral ivermectin is indicated for heavily infested individuals.

\medterm{Boil}-- Furuncle: an acute, deep infection of a hair follicle (folliculitis that extends to the dermis) by Staphylococcus aureus, forming a painful, red, swollen nodule with central purulent core that may spontaneously drain. Multiple or recurrent boils = furunculosis; coalescing boils form a carbuncle (larger, deep, multiple follicles, usually on the neck/back, with fever/systemic symptoms). Risk factors: diabetes, immunosuppression, obesity, poor hygiene, skin trauma, nasal MRSA carriage. Diagnosis: clinical; culture of pus. Management: warm compresses, incision and drainage of fluctuant lesions, antibiotics for surrounding cellulitis/systemic infection or recurrent disease (cover MRSA), and hygiene measures; treat predisposing factors.

\medterm{Boil (Furuncle)}-- A painful, deep-seated infection of a hair follicle caused most commonly by Staphylococcus aureus, extending into the surrounding subcutaneous tissue. Presents as a tender, erythematous, fluctuant nodule that eventually discharges pus. When multiple adjacent follicles coalesce, it forms a carbuncle, often with systemic symptoms. Common sites: neck, face, armpits, and buttocks. Treatment: warm compresses, incision and drainage for fluctuant lesions, and antibiotics for surrounding cellulitis, large lesions, or immunocompromised patients. Recurrent furunculosis may require decolonization (mupirocin nasal ointment, chlorhexidine washes).

\medterm{Boils}-- An acute, painful, necrotizing bacterial infection of a hair follicle and adjacent perifollicular subcutaneous tissue (furuncle) that extends deeply into the reticular dermis and subcutaneous fat, predominantly caused by \textit{Staphylococcus aureus} (including community-acquired methicillin-resistant \textit{S. aureus}, CA-MRSA, frequently harboring the Panton-Valentine leukocidin toxin). Boils present initially as a firm, tender, erythematous inflammatory nodule that rapidly fluctuates, developing central necrosis and purulent pointing. Coalescence of several adjacent furuncles forms an extensive, multiloculated carbuncle, which is accompanied by systemic fever, chills, and regional lymphadenopathy. Predisposing factors include anterior nares \textit{S. aureus} colonization, diabetes mellitus, immunosuppression, obesity, and recurrent skin friction in hair-bearing sites (neck, axillae, buttocks, thighs). Management involves moist heat compresses to promote spontaneous drainage for small lesions ($<5\text{ mm}$), and surgical incision and drainage for fluctuant collections. Systemic oral antimicrobials (cephalexin, or trimethoprim-sulfamethoxazole/doxycycline for MRSA coverage) are mandated when systemic signs, surrounding cellulitis, or facial danger-triangle localization are present.

\medterm{Breslow Thickness}-- Microstaging measurement of melanoma invasion depth from the top of the granular layer of the epidermis to the deepest point of tumor invasion, in millimeters. Most important prognostic factor in primary cutaneous melanoma, used in the AJCC TNM staging system (T category). Stratification: $\le 1.0$ mm (T1, local), 1.01-2.0 mm (T2), 2.01-4.0 mm (T3), $> 4.0$ mm (T4). Combined with ulceration status, guides wide local excision margins (in situ: 0.5 cm; $\le 1$ mm: 1 cm; 1-2 mm: 1-2 cm; 2-4 mm: 2 cm; $> 4$ mm: 2-5 cm) and sentinel lymph node biopsy decision (offered if $> 0.8$ mm or $> 1$ mm with ulceration).

\medterm{Bruising}-- Ecchymosis/contusion: a purplish-blue discoloration from blood extravasation into skin/subcutaneous tissue following trauma; color changes with age (red-blue $\rightarrow$ purple $\rightarrow$ green/yellow—bilirubin breakdown). Spontaneous or easy bruising suggests: bleeding disorders (thrombocytopenia, hemophilia, von Willebrand), anticoagulant/antiplatelet use, corticosteroid/aging (skin fragility), liver disease, vitamin C/K deficiency, or physical abuse (pattern—hand/finger, belt marks, in children). Small punctate bruises = petechiae (thrombocytopenia, vasculitis). Diagnosis: history (including medications and injury mechanism), exam, and targeted labs (CBC, coagulation, von Willebrand panel) for unexplained bruising. Management: treat cause, protect skin.

\medterm{Bullae}-- Large (>1 cm) fluid-filled blisters; see also vesicle (<1 cm). Classification by level of skin cleavage (important for diagnosis): subcorneal (impetigo, pustular psoriasis), intraepidermal (pemphigus—flaccid bullae, Nikolsky positive; staphylococcal scalded skin), subepidermal (bullous pemphigoid—tense bullae, dermatitis herpetiformis, porphyria cutanea tarda, erythema multiforme, drug reactions). Causes: autoimmune, infection (bullous impetigo, SJS/TEN), drugs, trauma (friction blisters), and inherited (epidermolysis bullosa). Diagnosis: clinical, Nikolsky sign, biopsy with immunofluorescence. Management: treat cause—immunosuppression for autoimmune; wound care; avoid rupture/ infection.

\medterm{Bullous Impetigo}-- A superficial staphylococcal skin infection caused by exfoliative toxins that produce fragile, flaccid bullae and honey-coloured or varnish-like crusts, especially in infants and young children. Diagnosis is usually clinical; localised disease is treated with topical antibiotics and extensive disease with systemic anti-staphylococcal therapy. It is contagious and requires hygiene measures and assessment for outbreaks.

\medterm{Bullous Pemphigoid}-- An autoimmune subepidermal blistering disease affecting elderly patients. It presents with tense blisters on erythematous or normal-appearing skin, often preceded by pruritic urticarial plaques. Autoantibodies target BP180 and BP230 (hemidesmosome proteins). Diagnosis is by direct immunofluorescence showing linear IgG and C3 at the basement membrane. Treatment includes high-potency topical or systemic corticosteroids and steroid-sparing agents (doxycycline, azathioprine, mymycophenolate mofetil). Rituximab is used for refractory disease. Prognosis is generally good but disease can be chronic with relapses.

\medterm{Cafe-au-lait Macules}-- Flat, uniformly hyperpigmented, light-brown to dark-brown (cafe-au-lait, "coffee with milk") macules, $> 5$ mm in prepubertal children or $> 15$ mm in postpubertal individuals. Most are isolated and benign; $\ge 6$ cafe-au-lait macules are a major diagnostic criterion for neurofibromatosis type 1 (NF1), and they occur in McCune-Albright syndrome (large, unilateral, segmental), Legius syndrome, tuberous sclerosis, Noonan syndrome, Fanconi anemia, and Russell-Silver syndrome. Histology: giant melanosomes with increased melanin. Wood lamp examination highlights macules. Asymptomatic; management: clinical documentation and surveillance for NF1 features (Lisch nodules, axillary freckling, neurofibromas, optic glioma).

\medterm{Calcinosis}-- The pathological deposition of calcium salts in cutaneous and subcutaneous tissues, presenting as firm, white to yellowish nodules or plaques that may ulcerate and extrude chalky material. It is commonly associated with connective tissue diseases such as dermatomyositis and systemic sclerosis, and can cause pain, limited mobility, and secondary infection. Management includes medical therapy and surgical excision of symptomatic deposits.

\medterm{Candidal Intertrigo}-- A Candida infection of skin folds (axillae, inframammary, inguinal, intergluteal, pannus) characterized by erythematous, macerated plaques with satellite pustules at the periphery. Most common in obese, diabetic, and immunocompromised patients, and with prolonged antibiotic use. Presents with itching, burning, and soreness. Diagnosis: clinical; KOH prep shows pseudohyphae and budding yeasts. Treatment: topical antifungals (clotrimazole, miconazole, nystatin); drying powders; keep skin folds dry. Severe or refractory cases: oral fluconazole. Prevention: weight loss, moisture management, glycemic control.

\medterm{Carbuncle}-- A deep, coalescent infection of multiple adjacent hair follicles by Staphylococcus aureus (usually MRSA), forming a larger, swollen, erythematous, painful mass with multiple draining sinuses, most often on the neck, back, or thighs; more serious than a furuncle, commonly with fever, malaise, and systemic infection. Risk factors: diabetes, immunosuppression, obesity. Complications: cellulitis, abscess, bacteremia, sepsis. Diagnosis: clinical; cultures (blood, pus). Management: incision and drainage, systemic antibiotics (cover MRSA), glycemic control, and treatment of risk factors.

\medterm{Cellulitis}-- An acute, spreading bacterial infection of the deep dermis and subcutaneous tissue, most commonly caused by beta-hemolytic streptococci (particularly Streptococcus pyogenes) or Staphylococcus aureus. It typically presents as an ill-demarcated, expanding area of erythema, warmth, edema, and tenderness, frequently accompanied by systemic symptoms such as fever, chills, and lymphadenopathy. A portal of entry (e.g., skin fissure, tinea pedis, wound) is often present. Diagnosis is clinical. Treatment involves systemic antibiotics targeting gram-positive cocci, elevation of the affected limb, and management of predisposing factors.

\medterm{Cerebriform nevus sebaceous}-- A variant of nevus sebaceous presenting as a soft, cerebriform (brain-like), pedunculated or sessile nodule, often found on the scalp or face. It is a benign adnexal hamartoma composed of sebaceous glands, but like other nevus sebaceous lesions, it carries a low risk of secondary neoplasms and may warrant monitoring or excision.

\medterm{Chemical Peels}-- Chemical peels use acids to exfoliate skin layers. Superficial peels (glycolic acid 20 to 70 percent, salicylic acid 20 to 30 percent) affect the epidermis. Medium peels (trichloroacetic acid 35 to 50 percent) reach the papillary dermis. Deep peels (phenol 88 percent) reach the reticular dermis. Indications include acne, photoaging, hyperpigmentation, and fine wrinkles. Recovery time increases with peel depth. Post-peel care includes sun protection and emollients. Complications include infectioninfection, post-inflammatory hyperpigmentation, and scarring. Proper patient selection and pre-treatment preparation minimise adverse effects.

\medterm{Cherry Angioma}-- A common benign proliferation of small dermal blood vessels that appears as a bright-red to purple papule, usually on the trunk or limbs. Lesions may bleed after trauma but generally require no treatment unless diagnosis is uncertain or removal is desired.

\medterm{Cold Sores}-- Recurrent herpes labialis; common, localized cutaneous and mucocutaneous infections caused by Herpes Simplex Virus, predominantly HSV-1 (and less frequently HSV-2), affecting the perioral skin and vermilion border of the lips. Following primary oral inoculation in childhood (which is often subclinical or manifests as acute herpetic gingivostomatitis), the virus ascends sensory axons to establish lifelong latent infection within the trigeminal ganglion. Reactivation is triggered by psychological stress, febrile illness, immunosuppression, localized tissue trauma, menstruation, or ultraviolet light exposure. The clinical episode characteristically begins with a localized sensory prodrome of tingling, burning, or pruritus lasting 6 to 24 hours. This is rapidly followed by the eruption of grouped, tense, painful, umbilicated vesicles on an erythematous base at the vermilion-cutaneous junction ('fever blisters'). The vesicles rupture within 24 to 48 hours into shallow, painful erosions that form a golden-honey crust and heal completely without scarring over 7 to 10 days. The lesion is highly contagious from the onset of the prodrome until complete re-epithelialization. Early episodic oral antiviral therapy (valacyclovir $2\text{ g}$ twice daily for one day, or famciclovir) initiated at the earliest prodromal sign significantly reduces pain and aborts or shortens lesion duration.

\medterm{Common wart}-- Verruca vulgaris; the most frequent cutaneous manifestation of human papillomavirus (HPV) infection, most commonly caused by HPV genotypes 1, 2, 4, 27, and 57. Common warts present as discrete, firm, exophytic, hyperkeratotic papules or nodules with a rough, verrucous, 'cauliflower-like' surface, ranging from 1 mm to over 1 cm in size. They arise most frequently on the hands, periungual skin, fingers, elbows, and knees, particularly in children and adolescents, though they may occur anywhere on the cutaneous surface. A pathognomonic diagnostic sign upon paring or dermoscopy is the visualization of multiple pinhead-sized black dots, which represent thrombosed, dilated capillary loops within elongated dermal papillae. Transmission occurs via direct skin-to-skin contact or autoinoculation through minor epithelial abrasions. Histopathology exhibits marked hyperkeratosis, acanthosis, papillomatosis, and koilocytosis in the upper stratum spinosum and granulosum. In immunocompetent individuals, approximately two-thirds of common warts resolve spontaneously within two years due to acquired cell-mediated immunity. Active treatment options include daily topical salicylic acid ($17\text{--}40\%$), cryotherapy with liquid nitrogen, topical 5-fluorouracil, pulsed-dye laser, and intralesional immunotherapy (candida antigen).

\medterm{Congenital melanocytic nevus}-- A melanocytic nevus present at birth or appearing within the first few weeks of life, ranging from small (< 1.5 cm) to giant lesions covering large body surface areas. It results from proliferation of melanocytes during embryonic development and carries a small but important risk of melanoma, which is higher in larger lesions. Management depends on size, location, and risk of malignant transformation.

\medterm{Congenital nevus}-- See congenital melanocytic nevus. A melanocytic nevus present at birth or arising in early infancy, composed of proliferating melanocytes within the epidermis and dermis. Small congenital nevi are common and carry minimal malignant risk, while medium and giant lesions require closer surveillance due to increased melanoma potential.

\medterm{Contact Allergen Testing}-- Patch testing identifies allergens causing allergic contact dermatitis. Standard panel includes 35 common allergens (nickel, fragrance mix, balsam of Peru, formaldehyde, rubber accelerators, topical antibiotics, preservatives). Patches are applied to the back for 48 hours and read at 48, 72, and 96 hours. Positive reactions indicate allergens to avoid. Contact allergen databases and patient education are essential for management. False positives and negatives can occur.

\medterm{Contact Dermatitis}-- An inflammatory skin reaction caused by direct contact with an external agent,divided into two main two main types: allergic contact dermatitis (ACD, type IV delayed hypersensitivity reaction, mediated by sensitized T lymphocytes) and irritant contact dermatitis (ICD, non-immunologic, dose-dependent direct chemical or physical damage to the epidermis, accounting for 80 percent of cases). Allergic contact dermatitis requires prior sensitization to the allergen (which can occcur over hours to days, while irritant contact dermatitis occurs immediately after exposure to the inciting agent). Clinical presentation varies by duration and severity, with acute ACD presenting with erythema, vesiculation, oozing, and crusting, and chronic ACD presenting with lichenification, fissuring, scaling, and hyperpigmentation. The distribution of the rash provides important clues to the causative agent. Diagnosis is primarily clinical, supported by the history of exposure and confirmed by patch testing, which is the gold standard for identifying the specific allergen.

\medterm{Corns and Calluses}-- Corns (heloma) are focal areas of hyperkeratosis caused by pressure and friction, typically on toes. Hard corns (heloma durum) occur on dorsal toes; soft corns (heloma molle) occur between toes. Calluses (tyloma) are diffuse hyperkeratosis on weight-bearing surfaces. Treatment includes padding, shoe modification, keratolytic agents (salicylic acid), and debridement. Diabetics should avoid self-treatment due to ulcer risk.

\medterm{Cosmetic Dermatology}-- Cosmetic dermatology encompasses aesthetic procedures. Botulinum toxin injections (Botox, Dysport, Xeomin) relax facial muscles reducing wrinkles. Dermal fillers (hyaluronic acid, calcium hydroxylapatite, poly-L-lactic acid, polymethylmethacrylate) restore volume. Chemical peels (glycolic, salicylic, TCA) improve skin texture and pigmentation. Laser therapy (ablative CO2, Er:YAG, non-ablative fractional) for wrinkles, scars, and pigmentation. Intense pulsed light (IPL) for photorejuvenation.

\medterm{Cradle Cap}-- Infantile seborrheic dermatitis presenting in neonates and young infants (typically between 2 and 10 weeks of age) as thick, greasy, adherent, yellowish or brown scales and crusts overlying an erythematous base on the vertex and frontal scalp. The underlying pathogenesis is multifactorial, driven by maternal androgen stimulation of neonatal sebaceous glands, resulting in hyperseborrhea, altered sebum lipid composition, and secondary colonization by lipophilic \textit{Malassezia} yeast species (\textit{M. globosa}, \textit{M. restricta}). The condition is self-limiting and distinguished from infantile atopic dermatitis by the conspicuous absence of severe pruritus, lack of sleep disruption, and sparing of the extremities. Mild cases resolve with conservative management, including application of mineral oil or white petrolatum to soften scales, followed by gentle mechanical removal with a soft-bristle baby brush and frequent washing with a mild baby shampoo. For persistent or inflammatory lesions, short courses of low-potency topical corticosteroids (hydrocortisone $1\%$) or topical antifungal creams (ketoconazole $2\%$) are safe and highly effective.

\medterm{Crusted scabies}-- A severe, highly contagious form of scabies (Norwegian scabies) characterised by thick, hyperkeratotic crusts teeming with enormous numbers of Sarcoptes scabiei mites, typically occurring in immunocompromised, elderly, or neurologically impaired individuals. It presents with widespreadScaling, crusting, and fissuring of the skin and is a major source of institutional outbreaks. Treatment requires aggressive and repeated use of scabicides plus environmental decontamination.

\medterm{Cryotherapy}-- The therapeutic application of extreme cold (typically using liquid nitrogen at $-196^\circ\text{C}$) to destroy abnormal or diseased tissue. It is a common lesion-directed treatment in dermatology for benign lesions (seborrheic keratoses, warts, skin tags), precancerous lesions (actinic keratoses), and superficial malignancies (small superficial basal cell carcinomas). The mechanism of action involves cellular injury via intracellular and extracellular ice crystal formation, osmotic changes, and vascular stasis leading to tissue necrosis during the freeze-thaw cycle. Adverse effects include localized pain, blistering, transient or permanent hypopigmentation, and scarring.

\medterm{Cutaneous Drug Eruptions}-- a spectrum of reactions. Exanthematous (maculopapular) is the most common, appearing 4 to 14 days after drug initiation. Urticarial appears within hours. Photoallergic occurs at sun-exposed sites. Lichenoid resembles lichen planus. Pigmented purpuric dermatosis presents as cayenne pepper spots. Erythema nodosum presents as tender nodules on shins. Each pattern suggests different drug classes and mechanisms.

\medterm{Cutaneous Lupus Erythematosus}-- A diverse spectrum of autoimmune connective-tissue dermatoses characterized by photosensitive inflammatory cutaneous lesions, which may occur as an isolated skin disorder or as a prominent manifestation of systemic lupus erythematosus (SLE). Classified into three primary subtypes: (1) Acute Cutaneous Lupus Erythematosus (ACLE), manifesting as transient, non-scarring malar 'butterfly' erythema over the cheeks and nasal bridge that conspicuously spares the nasolabial folds, strongly correlated with active systemic disease and high-titer anti-dsDNA and anti-Sm antibodies; (2) Subacute Cutaneous Lupus Erythematosus (SCLE), presenting as non-scarring, annular polycyclic or papulosquamous plaques distributed across photo-exposed sites (shoulders, chest, upper back), strongly associated with anti-Ro/SSA antibodies and drug-induced etiologies (terbinafine, calcium channel blockers, thiazides); and (3) Chronic Cutaneous Lupus Erythematosus (CCLE), dominated by Discoid Lupus Erythematosus (DLE), characterized by indurated, disc-shaped erythematous plaques with follicular plugging ('carpet-tack' sign), central scarring atrophy, telangiectasia, and permanent cicatricial alopecia. Histopathology reveals vacuolar interface dermatitis, apoptotic basal keratinocytes (Civatte bodies), basement membrane thickening, and dermal mucin. Treatment requires strict broad-spectrum photoprotection, high-potency topical or intralesional corticosteroids, topical calcineurin inhibitors, and systemic antimalarials (hydroxychloroquine up to $5\text{ mg/kg/day}$ with routine ophthalmologic screening).

\medterm{Cutaneous Lymphoma}-- Cutaneous lymphomas are heterogeneous group of T-cell and B-cell malignancies primarily involving the skin. Primary cutaneous B-cell lymphomas include marginal zone lymphoma, follicle center lymphoma, and diffuse large B-cell lymphoma leg type. Primary cutaneous T-cell lymphomas include mycosis fungoides and Sézary syndrome. Diagnosis requires multiple biopsies and immunohistochemical staining. Treatment is stage-dependent, ranging from topical therapy to systemic chemotherapy and radiation.

\medterm{Cutaneous Small-Vessel Vasculitis}-- An acute or chronic inflammatory disease targeting post-capillary venules in the superficial dermis, characterized by immune-complex deposition, vessel wall destruction, and vascular occlusion. The cardinal clinical sign is palpable purpura---non-blanching, elevated, erythematous-to-violaceous papules and plaques that typically erupt symmetrically over dependent lower extremities (shins, ankles, feet) due to hydrostatic pressure. Mechanistically, type III hypersensitivity immune complexes deposit in the venular walls, triggering classical complement cascade activation ($C3a$, $C5a$), chemoattraction of polymorphonuclear neutrophils, and subsequent release of lysosomal enzymes and reactive oxygen species that induce fibrinoid necrosis. Etiologies encompass adverse drug reactions (antibiotics, NSAIDs), systemic infections (hepatitis B, hepatitis C, streptococcal pharyngitis), underlying autoimmune connective-tissue disorders (rheumatoid arthritis, SLE), IgA vasculitis (Henoch-Schönlein purpura), and hematologic neoplasms, though approximately $50\%$ are idiopathic. Skin biopsy of a fresh lesion ($<24\text{--}48\text{ hours}$) demonstrating leukocytoclasia (neutrophilic nuclear fragmentation or 'nuclear dust') and extravasated red blood cells confirms the diagnosis. Systemic evaluation must exclude renal, gastrointestinal, and pulmonary involvement. Management entails leg elevation, compression hosiery, removal of offending triggers, and systemic corticosteroids or colchicine for severe cases.

\medterm{Cutaneous T-Cell Lymphoma}-- A malignancy of skin-homing T lymphocytes. Mycosis fungoides is the most common subtype, presenting as patches, plaques, and tumors in a clothing distribution. Sézary syndrome is the leukemic variant with erythroderma, lymphadenopathy, and circulating malignant cells. Diagnosis requires serial skin biopsies and T-cell receptor gene rearrangement studies. Treatment ranges from topical steroids and phototherapy to systemic therapies and extracorporeal photopherephotopheresis. Prognosis is variable; early-stage mycosis fungoides has a favourable prognosis, while tumour-stage and Sezary syndrome have poorer outcomes.

\medterm{Dandruff and Scalp Conditions}-- Dandruff is a common scalp condition presenting as flaking, itching, and mild inflammation caused by overgrowth of the yeast Malassezia, dry skin, or seborrheic dermatitis. Other scalp conditions include psoriasis, contact dermatitis, ringworm tinea capitis, and head lice. Treatment typically involves medicated anti-dandruff shampoos containing zinc pyrithione, coal tar, salicylic acid, ketoconazole, or selenium sulfide.

\medterm{Dermal Structure}-- The dermis has two layers: papillary dermis (superficial, thin, containing capillaries, lymphatics, and Meissner corpuscles) and reticular dermis (deep, thick, containing dense irregular connective tissue, collagen, elastin, hair follicles, sweat glands, and sebaceous glands). The dermo-epidermal junction contains type IV collagen and laminin, anchoring the epidermis to the dermis. Dermis provides tensile strength, elasticity, and nourishment to the avascular epidermis.

\medterm{Dermatitis}-- Inflammation of the skin, often used interchangeably with eczema. Atopic dermatitis: chronic pruritic inflammation in flexural areas, associated with asthma and allergies, driven by skin barrier dysfunction and type 2 inflammation. Contact dermatitis: allergic (type IV hypersensitivity) or irritant (direct toxicity). Seborrheic dermatitis: erythematous patches with greasy yellow scale on scalp, face, and chest, associated with Malassezia yeast. Stasis dermatitis: due to chronic venous insufficiency, lower legs. Treatment: emollients, topical corticosteroids, calcineurin inhibitors, and avoidance of triggers.

\medterm{Dermatitis Herpetiformis}-- A chronic, intensely pruritic, autoimmune blistering skin condition characterized by grouped, symmetric erythematous papules, urticarial plaques, and vesicles occurring primarily on extensor surfaces (elbows, knees, buttocks, scalp). It is a cutaneous manifestation of gluten-sensitive enteropathy (celiac disease). Pathogenesis involves IgA autoantibodies targeting epidermal transglutaminase (eTG), leading to IgA deposition in a granular pattern within the dermal papillae. Diagnosis is confirmed by direct immunofluorescence showing granular IgA deposition at the dermo-epidermal junction. Treatment consists of a strict gluten-free diet and dapsone for rapid symptom control.

\medterm{Dermatofibroma}-- A common benign dermal nodule, typically on the lower extremities of young adults. It presents as a firm, hyperpigmented papule that dimples when pinched (Fitzpatrick sign). It is thought to result from reactive fibrosis following minor trauma or insect bites. Dermatofibromas are asymptomatic and require no treatment. Excision is performed for diagnostic uncertainty or cosmetic reasons. They may recur after excision.

\medterm{Dermatofibrosarcoma}-- Dermatofibrosarcoma protuberans is a locally aggressive fibrohistiocytic tumor of the skin. It typically presents as a slowly enlarging, indurated plaque on the trunk, which may develop nodules over time. It has a high local recurrence rate (up to 50 percent with simple excision) but low metastatic potential. Wide local excision with 2 to 3 cm margins or Mohs surgery is treatment. Imatinib is effective for unresectable or metastatic disease.

\medterm{Dermatofibrosarcoma Protuberans}-- A rare locally aggressive cutaneous sarcoma arising from dermal fibroblasts. It presents as a slow-growing, firm, reddish-purple plaque or nodule, typically on the trunk. It has a high local recurrence rate but low metastatic potential. Treatment is wide local excision with Mohs micrographic surgery or imatinib for unresectable disease. It is characterized by the COL1A1-PDGFB fusion gene.

\medterm{Dermatographism}-- Urticaria factitia: the development of linear wheals (hives) within minutes of stroking or scratching the skin, the most common form of physical urticaria. Mast cell hyper-reactivity releases histamine after minor mechanical pressure. Associated with chronic spontaneous urticaria and atopy. Features: red, itchy, raised lines following skin friction, resolving in 30-60 minutes. Diagnosis: clinical (stroking skin with a blunt object reproduces wheals). Management: avoid triggering friction/tight clothing, oral H1-antihistamines (cetirizine, fexofenadine), and reassurance; severe cases may need higher doses or add-on therapy.

\medterm{Dermatomyositis}-- An idiopathic inflammatory myopathy with characteristic cutaneous findings such as heliotrope eruption, Gottron papules, photosensitive poikiloderma, and the shawl or holster sign, accompanied by variable proximal muscle weakness. Assessment includes muscle enzymes, autoantibodies, imaging or electrophysiology, and evaluation for associated disease.

\medterm{Dermatopathology}-- The microscopic examination of skin specimens. Histopathologic features include acanthosis (epidermal thickening), spongiosis (intercellular edema), parakeratosis (retained nuclei in stratum corneum), hyperkeratosis (thickened stratum corneum), dyskeratosis (individual cell keratinization), and Civatte bodies (apoptotic keratinocytes). Special stains include PAS (fungi), Fontana-Masson (melanin), trichrome (collagen), and Congo red (amyloid). Direct immunofluorescence diagnosdiagnoses autoimmune blistering diseases (pemphigus, bullous pemphigoid, lupus erythematosus) by detecting in vivo bound immunoglobulins and complement at specific skin locations.

\medterm{Dermatophytosis}-- A widespread, superficial fungal infection of keratinized tissues (stratum corneum of the epidermis, hair shafts, and nail apparatus) caused by specialized filamentous fungi known as dermatophytes, comprising three anamorphic genera: \textit{Trichophyton}, \textit{Microsporum}, and \textit{Epidermophyton}. Dermatophytes are strictly keratinophilic and produce extracellular keratinase and protease enzymes that degrade and digest insoluble keratin, allowing mycelial invasion through the cornified outer layers while rarely penetrating living tissue due to host serum inhibitory factors and cell-mediated immunity. Classified taxonomically by natural ecological habitat as anthropophilic (human-adapted, causing mild chronic inflammation, e.g., \textit{T. rubrum}), zoophilic (animal hosts, inducing vigorous inflammatory responses in humans, e.g., \textit{M. canis}, \textit{T. verrucosum}), or geophilic (soil-dwelling, e.g., \textit{Nannizzia gypsea}). Clinical disease is categorized by anatomical site: tinea pedis (feet), tinea unguium (nails/onychomycosis), tinea corporis (body), tinea cruris (groin), tinea capitis (scalp), tinea barbae (beard), and tinea manuum (hands). Diagnosis is established by potassium hydroxide (KOH) microscopy showing septate branching hyphae, fungal culture, or molecular PCR. Treatment utilizes topical allylamines (terbinafine) or azoles for localized epidermal disease, and oral systemic antifungals (terbinafine, itraconazole, fluconazole, griseofulvin) for nail, hair, or widespread infections.

\medterm{Dermoscopy}-- Noninvasive examination of skin lesions with magnification and polarized or nonpolarized illumination that reveals structures not visible to the unaided eye. It improves assessment of pigmented lesions and selected inflammatory or infectious conditions but does not replace biopsy when malignancy is suspected.

\medterm{Diaper Rash}-- Irritant inflammation of the skin covered by a diaper, commonly caused by prolonged contact with urine or stool, friction, moisture, or Candida infection. It causes redness and soreness; persistent, severe, or spreading rash requires evaluation for infection or another skin condition.

\medterm{Digital Mucous Cyst}-- A benign ganglion cyst arising from the distal interphalangeal joint, presenting as a translucent, dome-shaped nodule on the dorsal aspect of the finger near the nail fold. It is associated with osteoarthritis of the DIP joint. Treatment includes observation, intralesional corticosteroid injection, cryotherapy, and surgical excision with joint capsule removal to prevent recurrence.

\medterm{DRESS Syndrome}-- Drug Reaction with Eosinophilia and Systemic Symptoms, a severe, potentially life-threatening adverse drug reaction. Typically occurs 2-8 weeks after drug initiation (latency longer than SJS/TEN). Causative drugs: anticonvulsants (phenytoin, carbamazepine, lamotrigine), allopurinol, sulfonamides, dapsone, minocycline. Presents with fever, facial edema, generalized rash (morbilliform, exfoliative dermatitis), lymphadenopathy, eosinophilia, and atypical lymphocytosis. Internal organ involvement: hepatitis (most common), nephritis, pneumonitis, myocarditis, thyroiditis. Diagnosis: RegiSCAR criteria. Treatment: immediate drug withdrawal, systemic corticosteroids, supportive care. Mortality 5-10\%.

\medterm{Drug Hypersensitivity Syndrome}-- A severe delayed drug reaction occurring 2 to 8 weeks after drug initiation. It presents with fever, rash (morbilliform, facial edema, exfoliative dermatitis), lymphadenopathy, eosinophilia, and internal organ involvement (hepatitis most common, also nephritis, pneumonitis, myocarditis). Culprit drugs include anticonvulsants (carbamazepine, phenytoin, lamotrigine), allopurinol, sulfonamides, and minominocycline, typically associated with a latency period of 2 to 8 weeks. Diagnosis is based on the RegiSCAR criteria. Treatment consists of immediate withdrawal of the culprit drug, systemic corticosteroids (prednisolone 0.5 to 2 mg/kg/day) with slow tapering over weeks to months, and supportive care including wound care, hydration, and treatment of organ-specific complications. The mortality rate is approximately 10 percent.

\medterm{Dry Skin}-- Xerosis cutis; a prevalent dermatologic condition characterized by rough, scaly, dull, and pruritic skin resulting from a critical reduction in the water content of the stratum corneum (normally $10\text{--}20\%$, falling below $10\%$). Pathophysiologically, dry skin arises from disruption of the epidermal barrier, marked depletion of intercellular lipid bilayers (ceramides, cholesterol, and free fatty acids), and deficiencies in natural moisturizing factors (NMF: urea, pyrrolidone carboxylic acid, lactic acid, and amino acids). The condition is exacerbated by environmental extremes (cold temperatures, low ambient humidity, central heating), frequent exposure to harsh alkaline soaps and hot water, biological aging (asteatosis or senile xerosis), and systemic illnesses such as hypothyroidism, chronic renal failure, and atopic diathesis. Complications include painful cutaneous fissuring, secondary bacterial infection, and asteatotic eczema (eczema craquelé). Treatment requires frequent, liberal application of barrier-repair emollients containing ceramides, humectants (urea $5\text{--}10\%$, glycerin), and occlusive agents (petrolatum), alongside lukewarm bathing and soap-free syndet cleansers.

\medterm{Dysplastic Nevus}-- Atypical nevus: a melanocytic lesion with clinical and histologic atypia, larger (>5-6 mm), with irregular borders, variable color (pink-tan-brown), and a macular/papular component; often multiple in the dysplastic nevus syndrome (familial atypical mole-melanoma, FAMMM). Individuals with dysplastic nevi have increased melanoma risk (especially with family history). Diagnosis: dermoscopy and histology (architectural disorder, cytologic atypia, lentiginous proliferation). Management: close surveillance with total body photography/dermoscopy, excision of suspicious lesions, and sun protection; risk stratification by number and family history.

\medterm{Ectodermal Dysplasia}-- A large, heterogeneous group of over 200 distinct inherited genetic disorders characterized by primary developmental abnormalities and structural defects involving two or more anatomical tissues derived from the embryonic ectoderm: namely, hair, teeth, nails, skin, and eccrine sweat glands. The most prevalent and clinically significant subtype is X-linked Hypohidrotic Ectodermal Dysplasia (Christ-Siemens-Touraine syndrome), caused by loss-of-function pathogenic variants in the \textit{EDA} gene (encoding ectodysplasin A) or its receptor \textit{EDAR}. Affected hemizygous males present with the classic clinical triad: (1) Anhidrosis or severe hypohidrosis (marked reduction or complete absence of functional eccrine sweat glands), causing severe heat intolerance and life-threatening hyperpyrexia in infancy; (2) Hypotrichosis (sparse, fine, slow-growing, lightly pigmented scalp hair and absent eyebrows and eyelashes); and (3) Hypodontia or anodontia (severely delayed dentition with conical, peg-shaped, or missing primary and secondary teeth). Facial dysmorphism includes frontal bossing, a saddle nose, and periorbital hyperpigmentation with fine wrinkling. Hidrotic ectodermal dysplasia (Clouston syndrome, caused by \textit{GJB6} mutations) spares sweat glands, presenting with severe nail dystrophy, alopecia, and palmoplantar keratoderma. Management requires lifelong multidisciplinary care: aggressive temperature regulation and cooling measures, pediatric dental prostheses, artificial tears for lacrimal gland hypoplasia, and dermatologic skin barrier care.

\medterm{Eczema}-- A group of inflammatory skin conditions characterized by pruritic, erythematous, scaly patches. Atopic dermatitis is the most common form. Nummular eczema presents as coin-shaped lesions. Dyshidrotic eczema affects palms and soles with vesicles. Eczema herpeticum is a disseminated HSV infection in eczema patients, a medical emergency. Stasis dermatitis occurs in lower extremities with venous insufficiency. Treatment includes emollients, topical corticosteroids, calcineurin inhibitors (tacrolimus, pimecrolimus), and systemic immunosuppressants for severe cases.

\medterm{Eczema (Atopic Dermatitis)}-- A chronic inflammatory skin condition characterized by dry, itchy, and inflamed patches of skin. It is driven by a combination of genetic predisposition, immune system dysfunction, and environmental triggers. Management involves moisturizers, topical corticosteroids, avoidance of irritants, and sometimes systemic medications for moderate to severe cases.

\medterm{Epidermal Inclusion Cyst}-- A common benign cyst lined by stratified squamous epithelium with keratin debris. It presents as a firm, mobile, subcutaneous nodule, typically on the face, neck, and back. A central punctum is characteristic. It may become inflamed or infected. Treatment includes complete excision (including the cyst wall) to prevent recurrence. Incision and drainage alone leads to recurrence.

\medterm{Epidermal Layers}-- The epidermis has five layers from deep to superficial: stratum basale (germinative layer with keratinocyte stem cells and melanocytes), stratum spinosum (prickle cell layer with desmosomes), stratum granulosum (keratohyalin granules), stratum lucidum (transparent layer, only in thick skin), and stratum corneum (dead, flattened keratinocytes providing barrier function). The epidermis renews every 28 to 40 days. Keratinocytes, melanocytes, Langerhans cells, and Merkel cells are the four main epidepidermal cell types in the epidermis. The basement membrane zone anchors the epidermis to the dermis and contains type IV collagen, laminin, and integrins.

\medterm{Epidermoid Cyst}-- The most common benign cutaneous cyst, also termed an epidermal inclusion cyst, epidermal cyst, or infundibular cyst, formed by the proliferation and cystic encapsulation of stratified squamous epithelium within the dermis, originating predominantly from the infundibulum of the pilosebaceous follicle. Epidermoid cysts manifest as slow-growing, painless, firm, mobile, dome-shaped subcutaneous nodules typically measuring from 1 to 5 cm in diameter, most commonly located on the face, neck, retroauricular regions, chest, and back. A pathognomonic diagnostic sign on inspection is the presence of a central comedo-like punctum representing the obstructed follicular orifice. Histopathologically, the cyst is lined by true stratified squamous epithelium complete with a distinct stratum granulosum, which continuously sheds lamellated flakes of keratin into the cystic cavity. If the cyst wall ruptures traumatically or spontaneously, keratin extrudes into the surrounding dermis, triggering a severe, acute, sterile foreign-body granulomatous inflammatory reaction characterized by sudden erythema, pain, swelling, and purulent drainage, which is frequently mistaken for a bacterial abscess. Uninflamed cysts are cured by elective surgical excision with complete removal of the intact cyst wall; inflamed cysts are managed with warm compresses, intralesional triamcinolone, and incision and drainage only when acutely fluctuant.

\medterm{Epidermolysis Bullosa}-- A group of genetic disorders causing skin fragility and blistering with minimal trauma. Simplex (intraepidermal blistering), junctional (blistering at lamina lucida), dystrophic (sublamina densa blistering), and Kindler syndrome (mixed). Dystrophic EB can cause scarring, contractures, and squamous cell carcinoma. Management includes wound care, pain control, nutritional support, and genetic counseling. New treatments include gene therapy and protein replacement.

\medterm{Erosion}-- A secondary cutaneous lesion defined as a superficial, circumscribed loss of part or all of the epidermal layer that does not extend beyond the dermal-epidermal junction into the underlying dermis. Because the basal lamina and hair follicle dermal sheaths remain intact, erosions heal entirely by re-epithelialization from surrounding epidermal keratinocytes and residual adnexal epithelia without forming a permanent scar. Clinically, erosions present as moist, red, glistening, well-demarcated shallow excavations that frequently exude clear serous fluid and form secondary crusts during healing. Pathophysiologically, erosions arise from the unroofing and rupture of intraepidermal or subcorneal blisters (as seen in pemphigus vulgaris, bullous impetigo, and herpes simplex), physical mechanical friction (abrasions, intense excoriation in atopic dermatitis), or acute necrotizing epidermal detachment (Stevens-Johnson syndrome, staphylococcal scalded skin syndrome). They are distinguished from cutaneous ulcers, which involve full-thickness loss of epidermis and dermis and inevitably heal with cicatrization. Clinical management involves non-adherent wound dressings, gentle saline cleansing, barrier ointments (white petrolatum), and targeted medical therapy for the primary bullous or inflammatory dermatosis.

\medterm{Eruptive xanthomas}-- Sudden eruption of numerous small, yellow-orange papules with an erythematous halo, typically on the buttocks, elbows, knees, and trunk, caused by severe hypertriglyceridemia. They represent lipid-laden macrophages in the dermis and are a marker of significantly elevated serum triglycerides, often exceeding 1000 mg/dL. Treatment of the underlying dyslipidemia leads to regression of the lesions.

\medterm{Erysipelas}-- A superficial form of cellulitis involving the upper dermis and superficial lymphatics, most commonly caused by Streptococcus pyogenes. It presents as a bright red, well-demarcated, raised plaque with sharp borders, typically on the face or lower extremities. Fever and chills precede the rash. Complications include abscess, bacteremia, and recurrent episodes. Treatment is penicillin or amoxicillin. Recurrent erysipelas can cause lymphedema (elephantiasis nostras).

\medterm{Erythema}-- Redness of the skin or mucous membranes caused by increased blood flow to dilated superficial capillaries, often a sign of inflammation, infection, or other underlying conditions. Erythema blanch on pressure (distinguishes it from purpura, which does not blanch). Causes: localized inflammation (cellulitis, dermatitis, insect bite, sunburn), systemic conditions (fever, polycythemia, carcinoid syndrome, rosacea), drug reactions, and autoimmune disease (lupus, dermatomyositis). Specific patterns include erythema multiforme (target lesions, HSV-related), erythema nodosum (painful red nodules on shins, panniculitis), erythema migrans (Lyme disease--expanding annular rash with central clearing), and erythema ab igne (reticulated hyperpigmentation from chronic heat exposure).

\medterm{Erythema ab igne}-- A reticulated, hyperpigmented erythematous patch caused by chronic, repeated exposure to moderate heat or infrared radiation, such as from heating pads, laptops, or fireplaces. Histologically it shows increased melanin and dermal amyloid deposition. It was historically known as 'fire-side erythema' and may become permanent if the heat source is not removed.

\medterm{Erythema gyratum repens}-- A distinctive, rapidly moving, concentric erythematous eruption with a wood-grain or serpiginous pattern, which is almost always a cutaneous paraneoplastic marker. The lesions migrate daily and are most commonly associated with underlying malignancy, particularly lung cancer. Prompt investigation for an internal malignancy is essential upon recognition.

\medterm{Erythema marginatum}-- A migratory, non-pruritic, erythematous rash with a serpiginous or annular pattern that blanches centrally, occurring as a major Jones criterion in acute rheumatic fever. It typically appears on the trunk and proximal extremities and may come and go over days or weeks. Identification requires clinical suspicion and correlation with evidence of preceding streptococcal infection.

\medterm{Erythema Multiforme}-- An acute, self-limited immune-mediated reaction, most commonly triggered by herpes simplex virus (90 percent of cases) and Mycoplasma pneumoniae. It presents as target lesions (iris lesions with three zones: central duskiness, middle pale edematous ring, outer erythematous halo) on extremities. Minor EM has few mucosal lesions; major EM has extensive mucosal involvement. Treatment is supportive, with acyclovir prophylaxis for HSV-triggered recurrent EM.

\medterm{Erythema Nodosum}-- The most common form of panniculitis, characterized by painful, erythematous, tender nodules and plaques on the bilateral shins, representing an inflammatory reaction of the subcutaneous fat. The condition is a hypersensitivity reaction to a variety of inciting antigens, most commonly infections (Streptococcal pharyngitis, the most frequent cause in children; tuberculosis; coccidioidomycosis; histoplasmosis; blastomycosis; Yersinia; Salmonella; Mycoplasma; and viral infections including Epstein-Epstein-Barr virus and hepatitis B and C), medications (oral contraceptives, sulfonamides, penicillins, and NSAIDs), inflammatory bowel disease (Crohn disease and ulcerative colitis), sarcoidosis, Behcet disease, pregnancy, and malignancies (lymphoma and leukaemia). The onset is acute, with the appearance of exquisitely tender, warm, erythematous nodules, approximately 1 to 10 cm in diameter, typically on the pretibial surfaces. Individual lesions undergo an evolution of color changes resembling a bruise, resolving over 2 to 6 weeks without ulceration or scarring. Recurrence is common, particularly with streptococcal infections.

\medterm{Erythrasma}-- A superficial infection caused by Corynebacterium minutissimum in moist intertriginous areas, producing well-demarcated reddish-brown patches with fine scale and sometimes mild pruritus. Wood-lamp examination may show coral-red fluorescence. Treatment is directed at the extent and site of disease and may include topical or systemic antibacterial therapy.

\medterm{Erythrodermic psoriasis}-- A severe, life-threatening variant of psoriasis characterised by widespread erythema and scaling involving more than 90\% of the body surface area, often with systemic symptoms including fever, chills, and hemodynamic instability. It can arise de novo, from abrupt withdrawal of systemic therapy, or as a flare of pre-existing plaque psoriasis. Emergency management includes fluid resuscitation, temperature regulation, and initiation of systemic therapy.

\medterm{Exanthema}-- A widespread skin eruption, usually macular or papular, associated with an acute systemic illness. Viral infections are common causes, but drug reactions, bacterial disease, immune disorders, and other systemic conditions may also produce exanthema; fever, mucosal involvement, blistering, or skin pain may indicate a severe process.

\medterm{Exanthematous pustulosis}-- See acute generalised exanthematous pustulosis. A severe drug-induced eruption with acute onset of hundreds of small, non-follicular, sterile pustules on an erythematous base, often accompanied by fever and neutrophilia. It typically resolves within two weeks after withdrawal of the offending medication.

\medterm{Fibrocystic Breasts}-- A common, benign spectrum of physiological and histological alterations in female breast parenchyma characterized by a combination of fibrous stromal proliferation, micro- and macrocyst formation, and ductal or lobular epithelial hyperplasia driven by cyclic hormonal fluctuations (relative estrogen dominance and progesterone deficiency). In dermatologic and clinical presentations, patients present with localized breast pain (mastalgia), sensitive fullness, and palpable, tender, rubbery nodules or plaques, often worsening during the premenstrual luteal phase. Cutaneous findings may include focal skin tenderness, localized edema, or inflammatory erythema, which must be carefully distinguished clinically from primary cutaneous pathologies (such as hidradenitis suppurativa of the inframammary fold, intertrigo, and epidermal inclusion cysts) as well as breast malignancies (inflammatory breast carcinoma and Paget disease of the nipple). Clinical breast examination, diagnostic ultrasound, and mammography (BI-RADS staging) are essential to exclude neoplasia. Management is conservative, focusing on firm supportive brassieres, nonsteroidal anti-inflammatory drugs (NSAIDs), and cyst aspiration for large, painful, fluctuant fluid collections.

\medterm{Fibrosing alopecia}-- A group of primary cicatricial (scarring) alopecias in which chronic inflammation leads to permanent destruction of hair follicles and replacement by fibrous tract remnants. It includes conditions such as lichen planopilaris, frontal fibrosing alopecia, and central centrifugal cicatricial alopecia. Early diagnosis and treatment are essential to prevent irreversible hair loss.

\medterm{First-Degree Burn}-- A superficial burn affecting only the epidermis: erythema, pain, and dryness without blistering (e.g., mild sunburn); heals in 3-7 days without scarring. Assessment: size (percent BSA—rule of nines), depth, and cause (thermal, chemical, electrical). Management: cool running water, analgesia, emollients/moisturizers, and sun protection; systemic care for large areas (fluid, analgesia). Distinguish from second-degree (blistering) and third-degree (full-thickness, insensate) burns. Prognosis excellent.

\medterm{Fixed Drug Eruption}-- A recurrent drug reaction appearing at the same site with each exposure to the offending drug. It presents as well-demarcated, round to oval, erythematous to violaceous plaques that may blister. Common culprits include sulfonamides, tetracyclines, NSAIDs, and barbiturates. After resolution, lesions leave post-inflammatory hyperpigmentation. Genitalia, lips, and hands are common sites. Treatment is identification and avoidance of the causative drug.

\medterm{Flushing}-- Transient reddening and warmth of the skin caused by increased cutaneous blood flow. It may result from heat, emotion, exertion, fever, alcohol, endocrine or neuroendocrine disease, mast-cell activation, medication exposure, or autonomic dysfunction; associated diarrhoea, wheeze, hypotension, or palpitations may indicate systemic disease.

\medterm{Folliculitis}-- Inflammation of hair follicles, presenting as pustules centered on follicles. Bacterial folliculitis (Staphylococcus aureus) is most common. Hot tub folliculitis (Pseudomonas) occurs after hot tub exposure. Gram-negative folliculitis complicates long-term antibiotic use. Pseudofolliculitis barbae (razor bumps) results from ingrown hairs. Treatment includes topical antibiotics (mupirocin), benzoyl peroxide, and avoiding shaving. Recurrent staphylococcal folliculitis may require dede-colonisation with mupirocin nasal ointment and chlorhexidine washes.

\medterm{Fungal Infection (Tinea/Dermatophytosis)}-- Superficial fungal infections of skin, hair, or nails caused by dermatophytes (Trichophyton, Microsporum, Epidermophyton). Tinea corporis (ringworm): annular scaling plaque with central clearing. Tinea pedis (athlete's foot): interdigital scaling, maceration. Tinea cruris (jock itch): groin rash sparing scrotum. Tinea capitis: scaly patches with broken hairs. Onychomycosis: thickened, discolored, dystrophic nails. Diagnosis: KOH prep, culture. Treatment: topical antifungals (terbinafine, azoles); systemic (terbinafine, itraconazole, griseofulvin) for tinea capitis and onychomycosis.

\medterm{Fungal Nail Infection}-- Onychomycosis (tinea unguium); a chronic, progressive fungal infection of the nail plate, nail bed, and surrounding periungual tissues, representing approximately $50\%$ of all nail disorders. The vast majority of cases ($>80\text{--}90\%$) are caused by anthropophilic dermatophytes, predominantly \textit{Trichophyton rubrum} and \textit{Trichophyton interdigitale}, with non-dermatophyte molds (\textit{Scopulariopsis}, \textit{Fusarium}, \textit{Aspergillus}) and \textit{Candida} species responsible for the remainder. Clinical presentations are categorized into: (1) distal and lateral subungual onychomycosis (DLSO, the most common form, exhibiting subungual hyperkeratosis, yellow-brown discoloration, and distal onycholysis); (2) superficial white onychomycosis; (3) proximal subungual onychomycosis (a classic cutaneous marker of advanced HIV infection or severe immunosuppression); and (4) total dystrophic onychomycosis. Diagnosis requires mycological confirmation (KOH microscopy, fungal culture, or periodic acid-Schiff [PAS] staining of nail clippings) before systemic pharmacotherapy. First-line medical treatment is oral terbinafine ($250\text{ mg/day}$ for 6 weeks in fingernails, 12 weeks in toenails), which requires baseline liver enzyme surveillance due to potential hepatotoxicity.

\medterm{Fungal Skin Infections}-- Skin infections caused by dermatophytes or yeasts, including ringworm (tinea), athlete's foot, and candidiasis. These infections thrive in warm, moist environments and present with red, itchy, scaly, or ring-shaped lesions. Treatment typically involves topical or oral antifungal medications and measures to keep the skin clean and dry.

\medterm{Furuncle}-- See Boil. A deep follicular infection by Staphylococcus aureus forming a painful, erythematous nodule with a central necrotic core, often requiring incision and drainage. Recurrent furunculosis warrants screening for nasal carriage and immunosuppression (diabetes). Complications: cellulitis, bacteremia, scarring. Management: warm compresses, drainage, antibiotics if extensive (flucloxacillin or cover MRSA), and good hygiene.

\medterm{Furuncle and Carbuncle}-- Furuncle (boil) is a deep infection of the hair follicle caused by Staphylococcus aureus, presenting as a painful, erythematous nodule that may drain pus. Carbuncle is a cluster of interconnected furuncles causing a larger, more deeply invasive infection. Treatment includes warm compresses, incision and drainage, and antibiotics (dicloxacillin, cephalexin, or TMP-SMX for MRSA). Recurrent furunculosis suggests carrier state or immunodeficiency.

\medterm{Furuncular myiasis}-- A cutaneous infestation caused by larvae of certain fly species (such as Dermatobia hominis or Cuterebra) that develop within boil-like furuncular lesions in the skin, presenting with a painful, erythematous nodule with a central punctum through which the larva breathes. It is endemic in Central and South America and is acquired when eggs are deposited on the skin by intermediate insect vectors. Treatment involves manual removal of the larva.

\medterm{Furunculosis}-- Recurrent or multiple furuncles, which are deep infections of hair follicles and surrounding tissue, usually caused by Staphylococcus aureus. Lesions are painful, erythematous, and purulent and may occur with colonization, close-contact transmission, skin disease, diabetes, immune dysfunction, or other risk factors.

\medterm{Generalized granuloma annulare}-- A widespread variant of granuloma annulare presenting with numerous small, flesh-coloured to erythematous papules distributed over the trunk and extremities, lacking the classic annular arrangement of localized disease. It is often associated with diabetes mellitus and may be more persistent than localized forms. Treatment options include topical corticosteroids, phototherapy, and systemic agents.

\medterm{Genodermatoses}-- Genetic skin disorders. Inherited ichthyoses (lamellar, X-linked, congenital ichthyosiform erythroderma). Epidermolysis bullosa (simplex, junctional, dystrophic). Neurocutaneous syndromes (neurofibromatosis, tuberous sclerosis, Sturge-Weber). Ectodermal dysplasias (hypohidrotic, hypotrichosis, nail dystrophy). Incontinentia pigmenti (X-linked dominant, lethal in males). Bloom syndrome, Rothmund-Thomson syndrome, and Werner syndrome are DNA repair disorders with photosensitivitphotosensitivity and increased risk of skin cancer. The RASopathies (Costello, Noonan, cardiofaciocutaneous, LEOPARD syndromes) present with characteristic facies, cardiac defects, and papillomas. Multidisciplinary management and genetic counselling are essential.

\medterm{Gingival melanoma}-- A rare and aggressive mucosal melanoma arising from melanocytes in the gingival mucosa, presenting as a pigmented or amelanotic nodule or ulcer on the gums. Mucosal melanomas of the head and neck carry a poor prognosis compared to cutaneous melanoma, often presenting at an advanced stage. Treatment involves wide surgical excision with adjuvant therapy as indicated.

\medterm{Glomus Tumor}-- A rare benign neoplasm arising from glomus bodies (thermoregulatory shunts), presenting as a small, blue-red, exquisitely tender subungual or digital nodule. The classic triad is severe pain, point tenderness, and cold sensitivity. Diagnosis is clinical with positive Love pin test and Hildreth sign (tumor disappears with ischemia). Treatment is surgical excision. Recurrence is uncommon if completely removed.

\medterm{Gottron Papules}-- Pathognomonic dermatological sign of dermatomyositis, consisting of erythematous to violaceous lichenoid papules located over the dorsal aspects of the interphalangeal (PIP/DIP) and metacarpophalangeal (MCP) joints. May undergo necrosis or ulceration. Contrast with Systemic Lupus Erythematosus (SLE) skin lesions, which spare the joints and involve the interphalangeal spaces. Associated with proximal muscle weakness, elevated creatine kinase (CK), anti-Jo-1, and anti-Mi-2 autoantibodies.

\medterm{Gout nodulosis}-- The formation of tophi—large, firm, subcutaneous nodules composed of monosodium urate crystals—resulting from chronic, poorly controlled hyperuricemia. Tophi most commonly occur at the ears, fingers, toes, olecranon bursae, and Achilles tendon, and may ulcerate and discharge chalky white material. They are a hallmark of advanced gout and indicate the need for aggressive urate-lowering therapy.

\medterm{Granuloma Annulare}-- A benign inflammatory condition characterized by annular, erythematous to skin-colored plaques with raised borders and central clearing. It most commonly affects the hands, feet, and elbows. Localized form is most common; generalized form is associated with diabetes mellitus. Diagnosis is clinical with biopsy showing palisading granulomas around necrobiotic collagen. Self-resolving within 2 years. Treatment of persistent cases includes topical or intralesional corticosteroisteroids, phototherapy, and oral corticosteroids for severe disseminated disease.

\medterm{Granulomatous amebic encephalitis}-- A rare, usually fatal central nervous system infection caused by free-living amoebae such as Acanthamoeba species, typically affecting immunocompromised individuals. It presents with a subacute to chronic meningoencephalitis and forms granulomatous lesions in the brain. Cutaneous involvement may occur as the primary portal of entry, presenting as painless skin ulcers or nodules.

\medterm{Guttate Psoriasis}-- An acute eruption of numerous small, scaly, drop-like psoriatic papules and plaques, often following streptococcal pharyngitis in children or young adults. Diagnosis is clinical; the eruption may resolve spontaneously or evolve into chronic plaque psoriasis. Treatment includes topical therapy, phototherapy, and treatment of confirmed streptococcal infection when present.

\medterm{Hair Loss}-- Alopecia: reduction of scalp or body hair. Types: nonscarring (androgenetic—pattern hair loss, telogen effluvium—stress/illness/drugs, anagen effluvium—chemotherapy, alopecia areata—autoimmune patches, tinea capitis, trichotillomania) and scarring (lupus, lichen planopilaris, folliculitis decalvans). Female pattern hair loss is common. Evaluation: history (onset, medications, stress, family), exam (pattern, scarring, hair pull test), labs (TSH, ferritin, DHEA/testosterone if indicated), biopsy for scarring. Management: treat cause; minoxidil (topical), finasteride/spironolactone (androgenetic), immunotherapy/corticosteroids (alopecia areata), antifungals (tinea), and supportive measures.

\medterm{Hair Loss (Alopecia)}-- The partial or complete loss of hair from areas where it normally grows, most commonly the scalp. Causes include genetic predisposition (androgenetic alopecia), autoimmune conditions (alopecia areata), hormonal changes, medications, stress, and nutritional deficiencies. Treatment depends on the underlying cause and may include topical minoxidil, oral finasteride, corticosteroid injections, or hair transplantation.

\medterm{Hair Loss Classification}-- Hair loss (alopecia) is classified as scarring (cicatricial) or non-scarring. Non-scarring alopecias are potentially reversible: androgenetic alopecia (male and female pattern), alopecia areata, telogen effluvium, anagen effluvium, and tinea capitis. Scarring alopecias cause permanent hair loss: lichen planopilaris, discoid lupus, frontal fibrosing alopecia, and folliculitis decalvans. Diagnosis requires clinical examination, trichoscopy, and sometimes biopsy.

\medterm{Halo nevus}-- A benign melanocytic nevus surrounded by a rim of depigmentation (Sutton nevus), caused by an immune-mediated T-cell response directed against melanocytes in the nevus. It is most common on the trunk of children and young adults and typically resolves by central depigmentation, leaving a white macule. It must be distinguished from melanoma, which can also develop a halo response.

\medterm{Henoch-Schönlein purpura}-- See IgA vasculitis. A small-vessel leukocytoclastic vasculitis characterised by palpable purpura predominantly on the lower extremities and buttocks, often accompanied by arthralgia, abdominal pain, and glomerulonephritis. It most commonly affects children following an upper respiratory infection and is mediated by IgA immune complex deposition. Most cases are self-limiting, but renal involvement requires monitoring.

\medterm{Herpes Simplex Skin Infection}-- Infection of skin or mucosa by herpes simplex virus, producing grouped vesicles or erosions on an erythematous base. Recurrent disease commonly affects the oral or genital region; widespread infection in patients with eczema is eczema herpeticum and requires urgent medical assessment.

\medterm{Herpes Zoster}-- A painful, vesicular eruption resulting from reactivation of latent varicella-zoster virus (VZV, human herpesvirus 3) in the dorsal root ganglia or cranial nerve ganglia, where the virus has remained dormant since the primary varicella (chickenpox) infection, with an incidence of approximately 1 to 3 per 1,000 person-years overall and up to 10 per 1,000 person-years in adults over 60. Risk factors include advancing age (over 50 years, 50 percent of cases), immunosupprimmunosuppression (HIV, organ transplantation, haematological malignancies, chemotherapy, and long-term corticosteroid or anti-TNF therapy). The prodromal phase (2 to 3 days) consists of headache, photophobia, and malaise, followed by the eruption phase, where painful, grouped, vesicular lesions on an erythematous base appear in a unilateral, single dermatomal distribution, most commonly the thoracic dermatomes (50 to 60 percent), followed by the trigeminal nerve distribution, particularly the ophthalmic division (V1, 12 to 15 percent, which requires urgent ophthalmology referral due to the risk of herpes zoster ophthalmicus). The acute neuritis and the zoster rash typically resolve over 2 to 4 weeks.

\medterm{Hidradenitis}-- Hidradenitis suppurativa is a chronic inflammatory follicular occlusion disease affecting apocrine gland-bearing areas (axillae, groin, inframammary). It presents with recurrent painful nodules, abscesses, sinus tracts, and scarring. Hurley staging: Stage I (single or multiple abscesses without sinus tracts), Stage II (recurrent abscesses with sinus tracts and bridging scars), Stage III (diffuse or multi-tract involvement). Smoking and obesity are risk factors. Treatment includes lifestyle modiflifestyle modification (weight loss, smoking cessation), topical clindamycin, oral antibiotics (tetracyclines, rifampicin-clindamycin combination), biologic therapy (adalimumab, the only FDA-approved biologic for HS), and surgical deroofing or excision for advanced disease.

\medterm{Hidradenitis Suppurativa}-- A chronic, inflammatory, follicular occlusion disease characterized by recurrent, painful nodules, abscesses, sinus tracts, and scarring in apocrine gland-bearing areas (axillae [most common, 70 to 80 percent], groin and inguinal folds [30 to 40 percent], perineum and perianal region [20 to 35 percent], buttocks [20 to 30 percent], and inframammary folds [10 to 20 percent]), with a prevalence of 1 to 4 percent worldwide (likely underdiagnosed, with average diagnoaverage diagnosis delay of 7-10 years). The pathogenesis begins with follicular occlusion (hyperkeratinisation of the terminal hair follicle infundibulum leading to dilation and rupture of the follicle), which triggers an intense inflammatory response with release of keratin and cell debris into the dermis, activating neutrophils, macrophages, and TH1/TH17 lymphocytes with elevated levels of TNF-alpha, IL-17, and IL-1 beta. Risk factors include smoking (the strongest environmental risk factor, present in 70 to 90 percent of patients), obesity, hormonal influences, and genetic factors (mutations in NCSTN, PSENEN, and PSEN1 encoding components of the gamma-secretase complex).

\medterm{Hives}-- Urticaria; an acute or chronic vascular reaction of the skin characterized by the rapid development of transient, intensely pruritic, raised, erythematous wheals with central pallor, occasionally accompanied by angioedema (deeper dermal, subcutaneous, or submucosal swelling). Hives are mediated by the immunological or non-immunological activation and degranulation of cutaneous mast cells and basophils, releasing histamine, leukotrienes ($LTC_4, LTD_4, LTE_4$), prostaglandin $D_2$, platelet-activating factor, and cytokines. These mediators trigger local sensory nerve stimulation (pruritus), arteriolar vasodilation (erythema), and increased post-capillary venule permeability, producing localized dermal edema (wheals). Classified as acute ($<6\text{ weeks}$, often triggered by viral illnesses, IgE-mediated food or medication allergies, or insect stings) or chronic ($\ge 6\text{ weeks}$, often autoimmune with IgG antibodies targeting $Fc\epsilon RI$ or IgE, or inducible by cold, heat, pressure, or solar stimuli). Individual hives characteristically resolve completely without residual marks within 24 hours, often shifting locations. Angioedema involving the tongue, pharynx, or larynx constitutes an airway emergency. First-line therapy is high-dose second-generation non-sedating $H_1$-antihistamines; refractory chronic cases benefit from omalizumab.

\medterm{Hives (Urticaria)}-- A common skin condition characterized by transient, well-circumscribed, intensely pruritic wheals that blanch with pressure and resolve within 24 hours without scarring, caused by mast-cell mediator release in the dermis. Acute urticaria lasts less than 6 weeks; chronic urticaria recurs for 6 weeks or longer and is often spontaneous. Triggers include infections, foods, medicines, and physical stimuli. Treatment commonly begins with a second-generation H1 antihistamine; refractory chronic disease may require specialist-directed therapy such as omalizumab.

\medterm{Hyperhidrosis}-- Excessive sweating beyond thermoregulatory needs. Primary focal hyperhidrosis affects palms, soles, axillae, and face without underlying cause. Secondary generalized hyperhidrosis is caused by medications, endocrine disorders, infections, or malignancies. Treatment includes topical aluminum chloride (first-line), oral anticholinergics (glycopyrrolate, oxybutynin), iontophoresis for palms and soles, botulinum toxin injection, and endoscopic thoracic sympathectomy for severe refrasevere refractory cases. Secondary hyperhidrosis requires treatment of the underlying cause.

\medterm{Hyperpigmentation}-- Darkening of the skin caused by increased melanin production, deposition of pigment in the dermis, or accumulation of other pigments. It may be localized or generalized and may result from ultraviolet exposure, inflammation or injury, hormonal changes, medications, endocrine disease, or systemic disorders such as haemochromatosis and Addison disease. Examples include melasma, post-inflammatory hyperpigmentation, acanthosis nigricans, and linea nigra. Evaluation is guided by onset, distribution, associated symptoms, and examination; management is directed at the underlying cause and prevention of further ultraviolet exposure.

\medterm{Hypertrichosis}-- Excessive hair growth in androgen-independent areas, affecting both sexes, either congenital or acquired. Generalized: genetic syndromes, medications (minoxidil, cyclosporine, phenytoin, glucocorticoids), malnutrition, porphyria, hypothyroidism, and malignancy (paraneoplastic). Localized: congenital nevi, hypertrichosis lanuginosa (paraneoplastic—abrupt onset of fine lanugo hair suggests GI/lung malignancy). Distinguish from hirsutism (androgen-dependent, patterned hair in women—PCOS, adrenal, ovarian causes). Evaluation: medication and family history, exam, hormone workup if hirsutism features. Management: treat cause; cosmetic hair removal (laser, electrolysis), and address psychosocial impact.

\medterm{Hypertrophic Scar}-- Raised scar tissue confined to the original wound margins, caused by excessive collagen production during healing. Unlike keloids, hypertrophic scars regress over time and do not extend beyond wound boundaries. They are common after burns, surgery, and trauma. Treatment includes pressure garments, silicone sheeting, intralesional corticosteroids, and laser therapy. Early intervention prevents progression.

\medterm{Hypohidrosis}-- Reduced or absent sweating, leading to heat intolerance and risk of hyperthermia. Causes include autonomic neuropathy (diabetes, Parkinson disease), skin diseases (scleroderma, burns), medications (anticholinergics), and genetic disorders (hypohidrotic ectodermal dysplasia). Diagnosis is confirmed by sweat test (starch-iodine test, thermoregulatory sweat test). Treatment includes cooling measures, hydration, and avoiding heat exposure.

\medterm{Ichthyosis}-- A group of disorders characterized by dry, scaly skin resembling fish scales. Ichthyosis vulgaris is the most common, autosomal dominant, associated with filaggrin mutation and atopy. X-linked ichthyosis is caused by steroid sulfatase deficiency. Lamellar ichthyosis is a severe congenital form. Treatment includes keratolytics (urea, lactic acid, ammonium lactate), emollients, and retinoids for severe cases.

\medterm{Ichthyosis Vulgaris}-- The most common inherited disorder of keratinization and epidermal cornification, transmitted in an autosomal semi-dominant pattern with an incidence of approximately 1 in 250 individuals, caused by loss-of-function mutations in the filaggrin gene (\textit{FLG}) located on chromosome 1q21. Filaggrin is an essential structural protein that aggregates keratin intermediate filaments in the stratum corneum and is enzymatically degraded into hygroscopic amino acids that constitute natural moisturizing factors (NMF). Deficiency leads to impaired corneocyte adhesion, severe barrier dysfunction, and transepidermal water loss. Clinically, ichthyosis vulgaris manifests during the first year of life with fine, polygonal, adherent whitish-gray scales that are most prominent over the extensor aspects of the extremities, particularly the pretibial shins, with characteristic sparing of the flexural creases (antecubital and popliteal fossae). Associated features include prominent hyperlinear palms and soles, keratosis pilaris, and an elevated risk for the atopic triad (atopic dermatitis, asthma, allergic rhinitis). Histology shows a diminished or absent epidermal stratum granulosum. Treatment utilizes daily keratolytics (topical urea $10\text{--}20\%$, lactic acid, alpha-hydroxy acids, salicylic acid) and emollient barrier creams.

\medterm{Immunomodulators in Dermatology}-- Topical immunomodulators treat inflammatory skin diseases without corticosteroid side effects. Pimecrolimus 1 percent cream and tacrolimus 0.03 and 0.1 percent ointments are calcineurin inhibitors used for atopic dermatitis, facial eczema, and intertriginous areas. They are steroid-sparing agents. Off-label uses include vitiligo, lichen planus, and psoriasis. JAK inhibitors (ruxolitinib cream for vitiligo and atopic dermatitis) represent a new class of topical immunomodulators.

\medterm{Immunotherapy in Dermatology}-- Cancer immunotherapy has revolutionized skin cancer treatment. Checkpoint inhibitors (pembrolizumab, nivolumab, ipilimumab) block PD-1 or CTLA-4, enhancing T-cell mediated tumor killing. They are first-line for advanced melanoma and metastatic SCC. Cemiplimab is approved for advanced BCC and SCC. Adverse events include dermatitis, colitis, hepatitis, pneumonitis, and endocrinopathies. Skin immune-related adverse events occur in 30 to 40 percent of patients and may be the first manifestation of ssystemic immunotherapy toxicity, particularly when the skin is involved as a target organ of the immune response.

\medterm{Impetigo}-- A highly contagious, superficial bacterial infection of the epidermis, primarily affecting children. Non-bullous impetigo (most common, typically caused by Staphylococcus aureus or Streptococcus pyogenes) presents as transient vesicles or pustules that rupture to form characteristic honey-colored crusts on an erythematous base, often on the face and extremities. Bullous impetigo (caused by epidermolytic toxin-producing S. aureus) presents with flaccid, clear fluid-filled bullae. Diagnosis is clinical, supported by wound culture. Treatment includes topical mupirocin for localized lesions or oral antibiotics (e.g., cephalexin) for extensive disease.

\medterm{Incontinentia Pigmenti}-- Bloch-Sulzberger syndrome; a rare, X-linked dominant genodermatosis and neurocutaneous disorder that is embryonic lethal in most hemizygous males, caused by loss-of-function mutations (most commonly an exon 4--10 genomic deletion) in the \textit{IKBKG} gene (inhibitor of nuclear factor kappa B kinase regulatory subunit gamma, also known as \textit{NEMO}) on chromosome Xq28. NEMO is essential for NF-$\kappa$B activation, which protects cells from cytokine-induced apoptosis; cellular mosaicism arising from lyonization (X-inactivation) allows affected females to survive. Cutaneous lesions evolve through four distinctive, sequential chronological stages that follow the developmental lines of Blaschko: (1) Stage 1 (vesicular stage, neonatal): linear crops of inflammatory vesicles and bullae on an erythematous base; (2) Stage 2 (verrucous stage, 2 to 6 months): linear hyperkeratotic, waxy, verrucous plaques; (3) Stage 3 (hyperpigmented stage, infancy through adolescence): striking, swirling, 'marble-cake' or reticulated slate-gray to brown hyperpigmentation along Blaschko's lines, reflecting melanin incontinence into the dermis; and (4) Stage 4 (atrophic stage, adulthood): pale, hairless, hypopigmented linear streaks or atrophic macules, especially on the calves. Extracutaneous manifestations occur in up to $80\%$ of patients and dictate overall prognosis: hypodontia and conical teeth, retinal neovascularization with potential retinal detachment (mandating urgent serial pediatric ophthalmologic examinations), and central nervous system defects (seizures, developmental delay, ischemic stroke).

\medterm{Ingrown Toenail}-- Ingrown toenail (onychocryptosis) occurs when the nail edge grows into the surrounding skin, causing pain, redness, and swelling. It is common in adolescents and adults, exacerbated by improper nail trimming, tight shoes, and trauma. Treatment includes partial nail avulsion with matrixectomy (phenol application) for recurrent cases. Conservative management includes cotton wedge under the nail edge and proper nail trimming.

\medterm{Inhalation Injury}-- Severe thermal and chemical trauma to the upper respiratory tract, tracheobronchial tree, and pulmonary parenchyma resulting from the inhalation of superheated air, steam, toxic combustion gases (carbon monoxide, hydrogen cyanide), and aerosolized particulate matter during structural or vehicle fires. In cutaneous burn management, inhalation injury is a major independent prognostic factor, significantly increasing mortality beyond that predicted by total body surface area (TBSA) cutaneous burns alone. Cutaneous and physical indicators include facial and perioral burns, singed nasal vibrissae and eyebrows, carbonaceous sputum in the pharynx, hoarseness, stridor, and a brassy cough. Superheated air causes immediate thermal edema of the supraglottic airway that can produce rapid, catastrophic airway obstruction. Inhaled chemical toxins (aldehydes, phosgene, hydrogen chloride) damage the distal respiratory epithelium, inducing ciliary paralysis, mucosal sloughing, severe bronchospasm, and non-cardiogenic pulmonary edema (ARDS). Concomitant carbon monoxide poisoning shifts the oxyhemoglobin dissociation curve to the left, while hydrogen cyanide inhibits mitochondrial cytochrome c oxidase. Management requires early prophylactic endotracheal intubation before airway edema matures, administration of $100\%$ normobaric oxygen, intravenous hydroxocobalamin, and aggressive pulmonary toilet.

\medterm{Intertrigo}-- An inflammatory rash of opposing skin surfaces (skin folds) caused by friction, moisture, heat, and lack of air circulation. Affects intertriginous areas: axillae, inframammary, inguinal, abdominal folds, and intergluteal cleft. Presents with erythema, maceration, fissuring, and erosions. Often secondarily infected with Candida albicans (satellite pustules), bacteria (erythrasma, Corynebacterium), or dermatophytes. Risk factors: obesity, diabetes, hyperhidrosis, incontinence, and poor hygiene. Diagnosis: clinical; KOH prep and Wood lamp exam for secondary infections. Treatment: keep area dry (drying powders, barrier creams), topical antifungals for Candida, low-potency corticosteroids for inflammation. Prevention: weight loss, moisture-wicking clothing.

\medterm{Isotretinoin}-- A systemic retinoid used for severe, recalcitrant nodular acne. It reduces sebum production, normalizes follicular keratinization, decreases Propionibacterium acnes, and has anti-inflammatory effects. Dose is 0.5 to 1 mg/kg/day for 15 to 20 weeks, targeting cumulative dose of 120 to 150 mg/kg. Side effects include cheilitis (almost universal), dry skin, photosensitivity, elevated triglycerides and liver enzymes, and teratogenicity (absolute contraindication in pregnancy). Two forTwo forms of contraception are required for women of childbearing potential one month before, during, and one month after therapy due to teratogenicity (iPLEDGE risk management programme in the US). Monitoring includes monthly pregnancy tests, lipid profile, and liver function tests. Other side effects include arthralgias, myalgias, headache, pseudotumor cerebri (rare), and decreased night vision. A second course may be considered after a washout period of at least 8 to 16 weeks.

\medterm{Itching}-- Pruritus, the unpleasant cutaneous sensation that provokes the desire to scratch. It may be acute or chronic and may arise from skin disease, systemic illness, medication, or neuropathic mechanisms. Treatment targets the cause and may include emollients, topical therapy, or selected antipruritic medicines; unexplained chronic pruritus requires evaluation for systemic disease.

\medterm{Juvenile Xanthogranuloma}-- The most common form of non-Langerhans-cell histiocytosis (non-LCH), representing a benign, typically self-limiting idiopathic proliferation of lipid-laden, histiocytic dendritic cells primarily affecting infants and young children (over $70\%$ presenting within the first year of life, with one-fifth present at birth). Lesions clinically present as solitary or multiple, firm, well-demarcated, dome-shaped papules or nodules that initially appear pink-to-red, rapidly acquiring a characteristic orange-yellow, xanthomatous, or 'hue-of-sunset' coloration with a glistening surface, most commonly distributed on the head, neck, and upper trunk. Histopathologically, lesions demonstrate a dense dermal infiltrate of mononuclear histiocytes, lymphocytes, eosinophils, and characteristic Touton giant cells (multinucleated giant cells with a central ring of nuclei surrounded by a peripheral rim of foamy, xanthomatous cytoplasm). Lesional cells stain positive for CD68, factor XIIIa, and fascin, but are characteristically negative for CD1a and Langerin (ruling out Langerhans cell histiocytosis). While cutaneous lesions undergo spontaneous involution over one to five years leaving an atrophic scar, screening is critical because ocular involvement (especially in the iris, which can cause hyphema and secondary glaucoma) occurs, particularly in young infants with multiple cutaneous lesions or concurrent neurofibromatosis type 1 (NF-1).

\medterm{Kaposi Sarcoma}-- A vascular tumor caused by human herpesvirus 8 (HHV-8), presenting as violaceous plaques and nodules on skin and mucous membranes. Four types: classic (elderly Mediterranean men), endemic (African), iatrogenic (immunosuppression post-transplant), and AIDS-associated. Cutaneous lesions are the most common presentation in AIDS. Treatment includes antiretroviral therapy, local therapy (radiation, cryotherapy), and chemotherapy for widespread disease.

\medterm{Keloid}-- An overgrowth of scar tissue extending beyond the original wound margins, caused by excessive collagen deposition during wound healing. Keloids are more common in darker skin types (Fitzpatrick IV to VI) and occur at sites of trauma, surgery, piercings, and burns. They are pruritic and painful. Treatment includes intralesional corticosteroid injections, silicone sheeting, cryotherapy, laser therapy, and surgical excision with adjuvant radiation. Recurrence is common.

\medterm{Keloids}-- Exuberant, benign fibroproliferative dermal scars that expand progressively beyond the original anatomical boundaries of cutaneous injury and fail to regress spontaneously over time. Keloid pathogenesis involves an exaggerated and sustained wound-healing response marked by persistent hyperactive fibroblast proliferation and massive, disorganized deposition of thick, hyalinized type I and type III collagen bundles, driven by elevated autocrine and paracrine signaling of transforming growth factor-beta (TGF-$\beta$) and vascular endothelial growth factor (VEGF). Predisposing factors include genetic susceptibility, high skin tension, and darker skin phototypes (Fitzpatrick types V and VI). Lesions commonly arise on the earlobes, presternal chest, deltoid shoulders, and upper back following ear piercing, surgical incisions, trauma, or inflammatory acne. Clinically, keloids present as firm, rubbery, shiny, pink to hyperpigmented nodules or plaques that can cause intense pruritus, tenderness, and functional impairment. First-line management entails intralesional corticosteroid injections (triamcinolone acetonide $10\text{--}40\text{ mg/mL}$) alone or combined with 5-fluorouracil, pulsed-dye laser, cryotherapy, silicone gel sheeting, and surgical excision with adjuvant postoperative superficial radiotherapy.

\medterm{Keratoacanthoma}-- A low-grade or well-differentiated variant of squamous cell carcinoma (SCC) originating from hair follicles. Characterized by rapid growth over 4-6 weeks into a dome-shaped, flesh-colored or erythematous nodular lesion with a central keratin-filled crater. Often resolves spontaneously over months via involution, leaving a depressed scar. Biopsy required to distinguish from invasive SCC (often treated with complete excisional surgery).

\medterm{Keratoderma}-- A group of disorders characterised by excessive thickening of the skin of the palms and soles, which may be hereditary or acquired. Acquired forms can be diffuse, focal, or punctate and may be associated with systemic diseases, malignancy, or medications. Treatment is directed at the underlying cause and includes keratolytics, emollients, and retinoids.

\medterm{Keratoderma blennorrhagicum}-- A keratoderma associated with reactive arthritis (formerly Reiter syndrome), characterised by waxy, yellowish, vesiculopustular papules and plaques on the palms and soles that may coalesce and extend to the extremities. It mimics pustular psoriasis histologically and clinically. Treatment includes treatment of the underlying arthritis and topical or systemic retinoids.

\medterm{Keratoderma blenorrhagicum}-- Keratoderma blennorrhagicum (or keratoderma blennorrhagica); a distinctive, severe cutaneous manifestation of reactive arthritis (historically known as Reiter syndrome), characterized by hyperkeratotic, waxy, crusty, yellow-brown papules and plaques predominantly distributed on the palms and soles. The condition is an immune-mediated reactive phenomenon triggered weeks after a distant mucosal infection, typically genitourinary (\textit{Chlamydia trachomatis}) or gastrointestinal (\textit{Salmonella}, \textit{Shigella}, \textit{Campylobacter}, \textit{Yersinia}), strongly associated with the human leukocyte antigen HLA-B27 allele. Lesions begin as clear, sterile, deep-seated vesicles on an erythematous base that rapidly evolve into pustules and subsequently desquamate into thick, hard, waxy, conical or limpet-like hyperkeratotic crusts resembling 'relief maps' on the plantar surfaces. Histopathologically, keratoderma blennorrhagicum is indistinguishable from pustular psoriasis, demonstrating marked hyperkeratosis, parakeratosis, elongation of rete ridges, and intraepidermal spongiform pustules of Kogoj filled with neutrophils. Accompanying features include circinate balanitis, urethritis, conjunctivitis, enthesitis, and asymmetric oligoarthritis. Management focuses on treating the underlying inciting infection (if active), alongside high-potency topical corticosteroids, keratolytics, topical calcipotriol, and systemic disease-modifying agents (methotrexate, biologic TNF inhibitors) for severe cases.

\medterm{Keratosis Pilaris}-- A common, benign genetic follicular condition characterized by keratotic follicular papules. Small, rough, flesh-colored or red bumps on the extensor surfaces of proximal arms, thighs, and sometimes buttocks and face. Caused by abnormal keratinization and plugging of hair follicles with retained keratin. Associated with ichthyosis vulgaris and atopic dermatitis. Onset in childhood, may improve with age. Diagnosis: clinical. Treatment: emollients with keratolytics (lactic acid, urea, salicylic acid, ammonium lactate), topical retinoids (tretinoin, adapalene), and gentle exfoliation. Laser therapy for cosmetically bothersome erythema.

\medterm{KOH Preparation}-- Potassium hydroxide preparation; a rapid, cost-effective, and indispensable point-of-care direct microscopic diagnostic test in clinical dermatology used to demonstrate fungal hyphae, pseudohyphae, and yeast cells in keratinized tissues (stratum corneum skin scrapings, subungual debris, or epilated hair shafts). The technique involves scraping the active, scaly advancing border of a suspected fungal lesion with a sterile scalpel blade onto a glass microscope slide, applying a drop of $10\text{--}20\%$ aqueous potassium hydroxide solution (often combined with dimethyl sulfoxide [DMSO] or fluorescent brighteners such as calcofluor white), and gently warming the slide without boiling. The alkaline KOH selectively dissolves and digests host keratinocytes, cellular lipids, and proteinaceous debris while leaving the rigid, chitin- and glucan-rich fungal cell walls intact. Microscopic examination under low and high power reveals distinctive fungal morphology: branching, uniform, septate hyphae with parallel walls (dermatophytes in tinea); short, stubby hyphae intermixed with round yeast clusters ('spaghetti and meatballs' appearance of \textit{Malassezia} in pityriasis versicolor); or budding yeast cells with pseudohyphae (\textit{Candida}). It provides definitive, rapid confirmation prior to initiating antifungal therapy.

\medterm{Koilonychia}-- Spoon nails: nails with a concave (spoon-shaped) contour, most commonly associated with iron deficiency anemia; also seen in hemochromatosis, Plummer-Vinson syndrome, malnutrition, and trauma (nail dystrophy). The nail becomes thin, brittle, and turns up at the edges so a drop of water rests in the concavity. Diagnosis: clinical; confirm iron status (ferritin, iron studies, CBC). Management: treat the underlying cause (iron supplementation), and nails typically recover with correction of deficiency. Consider other causes when iron studies are normal.

\medterm{Lepromatous leprosy}-- The disseminated, multibacillary pole of the leprosy spectrum characterised by poor cell-mediated immunity to Mycobacterium leprae, resulting in widespread skin nodules, plaques, and diffuse infiltration with thickening of the skin and facial features (leonine facies). High bacterial loads are present in skin and nerves, and systemic involvement including rhinitis, orchitis, and neuropathy is common. Diagnosis is confirmed by slit-skin smear or biopsy showing numerous acid-fast bacilli.

\medterm{Leukocytoclastic Vasculitis}-- A specific histopathological pattern and clinical manifestation of cutaneous small-vessel necrotizing vasculitis affecting postcapillary venules, characterized by intense neutrophilic infiltration, neutrophil apoptosis, and nuclear fragmentation (karyorrhexis, generating granular basophilic nuclear fragments termed 'nuclear dust' or leukocytoclasia). Additional diagnostic histological hallmarks include fibrinoid necrosis of the venular walls, endothelial swelling, and marked extravasation of erythrocytes into the perivascular dermis. Clinically, LCV manifests as palpable purpura---symmetrically distributed, non-blanching, elevated violaceous papules and plaques that predominantly involve dependent lower extremities (shins, ankles, feet) under hydrostatic pressure, occasionally progressing to hemorrhagic bullae, ulcerations, or digital necrosis. It is triggered by circulating antigen-antibody immune complexes (type III hypersensitivity) that activate complement and recruit neutrophils. Recognized triggers include medications (penicillins, sulfonamides, NSAIDs), systemic infections (streptococcal, hepatitis C), underlying autoimmune disorders (SLE, rheumatoid arthritis), and systemic vasculitides (IgA vasculitis, ANCA-associated vasculitis). A punch biopsy of an early active lesion ($<24\text{--}48\text{ hours}$) is the gold standard diagnostic procedure, complemented by direct immunofluorescence (DIF) to assess vascular IgA or IgG deposition.

\medterm{Leukonychia}-- White discoloration of the nails. True leukonychia: involves the nail matrix and moves distally with nail growth (punctate, striate/Mees lines from arsenic poisoning, or total). Apparent leukonychia: pathology in the vascular nail bed that does not move with nail growth and fades upon digital pressure (e.g., Muehrcke lines from hypoalbuminemia, Terry nails from hepatic cirrhosis/liver failure).

\medterm{Leukoplakia}-- A clinical term for a white patch or plaque on the oral mucosa that cannot be rubbed off and cannot be characterized clinically as any other definable lesion (WHO). Leukoplakia is a potentially malignant disorder of the oral cavity, with a risk of malignant transformation to squamous cell carcinoma of approximately 1-5\% per year. Risk factors: tobacco smoking, alcohol consumption, areca nut/betel quid chewing, and chronic friction (e.g., from sharp teeth or dentures). Clinical subtypes: homogeneous (flat, thin, uniform white, low risk) and non-homogeneous (speckled--red and white, nodular, verrucous, higher risk). Sites: buccal mucosa, alveolar ridge, floor of mouth (highest risk), tongue, palate. The term may also rarely be used for vulvar, penile, or laryngeal white patches. Diagnosis: biopsy (optional for homogeneous lesions, always indicated for non-homogeneous or suspicious lesions) to confirm histology (hyperkeratosis, acanthosis, dysplasia). The presence of epithelial dysplasia (mild, moderate, severe) predicts malignant potential. Management: elimination of risk factors (smoking cessation, alcohol reduction), observation for homogeneous lesions, surgical excision for dysplastic or non-homogeneous lesions, CO2 laser ablation, cryotherapy. Photodynamic therapy and topical retinoids have limited evidence. Long-term follow-up is essential.

\medterm{Lice Infestations}-- Pediculosis is infestation by head, body, or pubic lice. Head lice cause scalp pruritus and nits attached to hair shafts; body lice live mainly in clothing; pubic lice affect coarse hair-bearing regions. Treatment includes an appropriate pediculicide, mechanical nit removal when relevant, laundering of clothing and bedding, and examination of close contacts.

\medterm{Lichen}-- A descriptive dermatological term referring to a flat-topped, firm, papular skin lesion resembling the lichen plant (a symbiotic organism of algae and fungi growing on rocks and trees). Lichenoid describes lesions that resemble lichen planus or lichenoid drug eruptions. Specific lichenoid conditions include: lichen planus, lichen sclerosus, lichen simplex chronicus (localised neurodermatitis from chronic scratching), lichen striatus, lichen nitidus, lichen myxoedematosus (papular mucinosis), and lichenoid drug eruptions (e.g., ACE inhibitors, antimalarials, gold, quinidine, NSAIDs, thiazides). See individual entries.

\medterm{Lichen Planus}-- A T-cell mediated inflammatory condition affecting skin, mucous membranes, nails, and hair. Classic presentation: the 6 Ps (pruritic, purple, polygonal, planar, papules, plaques), often with Wickham striae (white lacy network on surface). Oral lichen planus affects buccal mucosa and may be erosive. Hypertrophic lichen planus affects shins. Nail lichen planus causes longitudinal ridging and pterygium. Treatment includes topical corticosteroids, systemic corticosteroids for severesevere disease, and acitretin. Oral and genital lichen planus may require topical calcineurin inhibitors (tacrolimus, pimecrolimus) as first-line therapy due to the risk of corticosteroid-induced atrophy of mucosal surfaces.

\medterm{Lichen Sclerosus}-- A chronic inflammatory skin disease most commonly affecting anogenital skin. It presents as ivory-white, atrophic plaques with wrinkled paper-like appearance. In women, it predominantly affects the vulva and perianal region (figure-eight distribution). In men, it affects the glans penis and foreskin (balanitis xerotica obliterans). Lichen sclerosus is associated with increased risk of squamous cell carcinoma. Treatment includes ultra-potent topical corticosteroids (clobetasolclobetasol propionate 0.05 percent ointment) as first-line therapy, applied once or twice daily for 12 to 24 weeks, with gradual tapering to maintenance. Topical calcineurin inhibitors (tacrolimus, pimecrolimus) are used for corticosteroid-sparing effect. Systemic therapy with methotrexate, mycophenolate mofetil, or phototherapy is reserved for refractory, extensive, or extragenital disease. Long-term monitoring is required due to the risk of malignant transformation (squamous cell carcinoma develops in 3 to 5 percent of patients with vulval lichen sclerosus).

\medterm{Lichen Simplex Chronicus}-- A thickened, leathery plaque caused by chronic scratching and rubbing (itch-scratch cycle). It commonly affects the neck, ankles, lower legs, and genital area. It is a form of neurodermatitis associated with stress and anxiety. Treatment includes breaking the itch-scratch cycle with topical corticosteroids, antihistamines, and behavioral modification. Underlying atopic dermatitis or contact dermatitis should be treated.

\medterm{Linear IgA bullous dermatosis}-- An autoimmune blistering disease characterised by linear deposition of IgA antibodies at the basement membrane zone, presenting with tense vesicles and bullae arranged in a "string of pearls" configuration on the trunk and extremities. It may be idiopathic or drug-induced, most commonly by vancomycin. Dapsone is the first-line treatment.

\medterm{Lipedema}-- A chronic, progressive, painful disorder of adipose tissue distribution and microvascular dysfunction almost exclusively affecting postpubertal females, characterized by the symmetrical, bilateral, abnormal deposition of subcutaneous fat in the lower extremities, from the hips and buttocks down to the ankles, with striking and abrupt sparing of the dorsal feet (creating a characteristic 'cuff' or 'bracelet' sign at the malleoli). Pathogenesis remains incompletely defined, involving genetic predisposition, estrogenic hormonal influences (often triggered at puberty, pregnancy, or menopause), altered adipocyte hypertrophy, and microvascular hyperpermeability leading to chronic interstitial fluid accumulation and low-grade inflammation. Unlike simple obesity or primary lymphedema, lipedematous adipose tissue is exquisitely tender to palpation, prone to spontaneous ecchymoses and easy bruising following minor trauma, and completely unresponsive to caloric restriction or intense exercise. Stemmer's sign (the inability to pinch a fold of skin at the base of the second toe) is negative in uncomplicated lipedema, distinguishing it from lymphedema, although chronic advanced lipedema can overwhelm lymphatic drainage and progress to secondary lipolymphedema. Management encompasses conservative decongestive therapy (compression garments, manual lymphatic drainage) and specialized tumescent water-assisted liposuction.

\medterm{Lipoma}-- The most common benign soft tissue tumor, composed of mature adipocytes. It presents as a soft, mobile, painless subcutaneous nodule, commonly on the trunk, neck, and proximal extremities. Lipomas grow slowly and may be multiple (familial multiple lipomatosis). Diagnosis is clinical; ultrasound or MRI for deep or atypical lesions. Treatment is observation or surgical excision for symptomatic or cosmetically concerning lipomas. Angiolipoma is a painful variant with vascular proliferatiovascular proliferation that is more symptomatic than typical lipomas.

\medterm{Livedo Reticularis}-- A net-like (reticulate) violaceous discoloration of the skin due to reduced blood flow in the dermal vasculature. Physiologic (cold-induced, benign) vs. pathologic: primary livedo reticularis (idiopathic, benign), and secondary—vasculitis (polyarteritis nodosa—livedo racemosa with nodules), SLE, antiphospholipid syndrome (livedo reticularis + thromboses + fetal loss), Sneddon syndrome, cryoglobulinemia, cholesterol emboli, and infections. Diagnosis: history (cold exposure, systemic symptoms, thromboses, pregnancy loss), exam, ANA/antiphospholipid antibodies, ESR/CRP, and biopsy if vasculitis suspected. Management: treat the underlying cause; avoid cold; anticoagulation for antiphospholipid syndrome; treat vasculitis.

\medterm{Macular}-- A macule: a flat, circumscribed change in skin color <1 cm with no change in texture or elevation (non-palpable). Color varies: brown (melanin—freckles, lentigines, cafe-au-lait, nevus), red (erythema—inflammatory, vascular—telangiectasia, petechiae), white (vitiligo), blue/gray. Macules are distinguished from papules (elevated). Examples of macular lesions: maculopapular rash (measles, drug eruption), pityriasis rosea, tinea versicolor, and early melanoma can present as a pigmented macule (ABCDE). Larger flat lesions >1 cm are patches. Diagnosis by morphology and distribution; biopsy for suspicion of malignancy.

\medterm{Macule}-- A fundamental primary cutaneous lesion defined as a flat, distinct, circumscribed area of cutaneous or mucosal discoloration that is completely flush with the surrounding skin surface, measuring less than $1.0\text{ cm}$ in diameter. Because macules involve no elevation or depression, they cannot be detected by palpation with the eyes closed. Pathophysiologically, macular discoloration can result from: (1) alterations in melanin pigmentation, including hyperpigmentation (ephelides/freckles, junctional melanocytic nevi) or hypopigmentation and depigmentation (vitiligo, tinea versicolor, post-inflammatory hypopigmentation); (2) vascular alterations, such as localized erythema from dermal capillary vasodilation (which blanches completely upon glass diascopy) or non-blanching petechiae resulting from erythrocyte extravasation into the superficial dermis; or (3) deposition of exogenous or endogenous pigments (hemosiderin, silver, tattoo ink). Flat lesions of altered color measuring greater than $1.0\text{ cm}$ in diameter are classified morphologically as patches. Identification of macular borders, color uniformity, and diascopy response forms a cornerstone of initial dermatologic examination.

\medterm{Malignant acanthosis nigricans}-- A paraneoplastic syndrome in which velvety, hyperpigmented, verrucous plaques develop rapidly in flexural and mucosal areas as a marker of an underlying internal malignancy, most commonly gastric adenocarcinoma. Unlike benign acanthosis nigricans, the malignant form is extensive, abrupt in onset, and may involve the oral mucosa. Recognition should prompt investigation for an occult cancer.

\medterm{Malignant melanoma}-- A highly aggressive malignancy arising from melanocytes, characterised by uncontrolled proliferation and potential for early metastasis via lymphatic and haematogenous spread. It typically presents as an enlarging, irregularly pigmented lesion with asymmetry, border irregularity, colour variation, and diameter exceeding 6 mm (ABCDE criteria). Staging by Breslow thickness, ulceration, and sentinel lymph node biopsy guides prognosis and management.

\medterm{Malignant syphilis}-- A severe, destructive form of secondary syphilis presenting with necrotic, gummatous, or ulcerative skin lesions, often in immunocompromised patients, particularly those with HIV co-infection. It represents an aberrant immune response to heavy Treponema pallidum infection and requires aggressive intravenous penicillin therapy. Without treatment it can cause extensive tissue destruction.

\medterm{Medical Tattooing}-- Medical tattooing (permanent makeup, micropigmentation) covers depigmented skin (vitiligo, vitiligo camouflage), restores areolas after mastectomy, and masks scars. Trichopigmentation simulates hair on bald scalp. Medical tattooing uses specialized pigments and techniques. Contraindications include keloid propensity, active skin disease, and pregnancy. Complications include infection, allergic reaction, pigment migration, and dissatisfaction. Removal requires laser therapy.

\medterm{Melanoma}-- A malignant neoplasm arising from melanocytes (pigment-producing cells derived from the neural crest, located in the basal layer of the epidermis, hair follicles, and mucous membranes), representing the most lethal form of skin cancer, with an incidence of approximately 20 to 30 per 100,000 persons per year in fair-skinned populations and a lifetime risk of approximately 2 to 3 percent. Risk factors include ultraviolet radiation exposure (intermittent intense sun exposure with blisteblistering sunburns during childhood or adolescence (the single most important modifiable risk factor), high density of melanocytic nevi (greater than 50 nevi or the presence of atypical or dysplastic nevi, which confer a 5- to 20-fold increased risk), red or blonde hair, fair skin (Fitzpatrick I and II), positive family history of melanoma, personal history of melanoma (5 to 10 percent risk of developing a second primary melanoma), and immunosuppression. The four main histopathological subtypes are superficial spreading melanoma (the most common, accounting for 70 percent of all melanomas), nodular melanoma (15 percent, presenting as a rapidly growing, elevated nodule without a clinically appreciable radial growth phase, often diagnosed late with a worse prognosis), lentigo maligna melanoma (10 percent, arising on sun-damaged facial skin of elderly patients), and acral lentiginous melanoma (5 percent, the most common melanoma in darker skin types, Fitzpatrick IV to VI, arising on the palms, soles, and subungual sites). The diagnosis is confirmed by excisional biopsy with a 1 to 2 mm clinical margin for histopathological evaluation.

\medterm{Melanoma Surveillance}-- regular skin examination, patient education on self-examination, and dermoscopic monitoring of atypical nevi. High-risk patients (personal or family history, numerous atypical nevi, immunosuppression) require 6-month follow-up. Standard patients require annual examination. Total body photography establishes baseline for change detection. Sentinel lymph node biopsy is recommended for melanoma 1 mm or greater or with ulceration.

\medterm{Melanoma Treatment}-- Stage-dependent. Stage 0 to IA: wide local excision. Stage IB to IIA: wide excision with sentinel lymph node biopsy. Stage IIB to IIC: wide excision, SLNB, and adjuvant therapy. Stage III: wide excision, SLNB or lymph node dissection, and adjuvant immunotherapy (pembrolizumab, nivolumab) or targeted therapy (dabrafenib plus trametinib for BRAF-mutant). Stage IV: systemic therapy (immunotherapy, targeted therapy, chemotherapy). BRAF and c-KIT mutations guide targeted therapytargeted therapy). Palliative radiation for brain or bone metastases. Regular surveillance with imaging and skin examinations is required.

\medterm{Melasma (Chloasma)}-- An acquired, chronic pigmentary disorder characterized by symmetric, hyperpigmented, light-to-dark brown macules and patches on sun-exposed areas of the face (malar, forehead, upper lip, chin). It primarily affects women of childbearing age and individuals with darker skin types (Fitzpatrick types III-V). Pathogenesis involves melanocyte hyperactivity stimulated by ultraviolet (UV) radiation, visible light, hormonal fluctuations (pregnancy, oral contraceptives), and genetic predisposition. Management relies on strict photoprotection (broad-spectrum sunscreen with iron oxide to block visible light), topical hypopigmenting agents (hydroquinone, azelaic acid, retinoids, or triple combination cream), chemical peels, and low-fluence laser therapy.

\medterm{Merkel Cell Carcinoma}-- An aggressive neuroendocrine skin cancer arising from Merkel cells. It presents as a rapidly growing, painless, red-violet nodule on sun-exposed skin of elderly patients. Approximately 80 percent are associated with Merkel cell polyomavirus. Risk factors include immunosuppression, UV exposure, and fair skin. Treatment includes wide local excision, sentinel lymph node biopsy, and radiation therapy. Prognosis is poor with high rates of local recurrence and metastasis.

\medterm{Methotrexate in Dermatology}-- A critical disease-modifying antirheumatic and immunomodulatory antimetabolite widely prescribed in clinical dermatology for the management of moderate-to-severe inflammatory, autoimmune, and hyperproliferative cutaneous disorders, notably chronic plaque psoriasis, psoriatic arthritis, severe atopic dermatitis, pityriasis rubra pilaris, dermatomyositis, cutaneous lupus, and bullous pemphigoid. Methotrexate is a structural folic acid analogue that competitively inhibits dihydrofolate reductase (DHFR) and other folate-dependent enzymes, depleting intracellular tetrahydrofolates required for purine and pyrimidine biosynthesis, thereby suppressing rapidly dividing lymphocytes and keratinocytes. At low dermatologic doses (typically $7.5\text{--}25\text{ mg}$ administered once weekly, orally or subcutaneously), its primary anti-inflammatory action is mediated by the intracellular accumulation of aminoimidazolecarboxamide ribonucleotide (AICAR), which elevates extracellular levels of adenosine, a potent endogenous inhibitor of neutrophil function and pro-inflammatory cytokine secretion (TNF-$\alpha$, IL-1). Routine daily folic acid supplementation ($1\text{--}5\text{ mg/day}$, omitted on the day of methotrexate) reduces gastrointestinal and hematologic toxicities. Comprehensive laboratory surveillance is mandatory, monitoring complete blood counts, renal function, and hepatic fibrosis markers (FibroScan or procollagen III peptide). It is strictly teratogenic (pregnancy category X).

\medterm{Milia}-- Small, white, keratin-filled cysts appearing on the face, particularly around the eyes. They are common in newborns (neonatal milia) and adults (primary milia). Secondary milia occur after blistering injuries, dermabrasion, or laser therapy. Milia are asymptomatic and may resolve spontaneously. Treatment includes extraction with a sterile needle, topical retinoids, and chemical peels.

\medterm{Mole}-- Nevus: a benign proliferation of melanocytes, the most common skin lesion; also used for other nevi (moles are melanocytic nevi). Types: junctional (flat, brown—childhood), compound (slightly raised), intradermal (dome-shaped, flesh/pigmented). Congenital nevi (large ones—risk of melanoma), dysplastic (atypical) nevi (risk factor for melanoma). Features of melanoma concern: asymmetry, border irregularity, color variation, diameter >6 mm, evolution (ABCDE). Management: routine examination; any changing, atypical, or suspicious lesion warrants dermatoscopy and biopsy/excision. Sun protection reduces melanoma risk.

\medterm{Molluscum Contagiosum}-- A viral skin infection caused by a poxvirus, presenting as firm, dome-shaped, umbilicated papules with central dimpling, 2 to 5 mm in diameter. It is transmitted by direct skin contact, autoinoculation, and sexual contact in adults. Immunocompetent hosts clear lesions spontaneously within 6 to 12 months. Immunocompromised patients may have hundreds of lesions. Treatment includes cryotherapy, curettage, topical cantharidin, potassium hydroxide, and imiquimod. Lesions resoLesions resolve within weeks to months after treatment. Spontaneous immunity is not lifelong.

\medterm{Morton Neuroma}-- A perineural fibrosis of the plantar digital nerve, most commonly between the third and fourth metatarsal heads. It presents as burning pain, numbness, and a feeling of walking on a pebble. It is not a true neuroma but a mechanical entrapment neuropathy. Treatment includes wide shoes, metatarsal pads, corticosteroid injection, and surgical excision for refractory cases. Alcohol injection sclerotherapy is an alternative.

\medterm{Mpox (monkeypox)}-- A zoonotic viral disease caused by the monkeypox virus, an orthopoxvirus in the same genus as smallpox (variola) and vaccinia. The virus has two clades: clade I (formerly Congo Basin, Central African) with higher virulence and mortality (up to 10 percent), and clade II (formerly West African) with lower mortality (less than 1 percent). The 2022 global outbreak involved clade IIb, which spread predominantly through intimate contact, including sexual contact, among men who have sex with men, with cases concentrated in men who have sex with men and their close contacts. The incubation period is typically 5 to 21 days. The prodrome consists of fever, headache, myalgia, and lymphadenopathy (the most distinguishing feature from smallpox), followed by the eruption of a deep-seated, well-circumscribed, umbilicated vesiculopustular rash that evolves synchronously through macular, papular, vesicular, pustular, and crusting stages over 2 to 4 weeks, with centrifugal distribution. Oral antiviral therapy with tecovirimat (TPOXX) is used for severe disease, immunocompromised patients, and those at high risk of complications.

\medterm{Myiasis}-- Infestation of living tissue by larvae (maggots) of various fly species, which feed on host tissue and can cause pain, secondary infection, and tissue destruction. It may be cutaneous, involving the skin and subcutaneous tissue, or affect other body sites including the nose, ears, and wounds. Treatment involves physical removal of the larvae and wound care.

\medterm{Nail}-- A specialized, convex plate of hard keratin located on the dorsal aspect of the terminal phalanx of each digit, serving to protect the distal extremity, enhance fine motor manipulation and tactile counter-pressure, and function as an accessible clinical window into systemic and cutaneous diseases. The nail unit comprises: (1) the nail plate (compacted, anucleated, hard keratinized cells); (2) the nail matrix (the proliferative germinative epithelium that synthesizes the nail plate, visible proximally as the white, crescent-shaped lunula); (3) the nail bed (vascularized supporting stroma beneath the plate, providing adherence); (4) the eponychium and cuticle (sealing the proximal nail fold against environmental pathogens); and (5) the hyponychium (distal barrier beneath the free edge). Pathological nail alterations are highly informative diagnostic signs: pitting and salmon 'oil-drop' patches indicate psoriasis; longitudinal grooving, thinning, and pterygium indicate lichen planus; transverse depressions (Beau lines) reflect transient systemic matrix arrest; clubbing signifies chronic cardiopulmonary disease or malignancy; koilonychia (spoon nails) reflects iron deficiency anemia; and longitudinal melanonychia requires evaluation to exclude subungual malignant melanoma.

\medterm{Nail Changes}-- Abnormalities of the nails indicating local or systemic disease: koilonychia (spooning—iron deficiency), clubbing (chronic hypoxia—lung cancer, bronchiectasis, cyanotic heart disease, IBD), onycholysis (separation—psoriasis, thyroid, trauma), Beau lines (transverse grooves—acute illness, chemotherapy), Terry nails (white with distal pink band—liver disease), Muehrcke lines (hypoalbuminemia), splinter hemorrhages (endocarditis), pitting (psoriasis), yellow nail syndrome, paronychia (infection), and onychomycosis (fungal). Diagnosis: history, exam, nail culture/biopsy, and systemic evaluation. Management: treat the underlying disease; nail care and antifungals for fungal infection.

\medterm{Nail Clubbing}-- Digital clubbing (Hippocratic fingers); an acquired or congenital structural deformity characterized by bulbous, drumstick-like expansion of the soft tissues of the terminal phalanges of the fingers and toes, accompanied by increased longitudinal and transverse nail curvature. Pathophysiologically, clubbing is driven by the impaction of megakaryocytes and platelet clumps in the distal digital capillary bed, having escaped pulmonary capillary fragmentation due to right-to-left cardiac shunts or pulmonary arteriovenous malformations. These trapped platelets release platelet-derived growth factor (PDGF) and vascular endothelial growth factor (VEGF), inducing local endothelial proliferation, vascular permeability, and stromal connective-tissue hypertrophy. Clinically, clubbing exhibits: (1) loss of the normal Lovibond angle (the angle between the proximal nail fold and the nail plate, normally $<160^\circ$, which increases to $>180^\circ$); and (2) obliteration of the diamond-shaped window formed when opposing terminal phalanges are placed dorsum to dorsum (positive Schamroth sign). Over $80\text{--}90\%$ of clubbing cases are secondary to serious thoracic pathology, including bronchogenic lung carcinoma, idiopathic pulmonary fibrosis, bronchiectasis, cystic fibrosis, cyanotic congenital heart disease, and subacute infective endocarditis; non-thoracic causes include inflammatory bowel disease, cirrhosis, and Graves' acropachy.

\medterm{Nail Pits}-- Discrete, punctate, cupuliform depressions or 'dimples' on the dorsal surface of the nail plate, representing one of the most common and clinically recognizable physical signs of proximal nail matrix pathology. Nail pitting develops when small foci of parakeratotic cells within the dorsal aspect of the proximal nail matrix become incorporated into the newly formed nail plate; as the plate grows distally out from under the proximal nail fold, these abnormal, loosely cohesive parakeratotic cell clusters desquamate prematurely, leaving characteristic depressions on the dorsal nail surface. Nail pits are most famously associated with psoriasis (nail psoriasis), where they appear as deep, irregular, randomly scattered or geometric pits, often accompanied by salmon 'oil-drop' patches, subungual hyperkeratosis, and onycholysis. Pitting is also observed in alopecia areata (where pits are classically smaller, shallower, highly geometric, and arranged in regular horizontal or grid-like rows across multiple nails), atopic dermatitis, lichen planus, and secondary syphilis. Clinical evaluation involves counting and characterizing the pit morphology across all twenty nail units to assist in differential diagnosis and monitoring response to systemic or topical therapies.

\medterm{Nail Psoriasis}-- Nail psoriasis affects 50 percent of patients with psoriasis and 80 percent with nail involvement having joint disease. Features include pitting (most common), onycholysis (nail separation), oil drop sign (salmon patch), subungual hyperkeratosis, and nail plate crumbling. Nail psoriasis is associated with psoriatic arthritis. Treatment is challenging: topical corticosteroids, intralesional corticosteroids, systemic therapy for psoriasis, and apremilast. Nail avulsion is rarely needed.

\medterm{Narrowband UVB Phototherapy}-- Targeted phototherapeutic irradiation utilizing a narrow, specific spectrum of ultraviolet B electromagnetic radiation centered precisely at $311\text{--}313\text{ nm}$ (emitted by specialized Philips TL-01 fluorescent lamps), established as a first-line, gold-standard physical therapy in clinical dermatology for widespread inflammatory and autoimmune dermatoses. Narrowband UVB (NB-UVB) significantly surpasses older broadband UVB ($290\text{--}320\text{ nm}$) by isolating the most therapeutically potent, immunosuppressive, and anti-inflammatory wavelengths while eliminating the shorter, burning, erythemogenic wavelengths below $300\text{ nm}$. Mechanistically, NB-UVB penetrates the epidermis and upper papillary dermis, inducing pyrimidine dimer formation in DNA, activating apoptosis in pathologically activated cutaneous T lymphocytes and hyperproliferative keratinocytes, depleting epidermal Langerhans cells, and downregulating pro-inflammatory cytokines (IL-17, IL-22, IFN-$\gamma$). Primary clinical indications include moderate-to-severe plaque psoriasis, atopic dermatitis, vitiligo (stimulating follicular melanocyte migration), generalized pruritus, cutaneous T-cell lymphoma (early-stage mycosis fungoides), and pityriasis lichenoides. Administered typically two to three times weekly in specialized irradiation cabins with incremental dosimetry guided by minimal erythema dose (MED) or standardized skin phototype protocols, accompanied by mandatory UV-blocking eye goggles and genital shielding.

\medterm{Necrobiosis Lipoidica}-- A chronic granulomatous skin disorder strongly associated with Diabetes Mellitus (present in 0.3-1.2\% of diabetic patients, though 20\% of affected patients do not have diabetes). Characterized by well-demarcated, shiny, yellow-brown shiny plaques with a central atrophic telangiectatic center and prominent violaceous border, most commonly on the pretibial area of lower extremities. Ulceration occurs in 35\% of cases. Histology: palisading granulomas surrounding degenerated collagen (necrobiosis) and blood vessel wall thickening. Management: topical or intralesional corticosteroids (first-line), calcineurin inhibitors, pentoxifylline, and glucose control.

\medterm{Necrolytic migratory erythema}-- A recurrent, erosive, erythematous eruption with a migrating, annular or serpiginous pattern, characteristically involving the perineum, groin, and extremities, and almost always associated with an underlying glucagon-secreting pancreatic alpha-cell tumour (glucagonoma). It results from nutritional deficiency secondary to the catabolic effects of excess glucagon. Diagnosis requires demonstration of elevated serum glucagon and imaging of the pancreas.

\medterm{Neonatal Skin Conditions}-- Common and usually benign. Erythema toxicum neonatorum (red papules and pustules on erythematous base, appears day 2 to 3, resolves by 2 weeks). Milia (tiny white cysts on nose and cheeks). Miliaria crystallina (clear vesicles from sweat duct obstruction). Sebaceous hyperplasia (yellow papules on nose and cheeks). Transient neonatal pustular melanosis (vesiculopustular rash with collarette of scale, resolves spontaneously). Acne neonatorum (comedones and papules, matmatures with sebaceous gland development and resolves without scarring). Treatment is reassurance and education for parents. Topical treatments are rarely needed.

\medterm{Nettle Rash}-- A colloquial and historical clinical term for acute urticaria (hives), derived from its classic morphologic resemblance to the transient erythematous, intensely pruritic wheals produced by cutaneous contact with the stinging nettle plant (\textit{Urtica dioica}), which inoculates histamine, acetylcholine, and serotonin into the superficial dermis. In modern clinical dermatology, nettle rash represents an acute mast-cell-mediated condition characterized by transient edema of the superficial dermis. Activation of dermal mast cells and basophils triggers rapid degranulation and the release of preformed histamine, leukotrienes, and prostaglandins, which bind to endothelial receptors to produce microvascular vasodilation, increased capillary permeability, and sensory nerve stimulation. Lesions manifest as edematous, circumscribed, pink-to-erythematous wheals with surrounding flares and central blanching, characteristically migratory and resolving without residual hyperpigmentation within 24 hours. Etiologies include viral illnesses, IgE-mediated food or medication allergies, insect stings, and physical stimuli (cold, pressure, heat, dermatographism). Management focuses on trigger avoidance and first-line treatment with high-dose second-generation non-sedating $H_1$-antihistamines (cetirizine, fexofenadine, bilastine).

\medterm{Neurofibroma}-- A benign peripheral nerve sheath tumor. Solitary neurofibromas are common and present as soft, flesh-colored papules or nodules. Multiple neurofibromas are characteristic of neurofibromatosis type 1 (NF1). Cutaneous neurofibromas are diagnostically significant in NF1 (greater than or equal to 2 neurofibromas is one of the diagnostic criteria). Plexiform neurofibromas are pathognomonic for NF1 and carry risk of malignant transformation to malignant peripheral nerve sheath tumor (MMPNST), occurring in 5 to 10 percent of patients with NF1, presenting as rapid growth, new or increasing pain, or a change in consistency of a previously stable plexiform neurofibroma.

\medterm{Neutrophilic dermatosis}-- A group of inflammatory skin conditions characterised by dense neutrophilic infiltrate in the dermis without evidence of infection, which includes Sweet syndrome (acute febrile neutrophilic dermatosis), pyoderma gangrenosum, and Behçet disease. These conditions may be idiopathic, drug-induced, or associated with underlying haematological malignancy or inflammatory bowel disease. Treatment typically involves systemic corticosteroids.

\medterm{Nevus}-- A benign hamartoma of the skin; most commonly melanocytic nevus (mole). Other nevi: congenital, dysplastic/atypical, Spitz, blue nevus, halo nevus; non-melanocytic: epidermal, sebaceous, vascular (strawberry hemangioma, port-wine stain). Melanocytic nevi arise from melanocyte proliferation; atypical (dysplastic) nevi are irregular and are markers of increased melanoma risk, especially with family history. Congenital melanocytic nevi have small malignant potential (larger = higher). Diagnosis: clinical/dermoscopy; biopsy for atypical or changing lesions. Management: monitor; excise suspicious lesions (ABCDE criteria).

\medterm{Nevus comedonicus}-- A congenital hamartoma of the pilosebaceous unit presenting as a localised cluster of dilated follicular ostia filled with keratinous plugs, resembling closely grouped open comedones. It typically appears at birth or in early childhood on the face, neck, or trunk, and may be associated with other developmental anomalies. Treatment is primarily for cosmetic concerns and includes topical retinoids and laser therapy.

\medterm{Nikolsky Sign}-- A clinical dermatological sign in which light lateral pressure applied to normal-appearing skin causes epidermal detachment or sloughing. Positive sign indicates intraepidermal loss of cell-cell cohesion (acantholysis). Pathognomonic for Pemphigus Vulgaris (IgG anti-desmoglein 1/3 autoantibodies) and Staphylococcal Scalded Skin Syndrome (SSSS, exfoliative toxins A/B). Also positive in Toxic Epidermal Necrolysis (TEN) and Stevens-Johnson Syndrome (SJS). Negative in Bullous Pemphigoid (subepidermal split, hemidesmosomes intact).

\medterm{Nonmelanoma Skin Cancer}-- A group of skin cancers that arise from the cells of the epidermis, excluding melanoma. The two most common types are basal cell carcinoma (BCC) and squamous cell carcinoma (SCC), which together account for the vast majority of all skin cancers diagnosed globally. The primary risk factor is cumulative exposure to ultraviolet (UV) radiation from sunlight or tanning beds. Other risk factors include fair skin type, chronic immunosuppression, exposure to ionizing radiation or arsenic, and genetic syndromes (e.g., xeroderma pigmentosum). Unlike melanoma, nonmelanoma skin cancers generally grow slowly and have a low risk of metastasis, though SCC carries a higher metastatic risk than BCC. Treatment modalities include wide local excision, Mohs micrographic surgery, electrodessication and curettage, cryotherapy, photodynamic therapy, and topical agents (e.g., 5-fluorouracil, imiquimod) for superficial lesions.

\medterm{Nutritional Dermatoses}-- Nutritional deficiencies cause distinctive skin findings. Vitamin C deficiency (scurvy): perifollicular hemorrhage, corkscrew hairs, gingival bleeding. Niacin deficiency (pellagra): photosensitive dermatitis in sun-exposed areas, diarrhea, dementia. Riboflavin deficiency (ariboflavinosis): cheilosis, angular stomatitis, glossitis, seborrheic dermatitis-like eruption. Zinc deficiency (acrodermatitis enteropathica): periorificial and acral dermatitis, alopecia, diarrhea. Vitamin A deficiency: follfollicular hyperkeratosis, night blindness, and xerophthalmia (dryness of the conjunctiva and cornea, Bitot spots, keratomalacia). Vitamin E deficiency: acanthocytosis, peripheral neuropathy. Vitamin K deficiency: ecchymoses, hemorrhage. Iron deficiency: kollonychia (spoon nails), angular stomatitis, glossitis, pallor. Zinc deficiency causes acral dermatitis, periorificial dermatitis, alopecia, diarrhea, and growth retardation in children (acrodermatitis enteropathica). Early recognition and appropriate replacement therapy are essential.

\medterm{Occupational Dermatoses}-- Work-related skin diseases. Occupational contact dermatitis accounts for 90 percent of occupational skin diseases. Irritant contact dermatitis is more common (hands, wet work) than allergic contact dermatitis. Occupations at risk include healthcare workers, hairdressers, construction workers, food handlers, and metalworkers. Common occupational allergens include rubber accelerators, fragrances, preservatives, and metals. Prevention includes hazard identification, protprotective equipment, barrier creams, and education. Management includes identification and avoidance of causative agents, topical corticosteroids, and emollients. Workers with occupational dermatitis may require job modification or retraining.

\medterm{Ocular myiasis}-- Infestation of the eye or ocular adnexa by fly larvae, which may involve the conjunctiva, cornea, or intraocular structures, presenting with pain, irritation, tearing, and visual disturbance. In severe cases, larvae can penetrate the globe causing endophthalmitis and blindness. Urgent removal of the larva and ophthalmological management are required.

\medterm{Onycholysis}-- A distinct physical sign and disorder of the nail apparatus defined as the painless, spontaneous separation of the nail plate from the underlying vascularized nail bed, typically commencing at the distal free edge or lateral nail margins and progressing proximally. As the nail plate detaches from the hyponychium and bed, the newly created subungual dead space fills with air, imparting a characteristic opaque, whitish discoloration to the separated portion, which may turn yellow, brown, or green if colonized by opportunistic secondary pathogens such as \textit{Pseudomonas aeruginosa} (pyocyanin pigment) or \textit{Candida} species. Pathophysiological etiologies are extensive: (1) localized mechanical and chemical trauma (aggressive manicuring, subungual debris cleaning, exposure to harsh solvents); (2) primary cutaneous diseases, most notably nail psoriasis (in which subungual parakeratosis induces detachment, often preceded by salmon 'oil-drop' spots), lichen planus, and atopic eczema; (3) infectious diseases, particularly dermatophyte onychomycosis; (4) systemic diseases, famously thyroid dysfunction (Plummer's nails in thyrotoxicosis and hyperthyroidism); and (5) drug-induced phototoxicity (photo-onycholysis induced by tetracyclines/doxycycline, psoralens, or fluoroquinolones upon sun exposure). Management requires identifying the root cause, clipping back the detached plate, keeping the nail bed clean and dry, avoiding manicuring trauma, and targeted antimicrobials or topical therapies.

\medterm{Onychomycosis}-- Fungal nail infection, most commonly caused by dermatophytes (Trichophyton rubrum), affecting toenails more than fingernails. It presents as nail thickening, discoloration (yellow-brown), subungual debris, and nail separation. Diagnosis is confirmed by KOH preparation or fungal culture. Treatment includes oral terbinafine (first-line for dermatophyte), itraconazole, or topical ciclopirox for mild disease. Cure rates are 50 to 80 percent with oral therapy.

\medterm{Ophthalmomyiasis}-- Infestation of the eye by fly larvae (most commonly Oestridae), which may be superficial (conjunctival) or deep (intraocular), with the latter carrying a risk of retinal detachment, endophthalmitis, and permanent vision loss. Superficial ophthalmomyiasis presents with foreign body sensation, irritation, and visible larvae on the conjunctiva. Prompt removal and ophthalmological assessment are essential.

\medterm{Ostraceous psoriasis}-- A severe, chronic variant of plaque psoriasis in which the lesions become thick, hard, and shell-like (ostraceous), with laminated, greyish-brown crusts resembling oyster shells. It is often refractory to standard therapies and may require systemic agents, biologics, or combination treatment. The condition causes significant morbidity due to pain, restriction of movement, and psychosocial impact.

\medterm{Pallor}-- Paleness of the skin and mucous membranes, caused by reduced cutaneous blood flow (vasoconstriction—cold, shock, anemia, fear) or reduced hemoglobin/pigmentation. Causes: anemia (all types—iron deficiency, B12/folate, hemolysis, acute blood loss), vascular (vasovagal, cardiogenic shock, hypovolemia), metabolic (hypothyroidism), and pigment (vitiligo, albinism). Distribution: generalized (anemia) vs. localized. Clinical assessment: mucosal pallor (conjunctivae, palms—pallor in skin creases), with anemia confirmed by CBC/hemoglobin. Management: treat the underlying cause—iron/blood transfusion for significant anemia, volume resuscitation in shock.

\medterm{Palmar erythema of pregnancy}-- Symmetrical redness of the palms occurring during pregnancy, thought to be related to elevated circulating oestrogen levels causing vasodilatation. It is a benign, physiological finding that typically resolves after delivery and requires no specific treatment. It should be distinguished from liver disease-related palmar erythema if other signs of hepatic dysfunction are present.

\medterm{Panniculitis}-- Inflammation of the subcutaneous fat, presenting as tender, erythematous or violaceous nodules and plaques, most commonly on the lower extremities. The classification is based on whether the septa or lobules of the fat are primarily affected, and the aetiology includes infections, autoimmune diseases, medications, and pancreatic disease. Diagnosis often requires deep incisional biopsy for histopathological examination.

\medterm{Papilloma}-- A benign epithelial neoplasm characterised by finger-like projections of fibrovascular tissue covered by hyperplastic epithelium, most commonly caused by human papillomavirus infection in cutaneous warts or mucosal sites. Cutaneous papillomas present as flesh-coloured, rough, or pedunculated growths. Treatment options include cryotherapy, salicylic acid, and surgical excision.

\medterm{Papillomatosis cutis lymphostatica}-- Diffuse, warty, papillomatous thickening of the skin resulting from chronic lymphoedema, most commonly affecting the lower extremities in the setting of long-standing venous or lymphatic insufficiency. The skin develops cobblestone-like hyperkeratotic papules and plaques due to lymph stasis and chronic inflammation. Treatment is directed at the underlying lymphoedema with compression and skin care.

\medterm{Papular}-- A papule: a small, solid, elevated skin lesion <1 cm in diameter (palpable), of various colors and etiologies. Types: inflammatory (papular urticaria, lichen planus—purple polygonal, psoriasis—scaly), infectious (verruca/wart, molluscum, folliculitis), neoplastic (basal cell carcinoma, nevus, keratoacanthoma), and deposition (xanthoma, amyloid). Papulopustular: acne. Description includes color, shape (dome, flat-topped), surface, and distribution. Biopsy for uncertain diagnosis. Contrast with macule (flat) and nodule (>1 cm).

\medterm{Papular Urticaria}-- A hypersensitivity reaction to insect bites, especially from fleas, bedbugs, or mosquitoes, causing recurrent intensely itchy papules or small vesicles, often clustered on exposed areas in children. Diagnosis is clinical. Management includes insect control, protective clothing, topical corticosteroids, and oral antihistamines; secondary bacterial infection requires treatment.

\medterm{Papule}-- A primary cutaneous lesion defined as a solid, elevated, circumscribed, palpable superficial skin lesion measuring less than $1.0\text{ cm}$ (or traditionally $<0.5\text{ cm}$) in diameter. Papules project distinctly above the surrounding plane of the skin and are palpable to the examining finger. Morphologically and pathophysiologically, papules can be formed by: (1) localized epidermal thickening and hyperkeratosis (verruca vulgaris, seborrheic keratosis); (2) cellular infiltration within the dermis, including inflammatory infiltrates (lymphocytes, histiocytes, neutrophils in lichen planus, granuloma annulare, or secondary syphilis), neoplastic infiltrates (melanocytes in intradermal nevi, basaloid cells in basal cell carcinoma); (3) localized dermal edema (urticarial wheals in their early small phase); or (4) accumulation of metabolic deposits (calcium, amyloid, mucin). Papules exhibit diverse surface characteristics and cross-sectional geometries, described clinically as umbilicated (molluscum contagiosum), flat-topped and polygonal (lichen planus), dome-shaped, filiform, verrucous, or follicular (keratosis pilaris, acne). The coalescence of adjacent papules forms a plaque.

\medterm{Paronychia}-- Inflammation of the nail fold. Acute paronychia is bacterial (Staphylococcus aureus, Streptococcus), presenting as painful, erythematous, swollen nail fold with possible abscess. Treatment includes warm compresses, antibiotics, and incision and drainage if abscessed. Chronic paronychia is fungal (Candida) or irritant-induced, lasting greater than 6 weeks, common in wet workers. Treatment includes topical antifungals, corticosteroids, and avoiding moisture.

\medterm{Patch}-- A primary cutaneous lesion defined as a flat, non-palpable, circumscribed area of cutaneous or mucosal discoloration that is completely level with the adjacent skin surface and measures greater than $1.0\text{ cm}$ in diameter. It represents the larger dimensional counterpart of the macule ($<1.0\text{ cm}$). Because it involves no localized elevation or depression, it cannot be detected by palpation with the eyes closed. Patches can display diverse pigmentary, vascular, or inflammatory origins: (1) melanocytic hyperpigmentation (such as large congenital café-au-lait macules/patches, congenital dermal melanocytosis, melasma, or the horizontal spreading phase of lentigo maligna); (2) melanocytic hypopigmentation or depigmentation (extensive vitiligo, post-inflammatory hypopigmentation); (3) dermal vascular ectasia or capillary malformations (port-wine stains/nevus flammeus, which blanch under glass diascopy); and (4) early macular or superficial inflammatory infiltrates, such as the patch-stage of cutaneous T-cell lymphoma (mycosis fungoides), secondary syphilis, or extensive tinea versicolor. Clinical evaluation focuses on borders, color uniformity, scale, and Wood's lamp enhancement.

\medterm{Patch Testing}-- Patch testing identifies allergens causing allergic contact dermatitis. The standard series includes 35 allergens. Additional custom trays include occupational and personal allergens. Patches are applied to the upper back for 48 hours and removed. Reading at 48 hours (initial) and 72 to 96 hours (delayed) identifies positive reactions. Positive reactions are graded (1+ to 4+) based on erythema, edema, and vesiculation. Results guide allergen avoidance and patient education.

\medterm{Pediatric Dermatology}-- Pediatric dermatology encompasses conditions unique to children. Congenital lesions (birthmarks, hemangiomas, port-wine stains, café-au-lait macules). Infantile hemangiomas (strawberry marks) proliferate in first months of life and involute over years. Propranolol is first-line treatment for problematic hemangiomas. Mongolian spots (congenital dermal melanocytosis) are benign and resolve spontaneously. Neonatal acne, erythema toxicum, and milia are self-limited.

\medterm{Pediatric Dermatology Conditions}-- Distinct from adult dermatology. Birthmarks include vascular (hemangiomas, port-wine stains) and pigmented (café-au-lait, congenital nevi, Mongolian spots). Infantile hemangiomas proliferate in first months then involute. Congenital melanocytic nevi carry melanoma risk proportional to size. Juvenile xanthogranuloma is a benign histiocytosis. Neonatal acne resolves spontaneously. Diaper dermatitis is common and treated with barrier creams.

\medterm{Pediatric Eczema}-- Atopic dermatitis in children commonly begins in infancy with facial and extensor involvement; flexural disease predominates in older children. Pruritus, xerosis, sleep disturbance, and scratching are characteristic. Management includes regular emollients, appropriate topical anti-inflammatory therapy, trigger reduction, skin-barrier care, and specialist-directed systemic treatment for severe disease.

\medterm{Pediatric Skin Tumors}-- Pediatric skin tumors differ from adult tumors. Infantile hemangioma is the most common benign tumor of infancy. Congenital melanocytic nevus carries risk of melanoma (larger nevi have higher risk). Juvenile xanthogranuloma is a benign histiocytic lesion. Pyogenic granuloma is a common vascular lesion. Dermoid cysts are congenital. Malignant melanoma is rare but increasing in children. Regular monitoring of congenital nevi is recommended.

\medterm{Pediculosis}-- Infestation with lice. Head lice spread mainly through close head-to-head contact and cause scalp pruritus with nits attached to hair shafts. Body lice live in clothing and can transmit certain infections; pubic lice infest coarse hair-bearing regions. Treatment and contact management depend on the species and local resistance patterns.

\medterm{Pemphigoid}-- Bullous pemphigoid is an autoimmune subepidermal blistering disease caused by antibodies against hemidesmosomal proteins, especially BP180 and BP230. It usually causes intensely pruritic tense blisters on urticarial or erythematous skin, with less mucosal involvement than pemphigus. Diagnosis uses lesional histology, perilesional direct immunofluorescence, and serology; treatment is specialist-directed and may include topical or systemic anti-inflammatory and immunomodulatory therapy.

\medterm{Pemphigoid gestationis}-- An autoimmune blistering disease of pregnancy characterised by intensely pruritic urticarial plaques and tense vesicles or bullae, typically arising in the second or third trimester around the umbilicus and spreading to the trunk and extremities. It is caused by IgG autoantibodies directed against the BP180 antigen at the basement membrane zone. Most cases resolve postpartum but may flare immediately after delivery.

\medterm{Pemphigus}-- A group of autoimmune blistering diseases with intraepidermal blisters due to antibodies against desmogleins (cell adhesion molecules): pemphigus vulgaris (PV—flaccid blisters, mucosal involvement, positive Nikolsky sign), pemphigus foliaceus (superficial, no mucosal), paraneoplastic pemphigus. Causes: autoantibodies (IgG) against desmoglein 1/3, drug-induced (penicillamine, ACE inhibitors). Features: flaccid blisters/erosions, painful oral ulcers; severe: fluid/electrolyte loss, secondary infection. Diagnosis: skin biopsy (intraepidermal acantholysis), direct immunofluorescence (intercellular IgG/C3). Management: high-dose corticosteroids, rituximab (first-line adjunct), immunosuppressants (azathioprine, mycophenolate), IVIG; monitor for infection. Untreated: high mortality.

\medterm{Pemphigus Vulgaris}-- An autoimmune intraepidermal blistering disease caused by IgG antibodies against desmoglein 1 and 3. It presents with flaccid blisters that rupture easily, leaving painful erosions of skin and mucous membranes. Nikolsky sign is positive (lateral pressure causes epidermal sloughing). Diagnosis is by biopsy showing acantholysis and direct immunofluorescence showing intercellular IgG deposition. Treatment includes systemic corticosteroids and rituximab (first-line steroid-sparsteroid-sparing agent). Prognosis is poor without treatment; immunosuppression has dramatically improved outcomes, with mortality reduced to less than 10 percent.

\medterm{Perioral Dermatitis}-- A papulopustular eruption around the mouth, sparing the vermilion border. It is more common in young women and associated with topical corticosteroid use. Other triggers include fluoridated toothpaste, cosmetics, and inhaler steroids. Treatment includes discontinuation of topical steroids, topical metronidazole, erythromycin, or oral tetracyclines. It may take weeks to months to resolve.

\medterm{Photosensitivity}-- An abnormal skin reaction to sunlight (UVA/UVB), manifesting as rash, erythema, or burning on sun-exposed areas. Causes: drug-induced (tetracyclines, NSAIDs, thiazides, sulfonamides, fluoroquinolones, amiodarone, psoralens, retinoids), photodermatoses (polymorphous light eruption, solar urticaria, photoallergic/photoirritant contact dermatitis), porphyria (PCT—fragile blistering skin), lupus (photosensitivity in >50 percent), and xeroderma pigmentosum. Features: exposed-area erythema, papulovesicular rash, burning/tingling. Diagnosis: history (drugs, timing, distribution), phototesting, ANA, porphyrins. Management: strict sun avoidance/photoprotection, stop culprit drugs, treat underlying condition.

\medterm{Phototherapy}-- The systematic, supervised therapeutic administration of artificial, non-ionizing electromagnetic radiation from specific portions of the ultraviolet (UV) or visible light spectrum to treat inflammatory, autoimmune, proliferative, and malignant cutaneous diseases. Phototherapy exerts profound local and systemic immunosuppressive, antiproliferative, and anti-inflammatory biological actions by penetrating the epidermis and dermis, causing selective photochemical reactions in cellular chromophores. Mechanisms include induction of keratinocyte and T-cell apoptosis, inhibition of epidermal antigen-presenting Langerhans cells, shift of Th1/Th17 immune responses toward Th2/Treg profiles, and suppression of pro-inflammatory cytokines. Clinical phototherapy modalities include: (1) Narrowband UVB ($311\text{--}313\text{ nm}$), the first-line treatment for widespread plaque psoriasis, atopic dermatitis, and vitiligo; (2) Broadband UVB ($290\text{--}320\text{ nm}$); (3) Photochemotherapy (PUVA), which pairs oral or topical photosensitizing psoralen (8-methoxypsoralen) with ultraviolet A ($320\text{--}400\text{ nm}$) radiation for recalcitrant dermatoses and cutaneous T-cell lymphoma; and (4) Targeted phototherapy (308 nm excimer laser). Treatment protocols require precise dosimetry, cumulative exposure tracking, eye protection, and surveillance for long-term photoaging and non-melanoma skin cancer risks.

\medterm{Phototoxic dermatitis}-- An acute, dose-dependent inflammatory skin reaction occurring when certain ingested or topically applied chemicals (such as psoralens, tetracyclines, NSAIDs, or thiazides) absorb ultraviolet radiation, leading to direct cellular damage and a sunburn-like eruption in exposed areas. It presents with erythema, edema, and sometimes blistering in a photodistribution pattern. Avoidance of the offending agent and sun protection are the mainstays of management.

\medterm{Pilar Cyst}-- A benign cyst arising from the outer root sheath of hair follicles, most commonly on the scalp. It presents as a firm, mobile, subcutaneous nodule without a central punctum. Unlike epidermoid cysts, pilar cysts have smooth walls and contain dense keratin. They may be multiple in familial cases. Treatment is simple excision. Malignant transformation is extremely rare.

\medterm{Pityriasis Rosea}-- An acute, self-limited inflammatory papulosquamous eruption predominantly affecting adolescents and young adults (aged 10 to 35 years). The etiology is strongly linked to the systemic reactivation of latent human herpesviruses 6 and 7 (HHV-6 and HHV-7), although the exact mechanisms of immunological triggering remain active areas of study. In approximately $80\%$ of cases, the condition begins with a solitary 'herald patch'---a $2\text{--}5\text{ cm}$ round or oval, salmon-pink (in fair skin) or hyperpigmented (in dark skin) plaque with a distinctive inward-facing collarette of fine scale located along its inner border, typically on the trunk or neck. Within one to two weeks, a secondary generalized eruption emerges, characterized by crops of smaller, oval macules and papules that orient along the Langer lines of cutaneous cleavage across the back and chest in a characteristic 'Christmas tree' or fir-tree pattern. Lesions are mildly to moderately pruritic and spontaneously resolve over 6 to 8 weeks without scarring. Because its presentation can closely mimic secondary syphilis, gut-stage psoriasis, or pityriasis versicolor, serologic testing for syphilis (RPR/VDRL) is indicated when atypical or palmar lesions occur. Management is supportive, utilizing oral antihistamines and low-to-medium-potency topical corticosteroids for pruritus.

\medterm{Plantar Wart (Verruca Plantaris)}-- A wart (verruca) located on the sole of the foot, caused by human papillomavirus (HPV, typically types 1, 2, 4, 63). Plantar warts are endophytic (grow inward due to pressure) rather than exophytic like common warts, and are often painful with direct pressure (walking, standing). They appear as well-circumscribed, hyperkeratotic papules or plaques on the sole, with a rough surface and loss of normal skin lines over the lesion. Paring reveals punctate black dots (thrombosed capillaries, distinguishing them from corns). Multiple warts may coalesce into a mosaic wart (mosaic verruca). Diagnosis: clinical, confirmed by paring. Plantar warts are frequently confused with corns (heloma) but are distinguished by the presence of thrombosed capillaries (black dots) and disruption of dermatoglyphics. Treatment: spontaneous resolution occurs in 30-50\% within 1-2 years. First-line: salicylic acid (40\% topical, daily with paring) and cryotherapy (liquid nitrogen, 10-15 second freeze, repeated every 2-4 weeks). Second-line: laser therapy (pulsed dye, CO2), intralesional immunotherapy (Candida or mumps antigen injection), and topical immunotherapy (imiquimod, 5-fluorouracil). Refractory or mosaic warts may require surgical excision or photodynamic therapy. Immunocompromised patients (transplant, HIV) may have more extensive and resistant disease.

\medterm{Plantar warts}-- Common warts (verrucae) occurring on the plantar surface of the foot, caused by human papillomavirus infection of the keratinocytes. They present as hard, thickened papules with a rough surface and disrupted dermatoglyphics (skin lines), often with pinpoint black dots representing thrombosed capillaries. They may be painful with lateral compression and can coalesce into mosaic warts.

\medterm{Plaque}-- A fundamental primary cutaneous lesion defined as a solid, elevated, palpable, plateau-like lesion that measures greater than $1.0\text{ cm}$ in diameter and possesses a surface area that significantly exceeds its height above the skin plane. Plaques frequently arise from the peripheral extension and confluence of multiple adjacent smaller papules, but may also arise \textit{de novo}. The elevated plateau may possess sharp, well-demarcated borders or gradually sloping margins, and its surface may be smooth, verrucous, crusted, fissured, or covered with scale. Pathophysiologically, plaques result from extensive epidermal hyperplasia, stratum corneum hyperkeratosis, and dense inflammatory infiltration of the upper reticular and papillary dermis. The prototypical example is chronic plaque psoriasis (psoriasis vulgaris), which presents as sharply demarcated, erythematous plaques covered by thick, silvery, micaceous lamellar scales that bleed upon mechanical removal (Auspitz sign). Other characteristic plaque-forming dermatoses include lichen simplex chronicus (lichenified plaques from chronic rubbing), mycosis fungoides (infiltrated plaques of cutaneous T-cell lymphoma), nummular eczema, lupus erythematosus, and granuloma annulare.

\medterm{Porphyria Cutanea Tarda}-- The most common porphyria, caused by a deficiency of hepatic uroporphyrinogen decarboxylase (UROD), leading to accumulation of uroporphyrinogens. Triggered by hepatitis C virus (HCV), alcohol abuse, iron overload (hemochromatosis), and estrogen therapy. Clinical features: severe cutaneous photosensitivity, fragile skin, tense bullae, erosions, hyperpigmentation, and hypertrichosis on sun-exposed areas (dorsum of hands, face). Diagnosis: elevated urinary and plasma porphyrins (urine fluoresces coral-red under Wood lamp). Treatment: phlebotomy (to reduce iron stores), low-dose hydroxychloroquine, and treatment of underlying HCV.

\medterm{Porphyrias}-- Metabolic disorders of heme synthesis causing cutaneous and neurological symptoms. Cutaneous porphyrias include porphyria cutanea tarda (blisters on sun-exposed skin, hypertrichosis, sclerodermoid changes), erythropoietic protoporphyria (immediate pain and burning on sun exposure without visible changes), and variegate porphyria. PCT is associated with hepatitis C, HIV, iron overload, and estrogen use. Treatment of PCT includes phlebotomy and hydroxychloroquine. Acute porphyrias (porphyrias (acute intermittent porphyria, variegate porphyria, hereditary coproporphyria) present with acute attacks of abdominal pain, neuropsychiatric symptoms, and autonomic instability. Diagnosis is made by porphyrin analysis in urine, plasma, and faeces. Treatment of acute attacks includes intravenous haem arginate, glucose loading, and avoidance of precipitating factors (drugs, alcohol, fasting, hormones). Long-term management includes sun avoidance, protective clothing, and regular monitoring for complications such as liver disease and hepatocellular carcinoma in PCT.

\medterm{Poststreptococcal pustulosis}-- A sterile, non-follicular pustular eruption occurring as a rare sequela of streptococcal pharyngitis, presenting with scattered pustules on an erythematous base, typically on the trunk and extremities. It is thought to be an immune-mediated hypersensitivity reaction to streptococcal antigens. The eruption is self-limiting and resolves with supportive care.

\medterm{Pressure Injury}-- Localized damage to the skin and/or underlying soft tissue caused by prolonged pressure, shear, or friction, most commonly over bony prominences. Stages: Stage 1 (non-blanchable erythema, intact skin), Stage 2 (partial-thickness skin loss with exposed dermis, shallow open ulcer), Stage 3 (full-thickness skin loss, subcutaneous fat visible), Stage 4 (full-thickness tissue loss with exposed bone, muscle, or tendon), Unstageable (full-thickness tissue loss with eschar or slough obscuring the wound base), and Deep Tissue Injury (persistent non-blanchable deep red/maroon discoloration). Common sites: sacrum, heels, ischial tuberosities, trochanters, and occiput. Risk factors: immobility, malnutrition, incontinence, sensory loss, and decreased perfusion. Prevention: repositioning q2h, pressure-relieving surfaces (air mattress), moisture management, and nutritional support. Treatment: wound debridement, infection control, moist wound healing, NPWT, and surgical closure (flap) for advanced stages.

\medterm{Pressure Ulcer}-- Localized damage to skin and underlying tissue caused by sustained pressure or pressure with shear, typically over a bony prominence. Staging includes non-blanchable erythema, partial-thickness loss, full-thickness loss, and full-thickness tissue loss with exposed or palpable deep structures; obscured wounds are unstageable. Prevention includes pressure redistribution, repositioning, moisture management, nutrition support, and treatment of contributing illness.

\medterm{Pressure Urticaria}-- A form of chronic inducible urticaria in which sustained pressure causes localized swelling, erythema, and pain, usually after a delay of several hours. Lesions may last 8 to 72 hours and commonly affect the hands, feet, buttocks, or areas compressed by clothing or sitting. Treatment includes trigger avoidance, antihistamines, and specialist-directed therapy for persistent disease.

\medterm{Pruritus}-- Itching: an unpleasant sensation provoking scratching, a common symptom of skin and systemic disease. Causes: dermatologic (eczema, psoriasis, urticaria, scabies, xerosis, dermatitis), systemic (cholestasis—biliary obstruction, renal failure/uremia, hematologic—polycythemia vera, lymphoma/Hodgkin, hyperthyroidism, diabetes, iron deficiency, pregnancy), drug-induced, and psychogenic. Diagnosis: history (onset, duration, triggers, rash), skin exam, and targeted labs (LFTs, renal, CBC, TSH, iron, HIV). Management: treat cause, emollients, antihistamines (histaminergic itch), topical steroids/capsaicin, phototherapy (uremia/cholestasis), and avoidance of triggers; severe/refractory warrants malignancy/endocrine workup.

\medterm{Pseudoxanthoma elasticum}-- An inherited connective tissue disorder caused by mutations in the ABCC6 gene, leading to progressive calcification and fragmentation of elastic fibres in the skin, eyes, and cardiovascular system. Cutaneous findings include yellowish, peau d'orange papules in flexural areas (neck, axillae, groin) resembling plucked chicken skin. Ocular angioid streaks and premature atherosclerosis are important systemic complications.

\medterm{Psoriasis}-- A chronic, immune-mediated inflammatory skin disease characterized by well-demarcated, erythematous plaques with silvery-white scale, affecting approximately 2 to 3 percent of the global population, with a bimodal age of onset (peak at 15 to 25 years and 50 to 60 years), a roughly equal male-to-female ratio, and strong genetic predisposition (concordance rate of 70 percent in monozygotic twins; HLA-Cw6 associated with early-onset psoriasis, present in 60 percent of psoriasis patientpatients), and environmental triggers (trauma, infection, stress, medications including lithium, beta-blockers, chloroquine, interferon, and ACE inhibitors, and smoking) acting on a background of genetic susceptibility. The pathophysiology is driven by the IL-23/TH17 axis (dendritic cell-derived IL-23 activates CD4+ and CD8+ T cells to produce IL-17 and IL-22, which stimulate keratinocyte proliferation and the production of antimicrobial peptides, chemokines, and proinflammatory cytokines, with tumour necrosis factor alpha playing a key amplifying role), leading to the characteristic histological triad of acanthosis (epidermal thickening), hyperkeratosis (thickened stratum corneum with retained nuclei, parakeratosis), and a dense dermal inflammatory infiltrate of T cells, dendritic cells, and neutrophils. The clinical presentation is typically of well-demarcated, erythematous, indurated plaques with a silvery-white micaceous scale, located symmetrically on the extensor surfaces (elbows, knees), scalp, lumbosacral area, and nails.

\medterm{Psychodermatology}-- Psychodermatology addresses the mind-skin connection. Stress exacerbates psoriasis, eczema, acne, and urticaria. Psychiatric disorders (depression, anxiety, body dysmorphic disorder) are associated with skin conditions. Dermatologic conditions cause psychiatric morbidity: psoriasis and depression, acne and anxiety, vitiligo and social phobia. Integrated care addressing both skin and psychological aspects improves outcomes. Cognitive behavioral therapy, mindfulness, and antidepressants may benefimay benefit patients with psychodermatological conditions.

\medterm{Punctate palmoplantar keratoderma}-- A hereditary keratoderma characterised by multiple, small, firm, circumscribed keratotic papules and plaques scattered across the palms and soles, often with an autosomal dominant inheritance pattern. It may present in adolescence or adulthood and can be associated with malignant transformation in rare cases. Treatment includes keratolytics, retinoids, and surgical debulking of large lesions.

\medterm{Purpura}-- A visible, non-blanching red, purple, or dark brown discoloration of the skin or mucous membranes caused by the extravasation of erythrocytes from cutaneous blood vessels into the surrounding dermis or submucosal connective tissue. Because the extravasated erythrocytes lie outside the vascular lumen, purpuric lesions do not blanch when direct pressure is applied with a glass slide (negative diascopy), definitively distinguishing purpura from inflammatory capillary erythema or telangiectasia. Morphemically categorized by diameter into petechiae ($<3\text{--}4\text{ mm}$ pinhead macules) and ecchymoses ($>4\text{--}5\text{ mm}$ or large bruises). Clinically and pathophysiologically, purpura is classified into: (1) Non-palpable (macular) purpura, resulting from platelet disorders (thrombocytopenia, immune thrombocytopenic purpura, thrombotic thrombocytopenic purpura), coagulation factor deficiencies, vascular fragility (senile purpura, scurvy/vitamin C deficiency, systemic amyloidosis, prolonged corticosteroid use), or mechanical trauma; and (2) Palpable purpura, the classic hallmark of cutaneous small-vessel vasculitis (leukocytoclastic vasculitis, IgA vasculitis), wherein vessel wall inflammation and fibrinoid necrosis produce elevated, indurated, purpuric papules. Urgent clinical evaluation is mandated to exclude life-threatening meningococcemia, purpura fulminans/DIC, and systemic necrotizing vasculitides.

\medterm{Purpura fulminans}-- A rapidly progressive, life-threatening condition characterised by extensive, confluent, non-blanching purpuric skin necrosis with haemorrhagic bullae, resulting from acute intravascular thrombosis and disseminated intravascular coagulation. It may be triggered by severe sepsis (particularly meningococcaemia), inherited protein C or S deficiency, or other coagulopathies. It constitutes a medical emergency requiring intensive care, anticoagulation, and treatment of the underlying cause.

\medterm{Pustule}-- A small, circumscribed elevation of skin containing pus (neutrophils) within the epidermis. Causes: bacterial (folliculitis, impetigo, furuncle), acne vulgaris (inflammatory comedones), pustular psoriasis (generalized pustular psoriasis—systemic symptoms), hidradenitis suppurativa, and drug reactions. Sterile pustules occur in psoriasis and palmoplantar pustulosis. Diagnosis: clinical, Gram stain/culture of pus (bacterial), biopsy for atypical. Management: treat cause—antibiotics (topical/systemic) for bacterial and acne, retinoids, steroids/immunosuppressants for pustular psoriasis, and drainage of furuncles/abscesses.

\medterm{Pyoderma Gangrenosum}-- A rare, non-infectious, chronic inflammatory neutrophilic dermatosis characterized by rapidly progressive, agonizingly painful cutaneous ulcerations with classic undermined, violaceous, ragged borders and necrotic, purulent bases. Despite its historical name, it is neither an infectious pyoderma nor a true gangrene. Pathophysiologically, it is driven by dysregulated innate immune signaling, aberrant neutrophil hyperreactivity, and massive dermal neutrophilic infiltration with subsequent tissue destruction, mediated by elevated pro-inflammatory cytokines (IL-1$\beta$, IL-6, IL-8, TNF-$\alpha$). Approximately $50\text{--}70\%$ of cases are associated with underlying systemic illnesses, notably inflammatory bowel disease (Crohn's disease, ulcerative colitis), rheumatoid arthritis, seronegative spondyloarthropathies, and hematologic malignancies (acute myeloid leukemia, myelodysplastic syndrome, monoclonal gammopathy). A hallmark clinical feature is pathergy---the induction or dramatic worsening of ulceration following minor trauma, needle sticks, or surgical debridement; surgical intervention is strictly contraindicated. Diagnosis is clinical and exclusionary, supported by skin biopsy showing dense dermal neutrophilic infiltrates without evidence of infection or vasculitis. First-line medical therapy involves high-dose systemic corticosteroids (prednisone $1\text{--}2\text{ mg/kg/day}$) or cyclosporine, with biologic TNF inhibitors (infliximab) for refractory disease.

\medterm{Pyogenic Granuloma}-- A benign vascular proliferation presenting as a rapidly growing, friable, red papule or nodule that bleeds easily. It commonly occurs on fingers, face, and lips, and may be triggered by trauma or pregnancy (epulis gravidarum). Diagnosis is clinical. Treatment includes shave excision with electrocautery, laser ablation, or topical timolol for small lesions. Recurrence rate is 15 percent.

\medterm{Rash}-- A general term for any visible skin eruption, change in skin color, texture, or appearance. Rashes can be localised or generalised, acute or chronic, and caused by a wide variety of aetiologies: infectious (viral—measles, rubella, varicella, roseola, erythema infectiosum, hand-foot-and-mouth disease, infectious mononucleosis, HIV seroconversion; bacterial—scarlet fever, impetigo, meningococcaemia, Lyme disease, syphilis; fungal—dermatophytosis, candidiasis), allergic/drug-induced (urticaria, maculopapular drug eruption, Stevens–Johnson syndrome/toxic epidermal necrolysis, fixed drug eruption, contact dermatitis), inflammatory (eczema/dermatitis, psoriasis, lichen planus, seborrhoeic dermatitis, pityriasis rosea), autoimmune (lupus erythematosus, dermatomyositis, scleroderma, vasculitis), and neoplastic (cutaneous T-cell lymphoma, Kaposi sarcoma, basal cell carcinoma). Morphological classification is essential for diagnosis: macule (flat, <1 cm), papule (raised, <1 cm), plaque (raised, >1 cm), patch (>1 cm flat), nodule (>1 cm deep), vesicle (<1 cm fluid-filled), bulla (>1 cm fluid-filled), pustule (pus-filled), wheal (transient, oedematous), and petechiae/purpura (non-blanching hemorrhagic). Distribution (central vs peripheral, flexural vs extensor, sun-exposed vs intertriginous) and configuration (annular, linear, reticulate, dermatomal) further narrow the differential. Diagnostic tools: history (onset, duration, fever, sick contacts, medications, travel), dermatoscopy, Wood lamp, skin scraping (KOH prep for fungus, scabies), Tzanck smear (herpes), patch testing, and skin biopsy. Treatment: directed at the underlying cause.

\medterm{Rashes}-- A broad, non-specific clinical term describing any visible, palpable, or symptomatic inflammatory, infectious, neoplastic, or vascular eruption affecting the cutaneous surface. In systematic clinical dermatology, the evaluation of a rash begins with meticulous morphologic characterization, distinguishing primary lesions (the immediate physical expression of the underlying disease process: macules, patches, papules, plaques, vesicles, bullae, pustules, wheals, and nodules) from secondary lesions (consequences of evolution, trauma, or regression: scales, crusts, erosions, ulcers, fissures, lichenification, and atrophy). Clinical evaluation mandates assessment of anatomical distribution (flexural, extensor, photo-exposed, intertriginous, or dermatomal), configuration (annular, linear, reticulated, or targetoid), evolution over time, and associated systemic manifestations (fever, mucosal involvement, lymphadenopathy, arthralgias, or hemodynamic compromise). The presence of mucosal erosions, epidermal sloughing, palpable purpura, or facial edema demands immediate emergency evaluation to rule out life-threatening dermatoses such as Stevens-Johnson syndrome/toxic epidermal necrolysis (SJS/TEN), DRESS syndrome, meningococcemia, necrotizing fasciitis, or staphylococcal scalded skin syndrome.

\medterm{Reactive perforating collagenosis}-- A perforating dermatosis in which dermal collagen is extruded through the epidermis, presenting as discrete, dome-shaped papules with a central keratotic plug, most commonly on the extensor extremities. It may be idiopathic (childhood onset) or acquired in the setting of diabetes mellitus, chronic renal failure, or other systemic diseases. It is often intensely pruritic and may worsen with scratching.

\medterm{Rhinophyma}-- A severe, end-stage subtype of phymatous rosacea occurring predominantly in elderly men. Characterized by progressive nodular thickening, hypertrophy, and lobular enlargement of the nasal skin with prominent, dilated pilosebaceous pores. Pathophysiology: sebaceous gland hyperplasia, fibrosis, and chronic inflammation. Not caused by alcohol abuse (a common historical myth), though alcohol may exacerbate flushing. Treatment: surgical debulking via electrosurgery, carbon dioxide ($CO_2$) laser ablation, or dermabrasion.

\medterm{Ringworm}-- The historical and colloquial lay designation for superficial dermatophyte fungal infections of the stratum corneum (tinea), so named because of its classic annular, ring-like configuration with an active, elevated, erythematous, scaly advancing border and central clearing, rather than infection by any helminthic parasite. The condition is caused by keratinophilic filamentous fungi belonging to three dermatophyte genera: \textit{Trichophyton}, \textit{Microsporum}, and \textit{Epidermophyton}. These organisms secrete specialized keratinases that digest keratin in hair, nails, and the stratum corneum. Clinically, ringworm is classified by anatomical location: tinea corporis (glabrous skin of the trunk and limbs), tinea capitis (scalp, presenting with alopecia and black dots in children), tinea cruris (groin, 'jock itch'), and tinea faciei (face). Transmission occurs through contact with infected individuals (anthropophilic), domestic animals (zoophilic, notably \textit{Microsporum canis} from cats and dogs), or contaminated soil (geophilic). Diagnosis is confirmed by direct microscopy of skin scrapings in $10\text{--}20\%$ potassium hydroxide (KOH), revealing branching, septate hyphae. Treatment consists of topical allylamines (terbinafine) or azoles, with oral antifungal therapy (terbinafine, griseofulvin) reserved for tinea capitis, follicular invasion (Majocchi granuloma), or widespread tinea corporis.

\medterm{Rosacea}-- A chronic inflammatory disorder characterized by persistent or recurrent centrofacial erythema, flushing, telangiectasia, inflammatory papules or pustules, phymatous change, or ocular involvement. It is classified by phenotype rather than by obligatory subtypes. Treatment is individualized and may include trigger avoidance, topical anti-inflammatory or vasoconstrictor therapy, oral tetracyclines, laser treatment, or ophthalmic care.

\medterm{Sarcoidosis Skin Manifestations}-- Cutaneous sarcoidosis affects 25 percent of patients with systemic sarcoidosis. Lesions include erythema nodosum (tender shin nodules, Löfgren syndrome with bilateral hilar lymphadenitis and fever), lupus pernio (violaceous plaques on nose, cheeks, ears), maculopapular eruptions, and scar sarcoidosis (granulomatous infiltration of old scars). Diagnosis requires biopsy showing non-caseating granulomas. Systemic evaluation is essential.

\medterm{Scabies}-- A highly contagious, intensely pruritic parasitic infestation of the skin caused by the mite Sarcoptes scabiei var. hominis. The hallmark symptom is severe, generalized pruritus that characteristically worsens at night, driven by a hypersensitivity reaction to the mite's proteins, eggs, and feces. Clinical findings include erythematous papules, vesicles, excoriations, and pathognomonic burrows (linear or wavy tracts) located in web spaces of the fingers, flexor wrists, axillae, umbilicus, and genitalia. Crusted scabies is a highly contagious hyperkeratotic variant occurring in immunocompromised patients, harboring millions of mites. Diagnosis is confirmed by identifying the mite, eggs, or scybala on microscopic examination of skin scrapings. First-line treatments include topical permethrin 5 percent cream applied from the neck down overnight (repeated in 1 week) or oral ivermectin. Simultaneous treatment of close contacts and hot laundering of clothing and bedding are mandatory to prevent reinfestation.

\medterm{Scar Management}-- Scar management aims to minimize hypertrophic and keloid formation. Early intervention includes silicone gel or sheeting (first-line evidence-based treatment), pressure therapy (especially for burn scars), and massage. Intralesional corticosteroids (triamcinolone) for hypertrophic scars and keloids. 5-fluorouracil injection for resistant keloids. Laser therapy (pulsed dye laser for erythema, fractional CO2 for texture). Surgical revision for mature, problematic scars. Sun protection prevents hypprevents hyperpigmentation of scars and improves overall appearance.

\medterm{Scarring Alopecia}-- Cicatricial alopecia; an umbrella term for a diverse group of inflammatory hair-loss disorders characterized by the permanent, irreversible destruction of hair follicles and their replacement by fibrous scar tissue, leading to permanent baldness and complete loss of visible follicular ostia. Pathophysiologically, the primary inflammatory insult targets the epithelial stem cell reservoir located within the follicular bulge region of the outer root sheath near the arrector pili muscle insertion; once bulge stem cells are destroyed, follicular regeneration is permanently abolished. Primary cicatricial alopecias are classified histologically by the predominant inflammatory infiltrate: (1) Lymphocytic: lichen planopilaris (LPP), frontal fibrosing alopecia (FFA), chronic cutaneous discoid lupus erythematosus (DLE), and central centrifugal cicatricial alopecia (CCCA); (2) Neutrophilic: folliculitis decalvans and dissecting cellulitis of the scalp; and (3) Mixed: acne keloidalis nuchae and acne necrotica. Secondary cicatricial alopecia results from physical trauma, chemical burns, ionizing radiation, severe infections (tinea capitis kerion), or morpheaform neoplasms. Early scalp punch biopsy (horizontal and vertical sectioning) and dermoscopy (peripilar casts, follicular erythema, absent ostia) are vital to arrest progression with potent topical/intralesional corticosteroids, antimalarials, or doxycycline before complete follicular obliteration occurs.

\medterm{Schnitzler Syndrome}-- A rare autoinflammatory syndrome characterized by chronic urticarial rash, intermittent fever, bone pain, arthralgia or arthritis, and elevated inflammatory markers, usually with a monoclonal IgM protein and less commonly IgG. Evaluation includes serum protein studies, inflammatory markers, blood counts, and assessment for bone involvement. Interleukin-1 blockade, such as anakinra or canakinumab, is specialist-directed treatment.

\medterm{Schwannoma}-- A benign peripheral nerve sheath tumor arising from Schwann cells. It presents as a solitary, slow-growing, painless nodule along a peripheral nerve course. It can be moved laterally but not longitudinally along the nerve. Vestibular schwannoma (acoustic neuroma) arises from CN VIII. Diagnosis is by MRI. Treatment is surgical enucleation preserving nerve function. Schwannomas are encapsulated and separable from the nerve.

\medterm{Scleromyxedema}-- A rare, chronic, progressive mucinosis characterised by generalised papular and sclerodermoid thickening of the skin with firm, waxy papules in a "peau d'orange" pattern, most prominently on the face, neck, and hands. It is almost always associated with a paraproteinaemia (monoclonal IgG-lambda) and may involve internal organs including the cardiovascular, gastrointestinal, and nervous systems. Treatment includes intravenous immunoglobulin and systemic immunosuppression.

\medterm{Scrotal calcinosis}-- The deposition of calcium within the scrotal skin, presenting as multiple, firm, whitish, painless nodules on the scrotum, which may be solitary or numerous and may ulcerate. The exact aetiology is debated and may involve dystrophic calcification of pre-existing epidermal inclusion cysts, dermal cysts, or eccrine duct abnormalities. Surgical excision is the definitive treatment for symptomatic or cosmetically distressing lesions.

\medterm{Sebaceous Gland}-- A specialized microscopic exocrine holocrine gland embedded in the dermis of mammalian skin, almost universally associated with hair follicles to form the pilosebaceous unit (with the exception of free sebaceous glands on glabrous sites, such as Fordyce spots on the oral mucosa and vermilion lips, Meibomian glands in the eyelids, and Tyson glands on the prepuce). Sebaceous glands are distributed throughout the cutaneous surface with their highest densities located on the face, scalp, upper chest, and mid-back. Histologically, lipid-laden sebocytes synthesize and store lipids until undergoing programmed apoptotic rupture (holocrine secretion) to discharge sebum---a complex lipid mixture composed of triglycerides ($40\text{--}60\%$), wax esters ($20\text{--}30\%$), squalene ($10\text{--}15\%$), free fatty acids, and cholesterol---into the follicular infundibulum. Sebum functions to waterproof the skin, maintain stratum corneum hydration, transport lipophilic antioxidants (vitamin E) to the skin surface, and exert innate antimicrobial activity. Sebaceous gland activity is androgen-dependent, surging during neonatal life and puberty under dihydrotestosterone (DHT) stimulation, playing a pivotal role in the pathogenesis of acne vulgaris, seborrheic dermatitis, and sebaceous hyperplasia.

\medterm{Seborrheic Dermatitis}-- A chronic inflammatory skin condition affecting sebum-rich areas: scalp, face (nasolabial folds, eyebrows), chest, and intertriginous areas. Presents with erythematous patches with greasy, yellowish scales. Dandruff is a mild form. More common and severe in Parkinson disease, HIV/AIDS, and Down syndrome. Related to Malassezia yeast colonization and immune response. Treatment: ketoconazole shampoo/cream, selenium sulfide, zinc pyrithione, topical corticosteroids, and calcineurin inhibitors for facial involvement.

\medterm{Seborrheic Keratosis}-- A benign epidermal tumor appearing as stuck-on, waxy, hyperpigmented plaques with a verrucous surface. They are extremely common in older adults and increase in number with age. Seborrheic keratoses are not premalignant. The sign of Leser-Tralat indicates sudden eruption of multiple seborrheic keratoses, which may signal internal malignancy (gastric adenocarcinoma). Treatment is cosmetic removal by cryotherapy, curettage, or shave excision.

\medterm{Seborrhoea}-- A clinical condition characterized by excessive production and secretion of sebum (hyperseborrhea) by hyperactive sebaceous glands, producing a visibly oily, glistening, and greasy appearance of the skin and hair, often accompanied by widened follicular orifices and increased susceptibility to comedone formation. Seborrhea is physiologically mediated by androgenic hormones (testosterone and $5\alpha$-dihydrotestosterone, DHT) binding to nuclear androgen receptors within sebocytes, stimulating cellular proliferation and lipogenesis. Hyperseborrhea represents a primary predisposing state for several cutaneous disorders, including acne vulgaris, seborrheic dermatitis, and \textit{Malassezia} folliculitis. It is most pronounced across anatomical areas with the highest density of sebaceous units: the face (forehead, nose, nasolabial folds, and chin, constituting the 'T-zone'), scalp, retroauricular sulci, presternal chest, and interscapular back. Pathological seborrhea is frequently observed in hyperandrogenic states (polycystic ovary syndrome [PCOS], congenital adrenal hyperplasia) and neurodegenerative conditions such as Parkinson disease ('seborrheic mask' resulting from autonomic dysregulation). Management involves topical sebum-regulating and keratolytic agents (niacinamide, salicylic acid, topical retinoids) and systemic anti-androgens (spironolactone, oral contraceptives) in appropriate clinical contexts.

\medterm{Secondary syphilis}-- The systemic dissemination stage of Treponema pallidum infection, occurring weeks to months after the primary chancre, characterised by a generalised maculopapular rash (including the palms and soles), mucous patches, condylomata lata, lymphadenopathy, and systemic symptoms such as fever and malaise. It is highly infectious and represents the peak of spirochaetaemia. Serological testing (RPR/VDRL and FTA-ABS) confirms the diagnosis, and treatment is intramuscular benzathine penicillin.

\medterm{Second-Degree Burn}-- Partial-thickness burn involving the epidermis and dermis: superficial partial-thickness (blisters, moist, painful, pink) vs. deep partial-thickness (blanching, painful, may have thrombosed capillaries). Causes: thermal (scalds, flames), chemical, electrical. Management: cool water, analgesia, blisters (debride/leave depending), sterile dressings, infection prevention (topical—silver sulfadiazine, mupirocin), tetanus prophylaxis, and fluids for significant surface area; deep partial-thickness may require excision and grafting. Complications: infection, scarring, contractures, and fluid loss. Full-thickness areas need surgical management.

\medterm{Sentinel Lymph-Node Biopsy}-- A specialized surgical and pathological staging procedure used in clinical cutaneous oncology to evaluate whether micrometastatic disease has spread from a primary cutaneous malignancy to the regional lymphatic basin. The sentinel lymph node (SLN) is defined as the first draining node or nodes in a regional lymphatic basin that receives direct afferent lymphatic drainage from the primary anatomical tumor site. In cutaneous malignant melanoma, SLNB is established as the standard of care for staging in patients with intermediate-thickness melanomas (Breslow depth $1.0\text{--}4.0\text{ mm}$) or high-risk thin melanomas (Breslow depth $0.8\text{--}1.0\text{ mm}$ or $<0.8\text{ mm}$ with ulceration or other high-risk features), as well as selected high-risk cutaneous squamous cell carcinomas and Merkel cell carcinomas. The technique utilizes preoperative lymphoscintigraphy with radioactive technetium-99m sulfur colloid combined with intraoperative injection of isosulfan blue or indocyanine green dye around the primary tumor scar; a handheld gamma counter and direct visualization guide targeted excision. Histopathologic serial sectioning and immunohistochemical staining (using S100, SOX10, Melan-A/MART-1, HMB-45) identify submicroscopic tumor deposits, providing critical prognostic risk stratification and guiding adjuvant systemic immunotherapies.

\medterm{Skin Aging}-- intrinsic (chronological) and extrinsic (photoaging) components. Photoaging causes wrinkle formation, telangiectasias, lentigines, leather-like texture, and actinic keratoses. Intrinsic aging causes thinning, wrinkling, and delayed wound healing. Prevention includes sun protection (broad-spectrum SPF 30 or greater daily). Treatment includes topical retinoids (tretinoin), vitamin C, alpha hydroxy acids, chemical peels, laser therapy, and fillers. Tretinoin is the most evidencemost evidence-based topical treatment for photoaging, with demonstrated improvement in fine wrinkles, roughness, and mottled hyperpigmentation after 3 to 6 months of regular use.

\medterm{Skin Aging and Appearance}-- Skin aging is a complex biological process influenced by both intrinsic factors, such as genetics and natural aging, and extrinsic factors, particularly ultraviolet radiation exposure. Chronological aging leads to thinning, dryness, and loss of elasticity, while photoaging from sun exposure causes wrinkles, pigmentation changes, and textural irregularities. Preventive measures include consistent sun protection, moisturizing, and avoiding smoking, while treatments such as topical retinoids, antioxidants, laser therapy, and injectable fillers can help improve the appearance of aging skin.

\medterm{Skin Allergies and Contact Dermatitis}-- Inflammatory skin reactions triggered by exposure to allergens or irritants that come into direct contact with the skin. Allergic contact dermatitis results from a delayed immune response to substances such as nickel, poison ivy, or fragrances, while irritant contact dermatitis occurs when chemicals or physical agents damage the skin barrier directly. Symptoms include redness, itching, swelling, and blistering; management involves identifying and avoiding triggers, using topical corticosteroids, and maintaining skin barrier function with moisturizers.

\medterm{Skin Allergy Testing}-- prick testing (immediate hypersensitivity, IgE-mediated) and intradermal testing (more sensitive, used for drug and venom allergy). Prick testing introduces allergen extract into the epidermis; wheal and flare reaction at 15 minutes indicates sensitization. Intradermal testing injects dilute allergen into the dermis. Both tests require stopping antihistamines 3 to 7 days before testing. Negative predictive value is high; positive results require clinical correlationclinical correlation to confirm relevance.

\medterm{Skin and Aging}-- Skin aging involves intrinsic and extrinsic changes. Intrinsic aging causes thinning, wrinkling, and fragility. Extrinsic (photo)aging causes wrinkles, telangiectasias, lentigines, and actinic keratoses. Collagen loss, elastin degradation, and decreased cell turnover contribute to aging skin. Prevention includes sun protection, smoking cessation, and antioxidants. Treatment includes retinoids, vitamin C, alpha hydroxy acids, chemical peels, laser therapy, and fillers. Skin cancer screening is esis essential as part of the regular check-up for elderly patients.

\medterm{Skin and Environment}-- Environmental factors significantly affect skin health. Pollution causes oxidative stress, inflammation, and accelerated aging. Particulate matter penetrates skin, generating free radicals. Blue light from screens causes hyperpigmentation, particularly in darker skin tones. Indoor heating and air conditioning cause dryness. Water quality (hard water) affects eczema. Environmental protection includes antioxidant skincare, cleansing, and sun protection.

\medterm{Skin Appendages}-- hair follicles, sebaceous glands, eccrine sweat glands, and apocrine sweat glands. Hair follicles cycle through anagen (growth), catagen (transition), and telogen (resting) phases. Sebaceous glands produce sebum, an oily substance lubricating skin and hair. Eccrine glands (distributed over entire body) produce watery sweat for thermoregulation. Apocrine glands (axillae, groin) produce viscous secretion active at puberty, responsible for body odor when metabolized by bacteby bacteria that break down fatty acids and proteins in the secretion.

\medterm{Skin Barrier Function}-- The skin barrier prevents water loss and protects against environmental insults. The stratum corneum acts as a physical barrier, while the acid mantle (pH 4 to 6) provides chemical defense. Lipids (ceramides, cholesterol, free fatty acids) in the intercellular space form a waterproof barrier. Tight junctions in the stratum granulosum prevent paracellular transport. Barrier dysfunction is central to atopic dermatitis, psoriasis, and ichthyosis. Emollients restore barrier function by filling interintercellular lipids, while moisturisers reduce water loss and improve barrier function.

\medterm{Skin Biopsy Methods}-- Skin biopsy methods vary by lesion characteristics. Punch biopsy (2 to 6 mm circular blade) provides full-thickness specimen for most diagnoses. Shave biopsy (razor or scalpel) removes superficial lesions; deep shave (saucerization) includes subcutaneous fat for suspected melanoma. Excisional biopsy removes entire lesion with narrow margins, preferred for suspected melanoma. Incisional biopsy removes representative portion of large lesions. Site and technique should be chosen to optimize diagnosdiagnostic yield. Proper anesthesia and haemostasis are important for optimal specimen quality.

\medterm{Skin Biopsy Techniques}-- Skin biopsy is selected according to the suspected diagnosis and lesion depth. Punch biopsy samples full-thickness skin; shave biopsy samples superficial lesions; excisional biopsy removes an entire small lesion; and incisional biopsy samples part of a large lesion. For suspected autoimmune blistering disease, perilesional skin is submitted for direct immunofluorescence, while lesional tissue is submitted for routine histology.

\medterm{Skin Cancer Epidemiology}-- Skin cancer epidemiology reveals increasing incidence worldwide. Basal cell carcinoma is the most common cancer (1 million cases annually in the US). Squamous cell carcinoma is second (approximately 200,000 cases annually). Melanoma is less common but most deadly (about 100,000 new cases annually in the US). Incidence of melanoma has increased 200 percent over the past 30 years. UV radiation is the primary risk factor. Skin cancer screening and sun protection are public health priorities.

\medterm{Skin Cancer Prevention}-- Skin cancer prevention focuses on UV radiation protection. Broad-spectrum sunscreen (SPF 30 or greater) applied daily and reapplied every 2 hours during sun exposure. Protective clothing, hats, and sunglasses. Seeking shade during peak UV hours (10 AM to 4 PM). Avoiding tanning beds. Vitamin D supplementation for those with strict sun avoidance. Sunscreen reduces melanoma risk by 50 percent with intermittent use. Sunscreen use in children is particularly important.

\medterm{Skin Cancer Screening}-- Skin cancer screening uses the ABCDE criteria for melanoma, the ugly duckling sign, and total body skin examination. Dermoscopy (dermatoscopy) magnifies skin lesions revealing structures invisible to the naked eye, improving diagnostic accuracy. Skin self-examination teaches patients to monitor moles for changes. Annual screening is recommended for high-risk patients (personal or family history of melanoma, numerous atypical nevi, immunosuppression).

\medterm{Skin Cancer Staging}-- Skin cancer staging guides treatment and prognosis. Melanoma: Breslow thickness (most important prognostic factor), ulceration, mitotic rate, sentinel lymph node status (AJCC 8th edition). Basal cell carcinoma: rarely metastasizes; high-risk features include size greater than 2 cm, location on face, perineural invasion, and recurrent tumors. Squamous cell carcinoma: AJCC staging based on size, depth, perineural invasion, and nodal metastasis. High-risk features include thickness greater than 2 mgreater than 2 mm for SCC and greater than 2 cm for BCC. Sentinel lymph node biopsy is indicated for high-risk SCC and melanoma.

\medterm{Skin Cancer Surgery}-- Skin cancer surgical techniques maximize cure while minimizing morbidity. Wide local excision with appropriate margins (BCC: 4 mm, SCC: 4 to 6 mm, melanoma: based on Breslow thickness). Mohs micrographic surgery provides 99 percent cure rate for BCC and 97 percent for SCC with tissue preservation. Sentinal lymph node biopsy for melanoma greater than 1 mm or with ulceration. Reconstruction ranges from primary closure to flaps and grafts.

\medterm{Skin Flaps}-- Skin flaps transfer tissue with its blood supply. Rotation flaps rotate adjacent tissue into defect. Advancement flaps slide tissue forward. Transposition flaps move tissue across intervening skin. Random pattern flaps rely on dermal plexus; axial pattern flaps have named vascular pedicle. Flaps provide better color and texture match than grafts for facial reconstruction. Design follows relaxed skin tension lines and aesthetic subunits.

\medterm{Skin Grafts}-- Skin grafts replace lost skin. Split-thickness grafts (epidermis and part of dermis) heal by re-epithelialization of donor site, allowing reuse. Full-thickness grafts (entire dermis) provide better cosmetic results but require primary closure of donor site. Sheet grafts provide smooth surface; mesh grafts cover larger areas but have quilted appearance. Graft take requires vascular ingrowth from wound bed (plasmatic imbibition day 1 to 2, inosculation day 2 to 3, revascularization day 3 to 5).

\medterm{Skin Immune System}-- The skin immune system (SALT: skin-associated lymphoid tissue) provides immunological surveillance. Langerhans cells (dendritic cells in epidermis) capture antigens and present them to T cells. Dermal dendritic cells, mast cells, and resident T cells form an innate and adaptive immune network. Cytokines (IL-1, IL-6, TNF-alpha, IL-17, IL-23) mediate inflammatory responses. The skin immune system is involved in contact dermatitis, psoriasis, atopic dermatitis, and skin cancer immunosurveillance.

\medterm{Skin Microbiome}-- The skin microbiome is the community of microorganisms (bacteria, fungi, viruses) residing on the skin. Staphylococcus aureus colonization is increased in atopic dermatitis and psoriasis. Cutibacterium acnes plays a role in acne. Diversity is decreased in diseased skin. Probiotics and prebiotics are being studied for dermatologic conditions. Topical antimicrobials and emollients modulate the microbiome. Disruption of the microbiome contributes to skin disease pathogenesis.

\medterm{Skin of Color}-- Skin of color (Fitzpatrick IV to VI) has unique dermatologic considerations. It is more prone to hyperpigmentation disorders (post-inflammatory hyperpigmentation, melasma, acanthosis nigricans) and keloid formation. It is less prone to skin cancer but presents with more advanced disease when it occurs. Dermatoscopic features differ in pigmented skin. Misdiagnosis is common: vitiligo appears as ash-gray macules, lupus may present with hyperpigmentation rather than erythema. Culturally competent dculturally competent dermatological care acknowledges these differences and ensures accurate diagnosis and appropriate management for patients with skin of color.

\medterm{Skin Pigmentation}-- Determined by melanin produced by melanocytes in the stratum basale. Eumelanin (brown-black) and pheomelanin (red-yellow) are the two types. Melanin production is stimulated by UV radiation through melanocyte-stimulating hormone (MSH). Fitzpatrick skin type classification (I to VI) describes pigmentation and sun response. Hyperpigmentation disorders include melasma, post-inflammatory hyperpigmentation, and lentigines. Hypopigmentation disorders include vitiligo, pityriasis apityriasis alba (post-inflammatory hypopigmentation) and albinism (congenital hypopigmentation, oculocutaneous or ocular). Treatment of hyperpigmentation includes sun protection, topical hydroquinone, azelaic acid, kojic acid, vitamin C, retinoids, chemical peels, and laser therapy.

\medterm{Skin Tags}-- Benign fibroepithelial polyps occurring in intertriginous areas (neck, axillae, groin, under breasts). They are associated with obesity, insulin resistance, and pregnancy. Skin tags are asymptomatic but may become irritated or bleed. Treatment includes snipping, cryotherapy, electrocautery, and shave excision. They are harmless and require no biopsy unless atypical.

\medterm{Soft-tissue calcinosis}-- The deposition of insoluble calcium phosphate salts in soft tissues such as skin, subcutaneous tissue, muscles, and tendons, which may be metastatic (due to hypercalcaemia or hyperphosphataemia), dystrophic (in damaged or devitalised tissue), or idiopathic. It presents as firm, palpable nodules that may ulcerate and extrude calcareous material. Treatment addresses the underlying metabolic disturbance when present.

\medterm{Solar Urticaria}-- A rare form of inducible urticaria triggered by ultraviolet or visible light, producing pruritic erythema and wheals within minutes on exposed skin. Symptoms usually improve after removal from the light source. Diagnosis may involve phototesting. Management includes strict photoprotection, protective clothing, antihistamines, and specialist-directed phototherapy or other treatment for refractory disease.

\medterm{Squamous Cell Carcinoma}-- Cutaneous squamous cell carcinoma is a malignant keratinocyte tumor commonly arising on chronically sun-exposed skin or in areas of chronic inflammation or immunosuppression. It may present as a scaly, indurated, hyperkeratotic papule, plaque, or ulcer. Diagnosis requires biopsy; risk of recurrence or metastasis depends on tumor size, depth, differentiation, perineural invasion, location, immune status, and other pathologic features.

\medterm{Staphylococcal scalded skin}-- See staphylococcal scalded skin syndrome. A toxin-mediated disease caused by exfoliative toxins of Staphylococcus aureus, resulting in widespread erythema and flaccid bullae that rupture easily, leaving painful, moist, erythematous denuded areas resembling a scald burn. It predominantly affects infants and young children. Treatment involves intravenous antistaphylococcal antibiotics and supportive wound care.

\medterm{Stevens-Johnson Syndrome}-- A rare, acute, and life-threatening dermatological emergency and severe adverse cutaneous drug reaction characterized by extensive epidermal necrosis, full-thickness detachment, and severe mucositis. Mechanistically, SJS represents a cell-mediated, type IV cytotoxic hypersensitivity response wherein drug-antigen-specific cytotoxic $CD8^+$ T lymphocytes and natural killer (NK) cells release massive quantities of granulysin, perforin, granzyme B, and Fas ligand (FasL), precipitating widespread keratinocyte apoptosis. Common culprit drugs include allopurinol, aromatic antiepileptics (carbamazepine, phenytoin, lamotrigine), sulfonamide antimicrobials (trimethoprim-sulfamethoxazole), nevirapine, and oxicam NSAIDs, strongly associated in certain populations with specific human leukocyte antigen alleles (e.g., HLA-B*1502, HLA-B*5801). Clinical presentation begins with a febrile flu-like prodrome, followed by painful, atypical, purpuric targetoid macules that rapidly coalesce into flaccid blisters and detach under light lateral pressure (positive Nikolsky sign). Severe mucosal erosions affect at least two anatomical sites (oral, ocular, urogenital). The disease spectrum is categorized by detached body surface area (BSA): SJS involves $<10\%$ BSA; SJS/TEN overlap involves $10\text{--}30\%$ BSA; and Toxic Epidermal Necrolysis (TEN) involves $>30\%$ BSA. Prognosis is assessed using the SCORTEN score. Emergency therapy requires prompt culprit drug cessation, burn-unit admission, fluid and electrolyte management, non-adherent wound dressings, systemic immunomodulation (cyclosporine, etanercept), and urgent ophthalmologic care to prevent cicatricial blindness.

\medterm{Subcorneal pustular dermatosis}-- See Sneddon-Wilkinson disease. A chronic, relapsing, sterile pustular eruption characterised by flaccid, pus-filled blisters that accumulate beneath the stratum corneum, typically on the trunk and intertriginous areas. It is often associated with IgA gammopathy or monoclonal gammopathy and may be mistaken for pustular psoriasis. Dapsone is the treatment of choice.

\medterm{Subcutaneous Tissue}-- Subcutaneous tissue (hypodermis, superficial fascia) lies beneath the dermis and connects skin to underlying structures. It contains adipose tissue (fat), blood vessels, nerves, and the superficial muscular aponeurotic system (SMAS). Subcutaneous fat provides insulation, cushioning, and energy storage. The thickness varies by body site, age, sex, and nutritional status. Liposuction removes subcutaneous fat. Subcutaneous injections (insulin, heparin, vaccines) deposit medication in this layer.

\medterm{Sun Protection (Photoprotection)}-- Sun protection involves measures taken to reduce exposure to ultraviolet radiation from the sun, which is the primary environmental risk factor for skin cancer and photoaging. Recommended strategies include applying broad-spectrum sunscreen with a sun protection factor of at least 30, wearing protective clothing and wide-brimmed hats, seeking shade during peak ultraviolet hours between 10 a.m. and 4 p.m., and avoiding tanning beds. Consistent sun protection throughout life significantly reduces the risk of melanoma, squamous cell carcinoma, and basal cell carcinoma, as well as preventing premature skin aging.

\medterm{Sunburn}-- Acute erythema (and later edema, pain, blistering, peeling) of skin from excessive ultraviolet (UV) radiation (UVB predominantly), causing DNA damage and inflammatory response. Risk: repeated sunburn increases skin cancer risk (melanoma, basal/squamous cell carcinoma) and photoaging. Features: erythema 3-5 hours post-exposure, maximal at 12-24 h, pain, tenderness; severe: blistering, systemic symptoms (fever, chills, dehydration). Management: cool compresses, emollients (aloe), NSAIDs/analgesics, hydration; severe with blistering: wound care, fluids; complications: infection, pigmentary change. Prevention: sunscreen (broad-spectrum SPF 30+), protective clothing, avoid peak sun, and photoprotection in photosensitive individuals.

\medterm{Sweat Gland}-- A specialized tubular epithelial exocrine appendage situated within the dermis and subcutaneous tissue of mammalian skin, responsible for thermoregulation, electrolyte balance, and chemical signaling, categorized into two anatomically and functionally distinct types: eccrine and apocrine glands. Eccrine sweat glands (numbering 2 to 4 million) are distributed across virtually the entire cutaneous surface, reaching their highest densities on the palms, soles, and forehead. Innervated by sympathetic postganglionic cholinergic fibers, eccrine glands secrete a clear, hypotonic watery solution (containing water, $Na^+$, $Cl^-$, potassium, urea, and antimicrobial peptides such as dermcidin) directly onto the epidermal surface, evaporating to dissipate body heat during physical exertion or thermal stress. Apocrine sweat glands are restricted to hair-bearing intertriginous zones (axillae, anogenital region, areolae, with modified ceruminous and ciliary Moll glands). Innervated by sympathetic adrenergic fibers, apocrine glands open into the hair follicle infundibulum and secrete a viscous, milky, protein- and lipid-rich fluid that is decomposed by commensal cutaneous bacteria (primarily \textit{Corynebacterium} species) to generate characteristic axillary body odor (bromhidrosis). Pathology includes hyperhidrosis, anhidrosis, and hidradenitis suppurativa.

\medterm{Syphilitic alopecia}-- Non-scarring hair loss occurring as a manifestation of secondary syphilis, presenting with diffuse thinning or patchy, moth-eaten alopecia of the scalp, eyebrows, or beard. It is thought to result from direct spirochaetal invasion of hair follicles during the spirochaetaemic phase. Diagnosis requires a high index of suspicion and serological confirmation, and hair regrows with appropriate penicillin treatment.

\medterm{Systemic Antifungals}-- Systemic antifungals treat extensive or refractory fungal infections. Terbinafine is first-line for onychomycosis (6 weeks for fingernails, 12 weeks for toenails) and tinea capitis. Itraconazole pulse therapy (200 mg twice daily for 1 week per month) treats onychomycosis. Fluconazole treats tinea corporis, cruris, and versicolor. Griseofulvin treats tinea capitis (requires 12 to 16 weeks). Voriconazole and posaconazole treat invasive fungal infections. Hepatic function monitoring is required.

\medterm{Systemic Sclerosis}-- Scleroderma; a complex, multisystem autoimmune connective-tissue disease characterized by the triad of widespread microvascular obliteration, dysregulated innate and adaptive immune activation with autoantibody production, and progressive cutaneous and visceral fibrosis. Pathogenesis involves endothelial injury, perivascular lymphocytic infiltration, and uncontrolled activation of dermal fibroblasts driven by transforming growth factor-beta (TGF-$\beta$), platelet-derived growth factor (PDGF), and connective tissue growth factor (CTGF), precipitating massive deposition of dense extracellular collagen. Cutaneously, it evolves through an early edematous phase ('puffy fingers') to an indurated, bound-down sclerotic phase, producing sclerodactyly, digital flexion contractures, 'salt-and-pepper' depigmentation, telangiectasias, calcinosis cutis, and facial changes (microstomia, radial furrowing around the lips). Divided clinically into: (1) Limited cutaneous systemic sclerosis (lcSSc, historically CREST syndrome: Calcinosis, Raynaud phenomenon, Esophageal dysmotility, Sclerodactyly, Telangiectasia), with skin thickening restricted to the distal extremities and face, strongly associated with anti-centromere antibodies (ACA) and pulmonary arterial hypertension; and (2) Diffuse cutaneous systemic sclerosis (dcSSc), featuring rapid proximal skin involvement (trunk, upper arms, thighs), anti-topoisomerase I (anti-Scl-70) or anti-RNA polymerase III antibodies, and high risks of interstitial lung disease and scleroderma renal crisis. Management entails vasodilators (calcium channel blockers, PDE-5 inhibitors) for Raynaud's, ACE inhibitors for renal crisis, and immunosuppressive therapy (mycophenolate mofetil).

\medterm{Telogen Effluvium}-- Diffuse hair shedding caused by premature entry of hair follicles into the telogen (resting) phase. Triggers include severe illness, surgery, childbirth, rapid weight loss, emotional stress, medications (beta-blockers, retinoids, anticoagulants), and nutritional deficiencies (iron, zinc, vitamin D). Acute TE resolves within 6 months; chronic TE lasts longer than 6 months. Diagnosis is clinical with positive hair pull test and normal scalp examination. Treatment addresses theaddresses the underlying trigger; recovery of normal hair growth occurs over 3 to 6 months, though the condition may become chronic if the triggering factor persists.

\medterm{Third-Degree Burn}-- Full-thickness burn destroying the epidermis and dermis (into subcutaneous tissue): wound is white, leathery, charred, painless (nerve destruction), and non-blanching; no spontaneous healing—requires skin grafting. Assessment: percent BSA (rule of nines), inhalation injury (facial burns, soot, hoarseness—intubation), and depth. Management: advanced burn care—fluid resuscitation (Parkland formula: 4 mL/kg x \%BSA crystalloid), analgesia, escharotomy (circumferential—compartment), wound debridement and grafting, infection control (topical antimicrobials), nutritional support, and ICU care. Complications: infection/sepsis, hypovolemic shock, contractures, hypertrophic scarring, and multiorgan failure.

\medterm{Tinea barbae}-- Dermatophyte infection of the beard and mustache area, presenting with erythematous, scaly, follicular papulopustules or deep inflammatory nodules, often acquired from contact with infected animals or through grooming. It may be inflammatory (kerion-like) or superficial, and differential diagnosis includes bacterial folliculitis and acne. Treatment involves topical or systemic antifungals depending on severity.

\medterm{Tinea Capitis}-- A dermatophyte infection of the scalp and hair shafts, most common in children. It may cause scaling, alopecia, broken hairs, black dots, or an inflammatory kerion. Diagnosis uses microscopy, culture, or other validated testing; treatment requires an oral antifungal because topical therapy alone does not adequately penetrate infected hair shafts.

\medterm{Tinea circinata}-- An older term for tinea corporis (ringworm) of the body, presenting as annular, erythematous, scaly plaques with active, raised borders and central clearing. It is caused by dermatophyte fungi and may be acquired through direct contact with infected humans, animals, or fomites. Treatment includes topical antifungals for limited disease and oral antifungals for extensive or refractory infections.

\medterm{Tinea Corporis}-- A dermatophyte infection of glabrous skin, presenting as annular, erythematous, scaly plaques with raised, active borders and central clearing. It is contagious through direct contact with infected persons, animals, or contaminated objects. Diagnosis is by KOH preparation showing branching septate hyphae. Treatment includes topical antifungals (terbinafine, clotrimazole, miconazole) for 2 to 4 weeks. Oral terbinafine or itraconazole for extensive or refractory disease.

\medterm{Tinea (Dermatophytosis)}-- Superficial fungal infection by dermatophytes. Tinea corporis (ringworm): annular, scaly, erythematous plaque with central clearing. Tinea pedis (athlete's foot): interdigital scaling, maceration, and fissuring. Tinea cruris (jock itch): erythematous, pruritic rash in groin with raised border. Tinea capitis: scaly scalp patches with broken hairs, more common in children; may cause kerion (inflammatory mass). Diagnosis: KOH preparation and culture. Treatment: topical terbinafine or azoles for mild to moderate; systemic terbinafine, itraconazole for extensive, scalp, or nail involvement.

\medterm{Tinea faciei}-- Dermatophyte infection of the facial skin, often presenting with erythematous, scaly, annular or ill-defined plaques that may mimic eczema or rosacea, leading to frequent misdiagnosis. It is commonly caused by Trichophyton or Microsporum species and may be acquired from pets, other humans, or self-inoculation from other body sites. Topical antifungal agents are usually effective.

\medterm{Tinea imbricata}-- A chronic dermatophyte infection caused by Trichophyton concentricum, endemic in the South Pacific and Southeast Asia, characterised by widespread, concentric, overlapping, polycyclic scaly rings producing a "wood-grain" pattern on the trunk and limbs. Unlike most tinea corporis, it is highly contagious within households and is difficult to eradicate, often requiring prolonged systemic antifungal therapy.

\medterm{Tinea Incognito}-- A classic dermatological diagnostic challenge wherein the typical clinical morphology of a superficial dermatophyte infection (tinea) is profoundly altered, masked, or suppressed due to the inappropriate, misdiagnosed application of topical corticosteroids, calcineurin inhibitors, or systemic immunosuppressive agents. Corticosteroids suppress the local cell-mediated immune response, reducing erythema, scaling, and the characteristic raised, active, annular advancing border. Consequently, the infection lacks its classic ringworm appearance, presenting atypically as ill-defined, diffuse, erythematous macules, papules, pustules, or lichenified plaques that mimic atopic dermatitis, discoid lupus, rosacea, or psoriasis. Furthermore, local immunosuppression permits unrestricted fungal proliferation, allowing the dermatophyte (most commonly \textit{Trichophyton rubrum} or \textit{Trichophyton mentagrophytes}) to invade deeply into hair follicles, inducing a follicular pustular dermatosis known as Majocchi granuloma. Cessation of the topical steroid triggers an intense inflammatory rebound flare. Diagnosis requires a high index of clinical suspicion and confirmation via potassium hydroxide (KOH) examination or fungal culture of scale and epilated hairs, which typically reveal abundant fungal hyphae. Management involves stopping the offending topical steroid and administering appropriate topical or oral antifungal therapy.

\medterm{Tinea nigra}-- A superficial fungal infection of the stratum corneum caused by Hortaea werneckii, presenting as a painless, brown to black, non-scaly macule or patch, most commonly on the palms or soles. It is acquired in tropical and subtropical regions and is often confused with acral melanoma. Diagnosis is confirmed by KOH preparation revealing brown, septate hyphae.

\medterm{Tinea Pedis}-- A dermatophyte infection of the feet presenting in interdigital, moccasin, vesiculobullous, or ulcerative patterns. It may coexist with tinea unguium. Diagnosis is clinical or mycologic when uncertain; treatment uses topical antifungals for limited disease and oral therapy for extensive, refractory, or nail-associated infection.

\medterm{Tinea Versicolor}-- A superficial fungal infection caused by Malassezia species, presenting as hypo- or hyperpigmented macules with fine scale on the trunk and proximal extremities. It is more noticeable after sun exposure due to unaffected tanning of surrounding skin. Diagnosis is clinical with KOH showing spaghetti and meatballs appearance (hyphae and spores). Treatment includes selenium sulfide shampoo, topical antifungals (ketoconazole, ciclopirox), and oral fluconazooral fluconazole (200 to 400 mg single dose or 300 mg weekly for 2 weeks) for extensive or recurrent disease, with a 50 to 80 percent response rate but high relapse rates (more than 50 percent within 2 years, requiring maintenance therapy).

\medterm{Topical Antibiotics}-- Topical antibiotics treat superficial bacterial skin infections. Mupirocin 2 percent ointment is first-line for impetigo and folliculitis (inhibits bacterial protein synthesis). Retapamulin 1 percent ointment is an alternative. Fusidic acid is used outside the US. Topical gentamicin and bacitracin are less effective and may cause contact dermatitis. Neomycin has high allergy rates and is no longer recommended. Topical antibiotics should be used for limited areas and short durations to prevent reprevent resistance. Widespread or deep infections require systemic antibiotics.

\medterm{Topical Antifungals}-- Topical antifungals treat superficial fungal infections. Allylamines (terbinafine, naftifine) inhibit squalene epoxidase, fungicidal. Azoles (clotrimazole, miconazole, ketoconazole, econazole) inhibit ergosterol synthesis, fungistatic. Ciclopirox has broad-spectrum activity. Tolnaftate is effective against dermatophytes. Application frequency: allylamines once or twice daily for 1 to 2 weeks; azoles once or twice daily for 2 to 4 weeks. Treatment should continue 1 to 2 weeks after clinical resolclinical resolution to prevent recurrence.

\medterm{Topical Calcineurin Inhibitors}-- Non-steroidal anti-inflammatory agents used for atopic dermatitis and facial eczema. They inhibit T-cell activation by blocking calcineurin, preventing cytokine production. Advantages over corticosteroids: no skin atrophy, safe for face and intertriginous areas. Side effects include transient burning and itching at application site. Black box warning for malignancy risk has not been substasubstantiated by real-world use, though FDA requires it. Intermittent therapy is safe for long-term use.

\medterm{Topical Corticosteroid Potency}-- Topical corticosteroids are classified into seven potency groups. Group 1 (ultra-high): clobetasol propionate 0.05 percent. Group 2 (high): betamethasone dipropionate 0.05 percent. Group 3 to 4 (medium): triamcinolone acetonide 0.1 percent. Group 5 to 6 (low): hydrocortisone 1 percent. Group 7 (lowest): hydrocortisone 0.5 percent. Face, groin, and intertriginous areas require low-potency steroids. Thick plaques on palms and soles may require ultra-high potency. Prolonged use causes skin atrophy,skin atrophy, telangiectasias, striae, and perioral dermatitis. Occlusion increases potency. Short-term use of appropriate potency minimises adverse effects while maximising efficacy.

\medterm{Topical Retinoids}-- Vitamin A derivatives used for acne, photoaging, and hyperpigmentation. Tretinoin (all-trans retinoic acid) is the most studied. Adapalene is more stable and less irritating. Tazarotene is the most potent. Mechanism: normalizes follicular keratinization, increases cell turnover, and stimulates collagen production. Side effects include irritation, erythema, peeling, and photosensitivity. Start with low concentration, alternate nights, and use moisturizer. Sunscreen is mandatmandatory during retinoid therapy. Results are seen over 4 to 6 months. Maintenance therapy with lower frequency is needed.

\medterm{Total Body Surface Area}-- The percentage of body surface involved by a burn or other extensive skin process. It is estimated using methods such as the rule of nines, Lund-Browder chart, or the patient's palm and is used to guide burn severity assessment and resuscitation.

\medterm{Tripe palms}-- Diffuse, velvety, thickened hyperkeratosis of the palmar skin with a cerebriform, corrugated appearance resembling the lining of a cow's stomach (tripe), which is strongly associated with underlying malignancy, most commonly gastric adenocarcinoma. It may occur alone or in conjunction with acanthosis nigricans. Recognition should prompt investigation for an internal malignancy.

\medterm{Tuberculoid leprosy}-- The paucibacillary pole of the leprosy spectrum characterised by strong cell-mediated immunity to Mycobacterium leprae, resulting in well-formed granulomas and few or no organisms in tissue. It presents with a few well-defined, hypopigmented or erythematous, anaesthetic skin macules or plaques with raised borders and associated peripheral nerve damage. Skin smear is typically negative for acid-fast bacilli.

\medterm{Ulcer}-- A primary or secondary cutaneous lesion characterized by a full-thickness excavation and discontinuity of the integument involving the complete loss of the epidermis and at least part of the papillary or reticular dermis, frequently extending into the subcutaneous adipose tissue, fascia, muscle, or bone. Because the basal lamina and follicular epithelial reservoirs are destroyed, ulcers cannot heal solely by superficial re-epithelialization; instead, they require granulation tissue formation, wound contraction, and fibrous scarring. Clinical evaluation focuses on anatomical site, size, depth, base appearance (slough, granulation, eschar), wound exudate, and margins (punched-out, sloping, rolled, or undermined). Major pathophysiological etiologies encompass: (1) Vascular insufficiency, comprising venous stasis ulcers (gaiter region, irregular, exudative, hemosiderin-stained), arterial insufficiency ulcers (distal digits, punched-out, pale, painful, worsened by elevation), and vasculitic ulcers; (2) Neuropathic ulcers (mal perforant on weight-bearing plantar surfaces in diabetic polyneuropathy); (3) Decubitus/pressure injuries; (4) Infectious ulcers (ecthyma, cutaneous leishmaniasis, atypical mycobacteria); (5) Inflammatory/neutrophilic ulcers (pyoderma gangrenosum, with undermined violaceous borders); and (6) Malignant ulcerations (ulcerated basal cell carcinoma, Marjolin ulcer arising in chronic scars). Management requires addressing the root systemic cause, wound bed preparation (debridement, moisture balance, infection control), and specialized dressings or compression therapy.

\medterm{Urticaria}-- A common vascular reaction of the skin characterized by the rapid appearance of transient, intensely pruritic, erythematous, edematous papules and plaques (wheals), frequently accompanied by deeper subcutaneous or submucosal swelling (angioedema). Pathophysiologically, urticaria is driven by the activation and degranulation of dermal mast cells and circulating basophils, which release preformed histamine, leukotrienes ($LTC_4$, $LTD_4$), prostaglandin $D_2$, and platelet-activating factor. These vasoactive mediators induce precapillary arteriolar vasodilation and postcapillary venular hyperpermeability, causing localized transudation of fluid into the superficial dermis. Chronologically classified into acute urticaria ($<6\text{ weeks}$, commonly triggered by viral infections, adverse drug reactions, IgE-mediated food allergies, or hymenoptera stings) and chronic urticaria ($\ge 6\text{ weeks}$, predominantly autoimmune, driven by IgG autoantibodies targeting the high-affinity IgE receptor $Fc\epsilon RI$ or IgE, or inducible by physical stimuli such as cold, delayed pressure, heat, vibration, or solar radiation). Individual wheals characteristically resolve without residual scarring within 24 hours. First-line medical therapy consists of second-generation non-sedating $H_1$-antihistamines (cetirizine, fexofenadine, bilastine) up to four times standard dosing; refractory chronic spontaneous urticaria responds to the monoclonal anti-IgE antibody omalizumab or cyclosporine.

\medterm{Urticaria multiforme}-- A hypersensitivity reaction in children, often triggered by infections or medications, characterised by widespread, annular or polycyclic erythematous plaques with dusky, ecchymotic, or targetoid centres and pseudovesication, differing from true urticaria by the presence of dusky discolouration and prolonged lesion duration. It is usually self-limiting and resolves within one to two weeks. Antihistamines and short courses of corticosteroids may provide symptomatic relief.

\medterm{Urticaria Pigmentosa}-- Maculopapular cutaneous mastocytosis; the most common clinical form of mastocytosis, characterized by abnormal, clonally driven accumulation and proliferation of mast cells within the skin and occasionally internal organs. In over $80\%$ of cases, it is driven by activating somatic gain-of-function mutations in the \textit{KIT} receptor tyrosine kinase proto-oncogene (most notably the D816V mutation). The condition predominantly manifests in infancy or early childhood with multiple, discrete, reddish-brown to tan, round-to-oval macules and papules symmetrically scattered across the trunk and extremities, sparing the palms, soles, and face. A pathognomonic diagnostic sign is Darier's sign: gentle rubbing, scratching, or stroking of a lesional macule triggers immediate localized mast cell degranulation, producing an intensely pruritic wheal and flare reaction strictly confined to the lesion within minutes. Histology reveals a dense, band-like or perivascular infiltration of mast cells in the papillary and reticular dermis, staining positive with Giemsa, toluidine blue, CD117 (c-kit), and tryptase. Pediatric cases typically resolve spontaneously by adolescence. Management involves avoiding mast-cell secretagogues (opioids, NSAIDs, radiocontrast media, alcohol, hot baths) and administering $H_1$- and $H_2$-antihistamines, mast cell stabilizers (oral cromolyn sodium), or leukotriene receptor antagonists.

\medterm{Uveal melanoma}-- The most common primary intraocular malignancy in adults, arising from melanocytes within the uveal tract (choroid, ciliary body, or iris), presenting with visual symptoms, a pigmented intraocular mass, or incidental finding on examination. It has a strong propensity for haematogenous metastasis, particularly to the liver, and carries a guarded prognosis. Treatment options include plaque brachytherapy, proton beam therapy, and enucleation.

\medterm{Verruca (Wart)}-- Benign epidermal hyperplasia caused by human papillomavirus (HPV). Common warts (verruca vulgaris): rough papules on hands/fingers. Plantar warts: painful lesions on soles, grow inward due to pressure. Flat warts (verruca plana): smooth, flat-topped papules on face, hands, legs. Filiform warts: finger-like projections, often on face. Genital warts (condyloma acuminata): sexually transmitted. Diagnosis: clinical. Treatment: cryotherapy, salicylic acid, cantharidin, imiquimod, podophyllin, and laser therapy. Most warts resolve spontaneously over months to years.

\medterm{Verruciform Xanthoma}-- A rare, benign mucocutaneous lesion characterized by a papillary or verrucous epithelial architecture accompanied by the selective accumulation of lipid-laden, foamy histiocytes (xanthoma cells) strictly confined to the dermal or submucosal papillae. First described by Shafer in 1971, verruciform xanthomas most commonly occur in the oral cavity (especially the masticatory mucosa of the gingiva and hard palate) and the anogenital region (vulva, scrotum, penis, perineum), with rare occurrences on extragenital skin. Clinically, it presents as a solitary, well-circumscribed, sessile or pedunculated, painless, yellow-orange, pink, or reddish-brown papillomatous plaque or nodule that is frequently misdiagnosed as a common wart (verruca vulgaris), condyloma acuminatum, or squamous cell carcinoma. Unlike conventional xanthomas, verruciform xanthoma is not associated with hyperlipidemia or systemic lipid metabolic disorders; instead, it represents a reactive phenomenon secondary to localized epithelial damage, basement membrane degeneration, and subsequent phagocytosis of lipid debris by stromal histiocytes. Histology demonstrates hyperparakeratosis with characteristic invaginations, regular acanthosis, and xanthoma cells filling the dermal papillae without extending into the reticular dermis. Conservative surgical excision is curative, with minimal recurrence risk.

\medterm{Vesicle}-- A small (<1 cm) fluid-filled blister within the epidermis, containing clear fluid (serum) or, if infected, pus. Causes: viral (herpes simplex, herpes zoster, varicella, hand-foot-mouth disease), dermatitis (contact, dyshidrotic eczema), autoimmune (pemphigus, bullous pemphigoid early), and friction burns. Grouped vesicles on an erythematous base suggest herpesvirus; dermatomal distribution suggests zoster. Diagnosis: clinical, Tzanck smear, viral PCR, biopsy (pemphigus—intraepidermal acantholysis; pemphigoid—subepidermal). Management: treat cause—antivirals for herpes, topical steroids for eczema, immunosuppression for autoimmune blisters; keep clean to prevent superinfection. Larger than 1 cm: bulla.

\medterm{Vestibular papillomatosis}-- A normal anatomical variant of the female vulva characterised by multiple, soft, finger-like or filiform papules on the labia minora and vestibule, often mistaken for genital warts (condylomata acuminata). It is a benign, asymptomatic condition requiring no treatment, and differentiation from HPV-related condylomata is important to avoid unnecessary intervention. The papules are typically symmetric and uniform in size.

\medterm{Vitiligo}-- An acquired disorder of depigmentation characterized by well-demarcated pale or depigmented macules and patches caused by loss or dysfunction of epidermal melanocytes. It is associated with autoimmune disease, may follow trauma through the Koebner phenomenon, and increases susceptibility to sunburn in depigmented areas; treatment is individualized and may include topical therapy, phototherapy, or surgery.

\medterm{Vitiligo and Pigmentation Disorders}-- Vitiligo is a chronic skin condition characterized by the loss of melanocytes, resulting in well-defined white patches on the skin and sometimes mucous membranes. The exact cause is not fully understood but involves an autoimmune process in which the body attacks its own pigment-producing cells, and it may be associated with other autoimmune diseases. Pigmentation disorders also include melasma, characterized by brown patches on the face often triggered by hormonal changes, and post-inflammatory hyperpigmentation, which follows skin injury or inflammation; treatments range from topical corticosteroids and phototherapy to depigmentation and cosmetic camouflage.

\medterm{Vitiligo Treatment}-- Vitiligo treatment options include topical corticosteroids (first-line for limited disease), topical calcineurin inhibitors (tacrolimus, pimecrolimus for face and intertriginous areas), narrowband UVB phototherapy (first-line for widespread disease), excimer laser for localized lesions, and systemic immunosuppressants for rapidly progressive disease. JAK inhibitors (ruxolitinib cream) are FDA-approved for vitiligo. Surgical options include autologous melanocyte transplantation for stable diseasestable disease (lesions stable for at least 6 to 12 months), including suction blister grafting, punch grafting, and cultured melanocyte transplantation, with repigmentation rates of 50 to 90 percent in properly selected patients.

\medterm{Warts}-- Benign epithelial proliferations of the skin and mucous membranes (verrucae) induced by infection of basal layer keratinocytes with specific genotypes of the human papillomavirus (HPV), a double-stranded non-enveloped DNA virus. HPV gains entry via microscopic disruptions in the cutaneous barrier, stimulating localized epidermal acanthosis, hyperkeratosis, and papillomatosis. Clinical variants reflect distinct viral tropisms: verruca vulgaris (common warts, HPV-1, 2, 4, 7), manifesting as hyperkeratotic, exophytic papules with punctate black dots representing thrombosed dermal capillaries, predominantly on the hands; verruca plantaris (plantar warts, HPV-1), forming painful, endophytic, thick calloused lesions on weight-bearing soles that interrupt normal skin skin-cleavage dermatoglyphic lines; verruca plana (flat warts, HPV-3, 10), presenting as multiple smooth-topped, flesh-colored papules on the face and extremities; and condyloma acuminata (anogenital warts, HPV-6, 11). Histopathology reveals marked hyperkeratosis, papillomatosis, and distinctive koilocytes (keratinocytes with pyknotic nuclei and perinuclear cytoplasmic halos). First-line therapies include chemical keratolytics (salicylic acid $17\text{--}40\%$), cryotherapy with liquid nitrogen, topical imiquimod $5\%$, intralesional antigen immunotherapy, and electrosurgical destruction.

\medterm{Warty Dyskeratoma}-- A rare, benign, solitary adnexal and keratinizing tumor originating from the hair follicle infundibulum, presenting clinically as an asymptomatic or pruritic, slow-growing, flesh-colored, yellowish-brown, or hyperkeratotic papule or nodule with a distinctive central umbilication filled with a keratotic plug. The lesion typically arises on the sun-exposed skin of the head and neck (scalp, face, neck) in middle-aged or elderly individuals, and rarely on oral or genital mucosa. Histopathologically, warty dyskeratoma displays striking similarities to Darier disease, characterized by an endophytic, cup-shaped, crater-like invagination of the epidermis opening to the surface via a follicular pore, containing extensive suprabasal acantholysis, prominent villi lined by a single layer of basal cells extending into the lacunae, and premature abnormal keratinization producing diagnostic dyskeratotic cells: \textit{corps ronds} (enlarged round cells with dark pyknotic nuclei and clear perinuclear halos in the stratum spinosum/granulosum) and \textit{grains} (flattened, grain-like cells in the stratum corneum). Despite the histopathologic resemblance, warty dyskeratoma is an isolated, sporadic lesion with no genetic linkage to the \textit{ATP2A2} mutation of Darier disease. Simple surgical excision, shave excision with curettage, or electrodessication is completely curative.

\medterm{Wen (Epidermoid Cyst)}-- A traditional and historical clinical lay term referring to an epidermoid cyst (epidermal inclusion cyst or infundibular cyst), the most common benign epithelial cyst of the skin. Epidermoid cysts arise from the traumatic implantation of epidermal elements into the dermis or the cystic dilatation and occlusion of the follicular infundibulum of the pilosebaceous unit. Clinically, they present as slow-growing, mobile, firm, dome-shaped subcutaneous nodules ranging from several millimeters to several centimeters in diameter, characteristically displaying a central comedo-like punctum that anchors the lesion to the overlying epidermis. Histologically, the cyst is encapsulated by a true stratified squamous epithelium complete with a distinct granular layer, which continuously sheds laminated sheets of keratin into the cystic lumen, creating a foul-smelling, cheesy, pasty content. Common anatomical sites include the face, neck, scalp, and upper trunk. Spontaneous or traumatic rupture of the cyst wall releases keratin into the surrounding dermis, triggering an intense, painful, sterile foreign-body granulomatous inflammation that mimics a bacterial abscess. Definitive cure requires complete surgical excision of the intact fibrous cyst wall; incomplete removal leaves residual lining that predictably results in local recurrence.

\medterm{Wood Lamp Examination}-- A non-invasive, point-of-care dermatological diagnostic technique utilizing a handheld lamp equipped with a high-pressure mercury arc and a nickel oxide / barium silicate filter ('Wood filter') that emits long-wave ultraviolet A radiation (peak wavelength $365\text{ nm}$). Performed in a completely darkened examination room, the Wood lamp excites specific fluorophores present within fungal, bacterial, pigmentary, or metabolic compounds, producing characteristic diagnostic visible fluorescence. Prominent clinical diagnostic applications include: (1) Fungal infections: hair infected by ectothrix dermatophytes (\textit{Microsporum canis}, \textit{Microsporum audouinii}) fluoresces bright blue-green due to pteridine metabolites, whereas pityriasis versicolor (\textit{Malassezia} species) fluoresces pale yellow to copper-orange; (2) Bacterial infections: erythrasma (\textit{Corynebacterium minutissimum}) displays striking coral-red fluorescence due to coproporphyrin III production, while \textit{Pseudomonas aeruginosa} wound colonization emits yellowish-green fluorescence via pyoverdine; (3) Pigmentary disorders: epidermal melanin absorbs UV light and accentuates contrast in hypopigmented lesions (revealing ash-leaf macules in tuberous sclerosis or vitiligo macules) while distinguishing epidermal melasma (accentuated) from dermal melasma (unaccentuated); and (4) Porphyria: urine of patients with porphyria cutanea tarda fluoresces pink-orange.

\medterm{Wound Healing Phases}-- Wound healing occurs in overlapping phases: hemostasis (immediate, platelet aggregation and clot formation), inflammation (days 1 to 4, neutrophils and macrophages clear debris and bacteria), proliferation (days 4 to 21, granulation tissue, angiogenesis, epithelialization, collagen deposition), and remodeling (day 21 to 2 years, collagen maturation, scar strengthening to 80 percent of original strength). Factors impairing healing include diabetes, infection, ischemia, malnutrition, and steroids.

\medterm{Xanthelasma}-- Xanthelasma palpebrarum are soft, yellow, lipid-filled papules or plaques on the medial canthus of the eyelids (upper more commonly than lower), the most common form of cutaneous xanthoma. They are composed of foam cells (lipid-laden macrophages) in the superficial dermis. Xanthelasma is associated with hyperlipidemia (elevated LDL, low HDL) in approximately 50\% of patients; the other 50\% are normolipidaemic. It may also be associated with diabetes, obesity, and cardiovascular disease (xanthelasma is an independent risk factor for atherosclerotic cardiovascular disease). Diagnosis: clinical. Treatment: options for cosmetic removal include surgical excision (elliptical excision with primary closure, risks of ectropion or scar), laser therapy (CO2, argon, or pulsed dye laser), chemical cautery (trichloroacetic acid, which should be used sparingly to avoid eyelid scarring), and radiofrequency ablation. Underlying hyperlipidemia should be screened and managed.

\medterm{Xanthoma}-- Cutaneous deposits of lipids (cholesterol, triglycerides) in the dermis or tendons, presenting as yellow-orange papules, nodules, or plaques. Xanthomas are classified by clinical morphology and location: eruptive xanthomas (sudden eruption of multiple small, red-yellow papules on buttocks, extensor surfaces — pathognomonic for hypertriglyceridaemia, chylomicronaemia syndrome, lipoprotein lipase deficiency), tuberous xanthomas (firm, painless, yellow-red nodules on elbows, knees, buttocks — seen in familial hypercholesterolaemia, Sitosterolaemia), tendinous xanthomas (subcutaneous nodules within tendons, most commonly the Achilles tendon and extensor tendons of the hands — pathognomonic for familial hypercholesterolaemia, cerebrotendinous xanthomatosis), and planar xanthomas (flat or slightly raised yellow plaques on palms, flexural creases, or intertriginous areas — associated with type III hyperlipoproteinaemia/dysbetalipoproteinaemia). Diagnosis: clinical, serum lipid panel, and biopsy (foam cells, Touton giant cells). Xanthomas are not harmful themselves but are markers of underlying lipid disorders that require treatment to reduce cardiovascular risk.

\medterm{Xeroderma}-- A clinical dermatological term describing abnormally severe dryness, roughness, and dehydration of the cutaneous surface (xerosis), which may present as an acquired condition or as the defining feature of rare congenital genodermatoses. In routine clinical practice, xeroderma denotes significant stratum corneum dehydration (water content $<10\%$) resulting from depleted intercellular lipid bilayers (ceramides, fatty acids) and deficient natural moisturizing factors (NMF), producing scaling, dullness, fine fissuring, and intractable pruritus that disrupts epidermal barrier integrity. In clinical genetics, the term forms the primary designation for \textit{Xeroderma Pigmentosum} (XP), a severe, life-threatening autosomal recessive disorder caused by defective nucleotide excision repair (NER) enzymes (mutations in \textit{XPA} through \textit{XPG} or \textit{POLH}). XP patients cannot repair ultraviolet-induced DNA photoproducts (cyclobutane pyrimidine dimers), developing extreme photosensitivity, early freckling, solar elastosis, and a $>1{,}000$-fold elevated incidence of cutaneous malignancies (basal cell carcinoma, squamous cell carcinoma, melanoma) during childhood. Management of simple xeroderma requires liberal application of barrier-repair emollients containing urea and ceramides, whereas XP demands total photoprotection and aggressive oncologic surveillance.

\medterm{Xerosis}-- Abnormally dry skin, common in elderly patients and winter months. It results from decreased sebaceous and sweat gland activity, impaired skin barrier function, and environmental factors. Treatment includes regular emollient application (cream or ointment rather than lotion), gentle cleansing, and humidification. Severe xerosis may progress to asteatotic eczema (eczema craquelé) with cracked, porcelain-like appearance.

\medterm{Zinc Ointment (Zinc Oxide)}-- Zinc oxide is a topical agent with mild antiseptic, astringent, and barrier-protective properties, used in a wide range of dermatological preparations. Mechanism: zinc is an essential trace element and cofactor for numerous metalloenzymes involved in wound healing (matrix metalloproteinases, DNA polymerases, superoxide dismutase). Zinc oxide has anti-inflammatory effects (inhibits neutrophil degranulation and mast cell degranulation), provides UV protection (physical sunscreen, blocks UVA and UVB), and forms a protective barrier over irritated or denuded skin that reduces moisture loss and protects from friction and irritants. Indications: nappy/diaper rash (the most common use — as a barrier cream that protects the perineal skin from moisture and faecal enzymes), minor wounds, abrasions, and burns, sun protection (physical sunscreen, often combined with titanium dioxide), and as a topical treatment for acne vulgaris and seborrhoeic dermatitis. Zinc sulphate is used in oral zinc supplementation for zinc deficiency and for treatment of Wilson disease (blocks intestinal copper absorption). Adverse effects: topical — skin irritation, contact dermatitis (rare). Oral zinc: nausea, gastric irritation, copper deficiency with long-term high doses (zinc inhibits copper absorption via metallothionein induction).
